%% OVERFORE-Formulario-T-12.tex — generado desde OVERFORE-Formulario-T-12.docx
%% Formato: clase overfore, modo formulario. El formato lo pone la clase; aquí sólo va el contenido.
\documentclass[formulario]{../../comun/overfore}
\numsubdoc{3}
\formulario{T-12}{Matriz de Cumplimiento Técnico y Trazabilidad}
\informe{Informe Preparatorio N° 2}
\fecha{12 de octubre de 2026}
\externos{A1,S1,S13,S3,S4,S5,S9,T10,T13,T14,T9}
\begin{document}

\apartado{Identificación del elemento declarado}
El proponente declara el cumplimiento de cada requisito transversal de las Bases técnicas transversales y de cada requerimiento propio del caso, e indica en qué sección de la propuesta se acredita y con qué evidencia se verifica. La matriz se presenta en la sección B.

\begin{tabla}{I{0.263}I{0.737}}{Ficha de registro del formulario.}{tab:ficha-registro-formulario}
  {\cab{Campo} & \cab{Contenido}}
  {\emph{Fuente:} elaboración propia a partir del formulario oficial T-12.}
Formulario & T-12. Matriz de Cumplimiento Técnico y Trazabilidad \\
Subdocumento asociado & Subdocumento 3. Esquema de solución y alcance \\
Requisitos declarados & 622: 374 requisitos transversales y 248 requerimientos propios del caso \\
Fuente de los requisitos & Bases técnicas transversales y Bases técnicas del caso Minería \\
Verificación & Cada fila indica la sección de la propuesta que acredita el requisito y la evidencia con que se verifica \\
Trazabilidad con la EDT & Cada fila cita el paquete del Formulario T-14 que ejecuta el requisito. 601 de los 622 requisitos tienen paquete \\
Corte de la matriz & Versiones de los subdocumentos de esta entrega \\
\end{tabla}
Cada requisito se declara en uno de cinco estados. La tabla siguiente indica qué significa cada uno y qué contiene la fila correspondiente.

\begin{tabla}{I{0.227}I{0.773}}{Estados de cumplimiento y su significado.}{tab:estados-cumplimiento-significado}
  {\cab{Estado} & \cab{Significado}}
  {\emph{Fuente:} elaboración propia.}
Cumple & La propuesta compromete de forma concreta todo lo que exige el requisito, y la sección indicada lo acredita. \\
Cumple parcialmente & La propuesta compromete una parte del requisito. La última columna indica lo que falta y, si corresponde, el paquete del Formulario T-14 que lo completa. \\
Por cumplir & El requisito aún no tiene tratamiento en la propuesta. La última columna indica lo que falta y, si corresponde, el entregable de la propuesta que lo cubrirá. \\
Se acredita en la instancia que fija el T-22 & El requisito corresponde al sitio web, al video o al prototipo, y se acredita en la instancia de entrega que fija el Formulario T-22. La última columna indica el entregable que lo acredita. \\
No aplica & El requisito es deseable o según caso y la propuesta no lo ofrece. La fila justifica la exclusión. \\
\end{tabla}
El resumen siguiente muestra cuántos requisitos hay en cada grupo y en qué estado se encuentran.

\begin{tabla}{I{0.300}D{0.110}D{0.140}D{0.110}D{0.220}D{0.120}}{Resumen de requisitos por grupo y estado de cumplimiento.}{tab:resumen-requisitos-grupo-estado-cumplimiento}
  {\cab{Grupo de requisitos} & \cab{Cumple} & \cab{Cumple parcialmente} & \cab{Por cumplir} & \cab{Se acredita en la instancia que fija el T-22} & \cab{No aplica}}
  {\emph{Fuente:} elaboración propia a partir de la matriz de la sección B.}
Requisitos transversales (capítulos 02 a 26) & 285 & 64 & 8 & 11 & 6 \\
Requerimientos funcionales propios del caso & 112 & 14 & 0 & 0 & 0 \\
Requerimientos no funcionales propios del caso & 82 & 14 & 0 & 0 & 0 \\
Requerimientos de transición y preparación & 14 & 1 & 0 & 0 & 0 \\
Requerimientos del proyecto & 7 & 4 & 0 & 0 & 0 \\
Total & 500 & 97 & 8 & 11 & 6 \\
\end{tabla}
\begin{horizontal}

\apartado{Listado detallado}
La matriz agrupa los requisitos por capítulo de las Bases técnicas transversales y por tipo de requerimiento propio del caso. En cada fila, la columna de sección y evidencia indica la sección de la propuesta que acredita el requisito, el paquete del Formulario T-14 que lo ejecuta y cómo se verifica.

\seccion{Requisitos transversales}
\begin{tabla}{I{0.132}I{0.266}I{0.115}I{0.178}I{0.310}}{Capítulo 02. Modelo de arquitectura de referencia.}{tab:capitulo-02-modelo-arquitectura-referencia}
  {\cab{ID requerimiento} & \cab{Descripción} & \cab{Estado} & \cab{Componente que lo satisface} & \cab{Sección de la propuesta y evidencia}}
  {\emph{Fuente:} elaboración propia a partir de las Bases técnicas transversales.}
RT-02.01 Obligatorio & La solución se organizará en las ocho capas del numeral 2.1. El Proponente presentará el diagrama de la arquitectura lógica identificando cada capa, sus componentes y las interfaces entre ellas. & Cumple & Arquitectura lógica de ocho capas de SITCO Aranda & Subdoc. 4, 4.1.3.1 (Figura \refext{S4}{fig:vista-general-sitco-aranda-ocho-capas}); Subdoc. 4, 4.1.3.2 (Tabla \refext{S4}{tab:requerimientos-funcionales-modulo-plataforma-etapa}); Subdoc. 4, 4.1.7.2 (Tabla \refext{S4}{tab:llamadas-sincronas-entre-componentes-umbral}); Formulario T-14, paquete 1.3.1.1. Evidencia: Diagrama de arquitectura lógica de ocho capas (Subdoc. 4, 4.1.3.1) \\
RT-02.02 Obligatorio & La arquitectura será modular, con límites de contexto explícitos y acoplamiento débil. Se rechazará toda arquitectura monolítica que no permita desplegar de forma independiente sus componentes críticos. & Cumple & Descomposición modular en 9 módulos y 2 plataformas con límites de contexto y acoplamiento solo por eventos o interfaces versionadas & Subdoc. 3, 3.2.1 y 3.3.2; Subdoc. 4, 4.1.6 (Tabla \refext{S4}{tab:interfaces-sincronas-eventos-publicados-datos-consumidos}); Subdoc. 4, 4.1.2.2 (Tabla \refext{S4}{tab:registro-decision-adr-01-servicios-modulares}); Subdoc. 4, 4.1.2.3 (Tabla \refext{S4}{tab:registro-decision-adr-01-servicios-modulares}, ADR-01); Formulario T-14, paquete 1.3.1.1. Evidencia: Documento de arquitectura ISO/IEC/IEEE 42010 (Subdoc. 4), con los límites de contexto de cada módulo y la decisión ADR-01 de despliegue separado \\
RT-02.03 Obligatorio & La descripción de la arquitectura se ajustará a ISO/IEC/IEEE 42010, con vistas lógica, de procesos, de despliegue, de datos y de seguridad. & Cumple & Vistas ISO/IEC/IEEE 42010: lógica, procesos, despliegue, datos y seguridad & Subdoc. 4, 4.1.2.4 (vistas ISO/IEC/IEEE 42010); Subdoc. 4, 4.1.3 a 4.1.6 (vista lógica); Subdoc. 4, 4.1.7 (procesos); Subdoc. 4, 4.1.8 (seguridad); Subdoc. 4, 4.2 (despliegue); Subdoc. 4, 4.1.4.6 (vista de datos, desarrollada en Subdoc. 5); Formulario T-14, paquetes 1.3.1.1, 1.3.1.2, 1.3.2.1 y 1.3.3.1. Evidencia: Documento de arquitectura con las cinco vistas (lógica, procesos, despliegue, datos, seguridad) \\
RT-02.04 Obligatorio & El Proponente mantendrá un registro de decisiones de arquitectura (ADR) fechado, con la alternativa escogida, las alternativas descartadas y el criterio de decisión. El registro es entregable contractual y se actualizará durante toda la ejecución. & Cumple & Registro ADR-01 a ADR-08 con fecha, alternativas descartadas y criterio; cambios como registro nuevo & Subdoc. 4, 4.1.2.3 (Tablas \refext{S4}{tab:registro-decision-adr-01-servicios-modulares} a \refext{S4}{tab:correspondencia-canales-modulos-plataformas-esquema-solucion}, ADR-01 a ADR-08); Formulario T-14, paquete 1.3.1.1. Evidencia: Registro de decisiones de arquitectura (ADR) fechado, con alternativas descartadas y criterio, entregable contractual (OVERFORE-Subdocumento4-Anexos, Anexo E, control 5.8) \\
RT-02.05 Obligatorio & La capa de servicios de negocio será sin estado. El estado de sesión y el estado de proceso residirán en almacenes externos con alta disponibilidad. & Cumple & Servicios de negocio sin estado con caché externa de sesión y de proceso de alta disponibilidad (Azure Managed Redis) & Subdoc. 4, 4.1.4.4; Subdoc. 4, 4.1.1.2 (Tabla \refext{S4}{tab:tecnologias-capas-4-5-alternativa-evaluada}); Subdoc. 4, Anexos (Formulario T-11, caché de sesión y de proceso); Formulario T-14, paquetes 1.5.3.2 y 1.6.1.1. Evidencia: Revisión de arquitectura y prueba de reinicio de instancias sin pérdida de sesión ni de proceso \\
RT-02.06 Obligatorio & Toda operación de escritura expuesta a reintentos será idempotente, con clave de idempotencia declarada por el cliente y ventana de deduplicación documentada. & Cumple parcialmente & Conectores del Integrador de Sistemas con tiempo de espera, reintento, cortacircuito, mamparo y límite de tasa; puertas de enlace con límites de tasa & Subdoc. 4, 4.1.7.2; Subdoc. 4, 4.1.2.4; Subdoc. 2, Anexos (RNF-TRV-02.06, ventana de deduplicación de 72 horas); Formulario T-14, paquetes 1.4.2.1 y 1.6.1.1. Evidencia: Prueba de reenvío de solicitudes con 0 efectos duplicados (RNF-TRV-02.06) \\
RT-02.07 Obligatorio & Los flujos de eventos garantizarán entrega al menos una vez, con deduplicación en el consumidor y orden garantizado dentro de la partición o del agregado cuando el proceso lo exija. & Cumple & Bus de eventos (Azure Service Bus Premium) con orden por agregado, deduplicación y cola de fallidos & Subdoc. 4, 4.1.7.3; Subdoc. 4, ADR-03; Subdoc. 4, 4.1.1.2 (Tabla \refext{S4}{tab:tecnologias-capas-4-5-alternativa-evaluada}); Subdoc. 2, Anexos (RNF-TRV-02.07); Formulario T-14, paquete 1.4.2.1. Evidencia: Prueba de inyección de fallas y reenvíos en el bus de eventos (RNF-TRV-02.07) \\
RT-02.08 Obligatorio & La solución implementará patrones de resiliencia demostrables: reintento con retroceso exponencial y variación aleatoria, cortacircuitos, mamparos de aislamiento, límites de tasa y tiempo de espera explícito en toda llamada remota. No se admiten llamadas remotas sin tiempo de espera. & Cumple & Conectores del Integrador de Sistemas con tiempo de espera, reintento, cortacircuito, mamparo y límite de tasa; puertas de enlace con límites de tasa & Subdoc. 4, 4.1.7.1; Subdoc. 4, 4.1.2.4; Subdoc. 4, 4.1.4.3; Subdoc. 4, 4.1.7.2 (Tabla \refext{S4}{tab:llamadas-sincronas-entre-componentes-umbral}); Subdoc. 2, Anexos (RNF-TRV-02.08); Formulario T-14, paquete 1.6.1.1. Evidencia: Revisión de configuración y pruebas de resiliencia (RNF-TRV-02.08). \\
RT-02.09 Obligatorio & La solución degradará de forma elegante: ante la indisponibilidad de un componente no crítico deberá continuar operando en modo reducido, informando la degradación a la persona usuaria, y nunca fallar de forma total. & Cumple & Nodo de Faena y dispositivos de terreno (modo sin enlace); conectores con cola y dato marcado con su antigüedad; vista única que indica la fuente que falta & Subdoc. 3, 3.3.6; Subdoc. 4, 4.1.7.1 (Tabla \refext{S4}{tab:conectores-integrador-sistemas-etapa-lugar-ejecucion}); Subdoc. 4, 4.1.9 (Tabla \refext{S4}{tab:modulo-plataforma-sostiene-cada-funcion-durante}); Subdoc. 2, Anexos (RNF-DIS-007); Formulario T-14, paquetes 1.5.3.2 y 1.6.1.1. Evidencia: Prueba de inyección de fallas y de corte controlado de 48 h (RNF-DIS-007, RNF-DIS-001). \\
RT-02.10 Obligatorio & Las capas de aplicación e integración escalarán horizontalmente de forma automática, con umbrales, límites superiores y costo asociado declarados en la oferta. & Cumple parcialmente & Grupo de nodos de servicios en AKS con réplicas por zona y máximo de nueve nodos & Subdoc. 4, 4.2.2.2 (escalado de servicios en AKS); Subdoc. 4, 4.2.4.3 (capacidad de cómputo); Subdoc. 4, 4.2.4.6 (umbrales de ampliación de disco y Nodo); OVERFORE-Subdocumento4-Anexos, Anexo A (T11-10); Formulario T-14, paquete 1.6.1.1. Falta: métrica y valor del umbral (CPU, memoria o cola) que gatilla el escalamiento de servicios de negocio e integración, y su costo asociado. Evidencia: Prueba de carga creciente en Preproducción hasta gatillar el escalamiento (RNF-CAP-008). \\
RT-02.11 Obligatorio & El Proponente declarará explícitamente los puntos únicos de falla que subsistan en su arquitectura y justificará por qué son aceptables. Omitir esta declaración cuando existan puntos únicos de falla se evaluará como observación grave. & Cumple & Puntos únicos de falla declarados y justificados: sala técnica, tramo común de enlaces, energía, red LTE o diodo de datos, y región primaria o identidad & Subdoc. 4, 4.2.3.4 (puntos que afectan a más de un componente y su acotamiento); Formulario T-14, paquete 1.3.1.2. Evidencia: Declaración de puntos únicos de falla con su justificación (Subdoc. 4, 4.2.3.4) \\
RT-02.12 Según caso & La solución admitirá su replicación a nuevas unidades, sitios, sucursales o filiales del Cliente sin rediseño arquitectónico, mediante parametrización o multi-tenencia. & Cumple & Módulo Administración y Servicios Comunes (parametrización) & Subdoc. 3, 3.4.3; Subdoc. 2, 2.5.1; Formulario T-14, paquete 1.4.3.1. Evidencia: Ejercicio cronometrado de configuración de una unidad operacional adicional, sin modificar código (RNF-POR-001) \\
RT-02.13 Obligatorio & El Proponente presentará un modelo de dominio del negocio del caso, con las entidades principales, sus relaciones y los eventos de negocio que las modifican. & Cumple & Modelo de datos de la solución, con nueve dominios de información & Subdoc. 5, 5.1.1 (Figura \refext{S5}{fig:vista-general-nueve-dominios-informacion-eventos}, Tabla \refext{S5}{tab:eventos-entre-dominios-vista-general}; Tablas A.1 a A.3 del Anexo A de OVERFORE-Subdocumento5-Anexos); Subdoc. 5, 5.1.6; Subdoc. 5, Anexos (A.3); Formulario T-14, paquete 1.3.2.1. Evidencia: Modelo de dominio del negocio: entidades, relaciones y eventos \\
RT-02.14 Deseable & Se valorará la aplicación documentada de patrones de arquitectura evolutiva que permitan sustituir un componente sin reescribir la solución: capa anticorrupción frente a sistemas heredados, estrangulamiento progresivo y abstracción de proveedores. & Cumple parcialmente & Integrador de Sistemas (un conector por sistema) como capa anticorrupción & Subdoc. 4, 4.1.7.1 (capa anticorrupción por conector); Subdoc. 4, 4.1.2.4 (puertos abstractos); Subdoc. 3, 3.3.2.2 y 3.4.3; Formulario T-14, paquetes 1.3.1.1 y 1.4.2.1. Falta: estrangulamiento progresivo de sistemas heredados; la migración de planillas es una carga única (Subdoc. 4, Tabla \refext{S4}{tab:conectores-integrador-sistemas-etapa-lugar-ejecucion}). Evidencia: Registro de decisiones de arquitectura (Subdoc. 4, Tabla \refext{S4}{tab:registro-decision-adr-05-lectura-historiador}, ADR-04). \\
\end{tabla}
\begin{tabla}{I{0.132}I{0.267}I{0.115}I{0.178}I{0.309}}{Capítulo 03. Modelo híbrido: nube y on-premise.}{tab:capitulo-03-modelo-hibrido-nube-on}
  {\cab{ID requerimiento} & \cab{Descripción} & \cab{Estado} & \cab{Componente que lo satisface} & \cab{Sección de la propuesta y evidencia}}
  {\emph{Fuente:} elaboración propia a partir de las Bases técnicas transversales.}
RT-03.01 Obligatorio & El Proponente declarará el proveedor de nube pública, la región primaria y la región secundaria utilizadas. El proveedor deberá contar con presencia de región o zona en Chile o en Sudamérica. & Cumple & Microsoft Azure; región primaria Azure Chile Central y secundaria Azure Brazil South & Subdoc. 3, 3.3.5; Subdoc. 1, sección 1,6; Formulario T-14, paquetes 1.3.1.4 y 1.6.1.1. Evidencia: Declaración de proveedor y regiones (Subdoc. 4, 4.3.1 y 4.3.2) y carta de alianza (Subdoc. 1) \\
RT-03.02 Obligatorio & Todos los componentes con requisito de alta disponibilidad se desplegarán en al menos dos zonas de disponibilidad. No se aceptará un diseño en una sola zona. & Cumple & Servicios centrales en Azure Chile Central, al menos dos zonas de disponibilidad & Subdoc. 3, 3.3.5; Formulario T-14, paquetes 1.5.3.2 y 1.6.1.1. Evidencia: Diagrama de despliegue (Subdoc. 4, 4.2) y prueba de caída de zona \\
RT-03.03 Obligatorio & La totalidad de la infraestructura se definirá como código, versionada en el repositorio del Cliente, revisable y reproducible. No se admite infraestructura creada manualmente por consola, salvo la cuenta raíz inicial, cuya creación deberá documentarse. & Cumple & Infraestructura como código con OpenTofu, versionada en el repositorio del Cliente y sin creación manual por consola & Subdoc. 4, 4.2.1.1; Subdoc. 4, Anexos (Formulario T-11, automatización de infraestructura); Subdoc. 6, 6.2.2.3; Subdoc. 3, 3.2.2.4; Formulario T-14, paquete 1.6.1.1. Evidencia: Repositorio de infraestructura como código en el repositorio del Cliente. \\
RT-03.04 Obligatorio & La red se segmentará por capas, con subredes privadas para aplicación y datos, y exposición pública restringida a la capa de borde. Ningún componente de datos será alcanzable desde Internet. & Cumple & Red de producción en subredes privadas por capa, con exposición pública limitada a Front Door y almacenes solo con dirección privada & Subdoc. 4, 4.2.2.2 (Figura \refext{S4}{fig:actores-caminos-acceso-sobre-topologia-fisica}); Subdoc. 4, 4.2.3.3; Subdoc. 4, Anexos (Formulario T-11, puntos de conexión privados); Formulario T-14, paquetes 1.3.1.3 y 1.6.1.1. Evidencia: Diagrama de red por capas y revisión de reglas de acceso \\
RT-03.05 Obligatorio & El Proponente privilegiará servicios administrados por sobre servicios autoadministrados cuando ello reduzca el riesgo operacional, y justificará cada excepción. & Cumple & Servicios administrados de nube; tres cargas autoadministradas justificadas; excepciones de faena justificadas por operación sin enlace & Subdoc. 4, 4.2.1.1 (preferencia por servicios administrados); Subdoc. 4, 4.2.2.1 (excepciones justificadas); Subdoc. 4, 4.1.1.1 a 4.1.1.4 (Tablas 4.1 a 4.06, alternativa descartada por tecnología); Formulario T-14, paquete 1.3.1.2. Evidencia: Justificación de cada excepción a servicios administrados (Subdoc. 4, 4.1.1 y 4.2.2.1) \\
RT-03.06 Obligatorio & Se aplicarán prácticas FinOps: etiquetado obligatorio de todos los recursos por ambiente, módulo y centro de costo; presupuestos con alertas de desviación; y reporte mensual de consumo desglosado entregado al Cliente. & Cumple & Prácticas FinOps de la nube: etiquetado por ambiente, módulo y centro de costo, presupuestos con alertas e informe mensual de consumo & Subdoc. 4, 4.2.1.1; Subdoc. 4, Anexos (Formulario T-11, nube principal alcance FinOps mensual); Subdoc. 6, 6.2.2.2; Subdoc. 3, 3.2.2.6; Formulario T-14, paquete 1.6.1.1. Evidencia: Reporte mensual de consumo desglosado entregado al Cliente \\
RT-03.07 Obligatorio & El Proponente declarará su estrategia de reversibilidad y de mitigación del bloqueo por proveedor, identificando qué componentes son portables, cuáles no lo son y cuál sería el esfuerzo estimado de una migración. & Cumple & Portabilidad declarada por componente: contenedores, motores de datos e IaC portables; entrada pública, puerta de exposición, bus, identidad, bóveda y observabilidad no portables sin trabajo de migración & Subdoc. 4, 4.2.2.1 (portabilidad por componente); Subdoc. 4, 4.1.1 (criterio de portabilidad); Formulario T-14, paquete 1.1.9.1. Evidencia: Plan de reversibilidad inicial con inventario, portabilidad por componente y esfuerzo estimado de migración (1.1.9.1). \\
RT-03.08 Deseable & La solución empleará instancias reservadas, planes de ahorro o capacidad comprometida cuando el perfil de carga lo justifique, y lo reflejará en la estructura de costos. & Cumple parcialmente & Capacidad base de producción reservada por tres años; crecimiento por uso en cambio de turno y cierres & Subdoc. 4, 4.2.1.1 (capacidad base reservada 3 años); OVERFORE-Subdocumento4-Anexos, Anexo A (T11-01); Formulario T-14, paquete 1.3.1.2. Falta: reflejar la capacidad comprometida en la estructura de costos (Oferta Económica). Evidencia: Estructura de costos con capacidad comprometida (Oferta Económica). \\
RT-03.09 Deseable & Se valorará el uso de cómputo sin servidor o de contenedores administrados para las cargas de perfil variable, con el análisis comparativo de costo frente a instancias permanentes. & Cumple parcialmente & Contenedores administrados en AKS con escalado para carga variable; base reservada y crecimiento por uso; sin comparación de costo & Subdoc. 4, 4.2.2.2 (escalado de servicios en AKS); Subdoc. 4, 4.2.1.1 (base reservada y crecimiento por uso); Formulario T-14, paquete 1.3.1.2. Falta: análisis comparativo de costo entre contenedores administrados o cómputo sin servidor e instancias permanentes. Evidencia: Análisis comparativo de costo entre cómputo sin servidor e instancias permanentes (Subdoc. 4, 4.2.1.1). \\
RT-03.10 Obligatorio & El componente on-premise operará de forma autónoma y degradada ante la pérdida total del enlace con la nube, durante un período mínimo de 24 horas continuas o el mayor que fije el caso. & Cumple & Nodo de Faena (sala técnica, Etapa 1) & Subdoc. 3, 3.3.6, Tabla \refext{S3}{tab:funciones-sitco-aranda-durante-corte-enlace} y \refext{S3}{tab:requerimientos-funcionales-modulo-etapa-entrega}.5; Formulario T-14, paquetes 1.4.2.2 y 1.5.5.2. Evidencia: Corte controlado de 48 horas ejecutado en faena (RNF-DIS-001) \\
RT-03.11 Obligatorio & Durante la operación desconectada, la solución continuará registrando las transacciones operacionales críticas de forma local, con integridad garantizada y sin pérdida de datos. & Cumple & Dispositivos de terreno y Nodo de Faena & Subdoc. 3, 3.3.6; Formulario T-14, paquete 1.4.2.2. Evidencia: Conciliación de registros generados y sincronizados en tres cortes controlados (RNF-DIS-003) \\
RT-03.12 Obligatorio & Restablecido el enlace, la sincronización será automática, con reconciliación determinista de conflictos, regla de resolución documentada y bitácora auditable de las decisiones aplicadas. & Cumple & Nodo de Faena (sala técnica, Etapa 1) y bitácora de conciliación en la nube & Subdoc. 3, 3.3.6; Formulario T-14, paquetes 1.4.2.2 y 1.4.5.1. Evidencia: Corte controlado de 48 h con cambio de turno, midiendo el tiempo hasta la consistencia completa (RNF-DES-012) \\
RT-03.13 Obligatorio & El Proponente declarará qué funciones NO estarán disponibles en modo desconectado y qué procedimiento manual las suple. La ausencia de esta declaración se evaluará como observación grave. & Cumple & Tabla de funciones durante un corte del enlace & Subdoc. 3, Tabla \refext{S3}{tab:funciones-sitco-aranda-durante-corte-enlace}; Formulario T-14, paquete 1.4.2.2. Evidencia: Declaración de funciones no disponibles y procedimiento manual (Tabla \refext{S3}{tab:funciones-sitco-aranda-durante-corte-enlace}); ratificación en el corte controlado de 48 h \\
RT-03.14 Obligatorio & Los equipos on-premise críticos serán redundantes. El almacenamiento local tolerará la falla de al menos un disco; el Proponente declarará el nivel RAID escogido y lo justificará frente a las alternativas. & Cumple & Nodo de Faena: tres servidores en clúster, dos discos NVMe en RAID 1 por servidor y PostgreSQL con réplica síncrona & Subdoc. 4, 4.2.3.2; Subdoc. 4, Anexos (Formulario T-11, servidores del Nodo servidores de gabinete); Subdoc. 3, 3.3.5; Formulario T-14, paquete 1.6.1.2. Evidencia: Especificación técnica de almacenamiento (Formulario T-11) \\
RT-03.15 Obligatorio & Los sistemas on-premise se endurecerán conforme a los CIS Benchmarks aplicables, con gestión centralizada de parches y ventana de aplicación acordada con el Cliente. & Cumple & Servidores de faena endurecidos con CIS (Ubuntu Pro) con parches centralizados vía Azure Arc & OVERFORE-Subdocumento4-Anexos, Anexo E (controles 8.7 y 8.8, CIS y Azure Arc); Subdoc. 4, 4.2.3.1 (entrega por GitOps con Azure Arc); Formulario T-11 (T11-26 y T11-59); Formulario T-14, paquete 1.6.1.2. Evidencia: Informe de línea base CIS y calendario de parches. \\
RT-03.16 Obligatorio & El monitoreo del componente on-premise se integrará a la misma plataforma de observabilidad que la nube, con alertamiento unificado. & Cumple & Capa de observabilidad con colector local en Nodo y gabinetes, Azure Monitor y consola operativa única con alertas por síntoma & Subdoc. 4, 4.1.4.8; Subdoc. 4, 4.2.2.2; Subdoc. 4, Anexos (Formulario T-11, plataforma de observabilidad); Subdoc. 3, 3.2.2.6; Formulario T-14, paquetes 1.4.2.2 y 1.6.1.1. Evidencia: Consola única de estado de componentes (RNF-OPE-001) \\
RT-03.17 Obligatorio & El enlace entre el sitio on-premise y la nube será redundante, con caminos físicos y proveedores distintos, y conmutación automática con tiempo de conmutación declarado. & Cumple & Fibras A y B de proveedores y trazados distintos más enlace satelital, con SD-WAN y conmutación automática en 5 s o menos & Subdoc. 4, 4.2.3.2; Subdoc. 4, Anexos (Formulario T-11, router SD-WAN fibra B); Subdoc. 3, 3.3.4 y 3.3.5; Formulario T-14, paquetes 1.3.1.3 y 1.6.1.3. Evidencia: Prueba de conmutación entre enlaces (Subdoc. 4) \\
RT-03.18 Obligatorio & Los dispositivos de borde y de terreno se administrarán de forma remota y centralizada: inventario, configuración, actualización de firmware y de aplicación, bloqueo y borrado remoto. & Cumple & Microsoft Intune para tabletas y consolas, Azure Arc para servidores de faena y gestión de firmware de las botoneras desde el Nodo de Faena & Subdoc. 4, 4.1.1.3 (Tabla \refext{S4}{tab:sistema-operativo-cada-dominio-ejecucion-fin}); Subdoc. 4, 4.1.3.3; Subdoc. 4, 4.1.5.1; Subdoc. 4, 4.2.1.1; Subdoc. 4, Anexos (Formulario T-11, Azure Arc); Subdoc. 3, 3.3.3; Formulario T-14, paquetes 1.3.1.3, 1.4.5.1 y 1.6.2.1. Evidencia: Prueba de bloqueo remoto en 15 min o al recuperar el enlace (RF-CNX-004) \\
RT-03.19 Deseable & Se valorará el procesamiento en el borde de las cargas que lo admitan —filtrado, agregación previa, inferencia local— reduciendo el volumen transferido y la dependencia del enlace. & Cumple & Procesamiento en faena: validación de acceso, balance del turno en el Nodo, alertas locales y series comprimidas 4:1 antes de transferir & Subdoc. 4, 4.2.2.1 (funciones que se ejecutan en la faena); Subdoc. 4, 4.1.6.1.3 (balance del turno en el Nodo); Subdoc. 4, 4.2.4.2 (Tabla \refext{S4}{tab:volumetria-referencia-capacidad-diseno}, \ref{tab:capitulo-10-disponibilidad-continuidad-resiliencia} GB a 3 GB en un corte de 48 h); Formulario T-14, paquetes 1.3.1.3, 1.4.2.2 y 1.4.5.1. Evidencia: Declaración de funciones de borde y de volumen transferido (Subdoc. 4, 4.2.2.1 y Tabla \refext{S4}{tab:volumetria-referencia-capacidad-diseno}) \\
RT-03.20 Obligatorio & El Proponente dimensionará el ancho de banda requerido por sitio, en régimen normal y en peak, y lo justificará con el cálculo de volumen de transacciones y de datos. & Cumple & Ancho de banda por sitio en régimen y en punta, con su cálculo de volumen; en faena, 8 Mbps en régimen y 18 Mbps en punta (54 con 3x) & Subdoc. 4, 4.2.4.5 (faena: régimen y punta); OVERFORE-Subdocumento4-Anexos, Anexo D (D.3.5, Tabla D.8); Formulario T-11 (T11-63 y T11-64); Formulario T-14, paquete 1.3.1.3. Evidencia: Cálculo de ancho de banda por sitio, en régimen normal y en punta (Anexo D, D.3.5). \\
RT-03.21 Obligatorio & La conexión entre la red del Cliente y la nube se establecerá mediante enlace privado dedicado o red privada virtual con cifrado, según lo que el volumen y la criticidad justifiquen. & Cumple & Túneles IPsec activo-activo (AES-256-GCM) por fibra hacia la puerta VPN de Azure; personal del CIO por acceso de confianza cero & Subdoc. 4, 4.2.3.2; Subdoc. 4, 4.2.2.2; Subdoc. 4, Anexos (Formulario T-11, puertas de red virtual); Formulario T-14, paquete 1.3.1.3. Evidencia: Diseño de conexión privada cifrada entre faena, CIO y nube \\
RT-03.22 Obligatorio & El acceso remoto de las personas trabajadoras del Cliente, incluido el trabajo desde el hogar, se resolverá con acceso a la red de confianza cero, con verificación de postura del dispositivo. No se admite exponer servicios internos directamente a Internet. & Cumple & Acceso de confianza cero (Entra Private Access) con doble factor y postura del dispositivo, sin VPN ni exposición de servicios internos & Subdoc. 4, 4.1.4.2; Subdoc. 4, 4.1.8.2 (Tabla \refext{S4}{tab:servicios-alcanzables-desde-fuera-red-compania}); Subdoc. 4, 4.2.2.2; Subdoc. 4, Anexos (Formulario T-11, acceso de confianza cero trabajo remoto); Subdoc. 3, 3.3.4; Formulario T-14, paquetes 1.3.1.3, 1.3.3.1 y 1.3.3.3. Evidencia: Revisión de la arquitectura de seguridad (Subdoc. 4) \\
RT-03.23 Según caso & La red inalámbrica de los sitios operacionales, cuando el caso la requiera, contará con segmentación por tipo de dispositivo, autenticación por certificado o credencial de empresa y cobertura verificada mediante estudio de sitio. & Cumple & Segmento dedicado de SITCO en la red LTE privada existente, con acceso a la puerta operacional por certificado del dispositivo y estudio de cobertura del rajo & Subdoc. 4, 4.2.3.2; Subdoc. 4, 4.2.3.3; Subdoc. 4, 4.2.3.4; Subdoc. 4, Anexos (Formulario T-11, estudio de cobertura LTE); Subdoc. 3, 3.3.4; Formulario T-14, paquetes 1.3.1.3 y 1.3.1.5. Evidencia: Estudio de cobertura del rajo con proyección a tres años (RP-004). \\
RT-03.24 Deseable & El Proponente declarará la calidad de servicio y la priorización de tráfico aplicada a las transacciones operacionales críticas frente al tráfico administrativo. & Cumple parcialmente & Calidad de servicio en el SD-WAN que prioriza la operación en vivo sobre la sincronización diferida & Subdoc. 4, 4.2.3.2; Subdoc. 4, 4.2.4.5; Subdoc. 4, 4.1.9.2; Subdoc. 4, Anexos (Formulario T-11, router SD-WAN, QoS por clase); Formulario T-14, paquete 1.3.1.3. Falta: clases de tráfico. Evidencia: Declaración de calidad de servicio y priorización de tráfico (Subdoc. 4, 4.2.3.2 y 4.2.4.5). \\
\end{tabla}
\begin{tabla}{I{0.139}I{0.263}I{0.112}I{0.201}I{0.286}}{Capítulo 04. Ambientes, entrega continua y configuración.}{tab:capitulo-04-ambientes-entrega-continua-configuracion}
  {\cab{ID requerimiento} & \cab{Descripción} & \cab{Estado} & \cab{Componente que lo satisface} & \cab{Sección de la propuesta y evidencia}}
  {\emph{Fuente:} elaboración propia a partir de las Bases técnicas transversales.}
RT-04.01 Obligatorio & Los cinco ambientes del numeral 4.1 estarán habilitados y operativos como condición del hito H3 del Formulario E-25. & Cumple & Ambientes de Desarrollo, QA, Preproducción y Producción habilitados con el hito H3 (mes 6);\newline Recuperación ante Desastres habilitado con el hito H3 (mes 6) & Subdoc. 4, 4.1.3.4 y 4.2.3.1 (Figura \refext{S4}{fig:nube-principal-chile-central-acceso-produccion}); Subdoc. 9, 9.3.1; Formulario T-17, hito H3; Formulario T-13, Tabla \refext{T13}{tab:ambientes-prueba-uso-datos-ubicacion} del Formulario T-13; Formulario T-14, paquete 1.6.1.1. Evidencia: Acta de habilitación de los cinco ambientes. \\
RT-04.02 Obligatorio & Preproducción será equivalente a producción en topología, versiones de componentes y configuración. Las diferencias que subsistan por costo se declararán expresamente y se justificarán. & Cumple & Clúster de Preproducción en tres zonas con las mismas versiones y configuración que Producción; banco de pruebas que reproduce el Nodo de Faena & Subdoc. 4, 4.1.3.4; Subdoc. 4, 4.2.3.1; Formulario T-13, Tabla \refext{T13}{tab:ambientes-prueba-uso-datos-ubicacion} del Formulario T-13; Formulario T-14, paquete 1.6.1.1. Evidencia: Comparación de topología y versiones entre Preproducción y Producción, con diferencias justificadas. \\
RT-04.03 Obligatorio & El código residirá en un sistema de control de versiones con ramas protegidas, revisión obligatoria por pares y prohibición de escritura directa sobre la rama principal. & Cumple & GitHub Enterprise Cloud en organización del Cliente, con rama principal protegida, solicitudes de integración y revisión obligatoria por pares & Subdoc. 3, 3.2.2.4; Subdoc. 9, 9.2.1 y 9.2.2; Subdoc. 6, 6.2.2.4 (Tabla \refext{T10}{tab:herramientas-evidencia-integracion} del Formulario T-10); Subdoc. 9, Anexos; Formulario T-14, paquete 1.4.1.1. Evidencia: Configuración del repositorio del Cliente \\
RT-04.04 Obligatorio & Existirá trazabilidad completa entre requerimiento, incidencia, cambio de código, prueba ejecutada y despliegue realizado. & Cumple & Cadena de trazabilidad requerimiento, incidencia, cambio de código, prueba y despliegue, con registro de despliegues y Formulario T-12 & Subdoc. 9, 9.2.5; Subdoc. 9, 9.2.1; Subdoc. 6, 6.2.1.1y 6.2.2.1; Formulario T-14, paquetes 1.2.2.2, 1.2.2.4, 1.2.5.2 y 1.2.5.4. Evidencia: Matriz de trazabilidad requerimiento a despliegue \\
RT-04.05 Obligatorio & El flujo de integración continua ejecutará, como mínimo: compilación, pruebas unitarias, análisis estático de código, análisis de composición de software, escaneo de secretos y escaneo de imágenes de contenedor, con criterios de bloqueo automático del despliegue. & Cumple & Flujo de integración continua con compilación, pruebas unitarias, análisis estático y de composición, escaneo de secretos y de imágenes de contenedor, y criterios de bloqueo & Subdoc. 6, 6.2.2.2 (integración continua); Formulario T-14, paquetes 1.1.4.2, 1.4.1.1, 1.5.1.1 y 1.5.1.3. Evidencia: Configuración del flujo de integración continua con criterios de bloqueo. \\
RT-04.06 Obligatorio & Los despliegues serán automatizados y reproducibles, con reversión automatizada y sin intervención manual en el paso a producción. & Cumple & Flujo de entrega continua con infraestructura reproducible, despliegue sin intervención manual y reversión automática & Subdoc. 6, 6.2.2.3; Subdoc. 3, 3.2.2.4; Subdoc. 4, 4.2.3.1; Formulario T-14, paquete 1.4.1.2. Evidencia: Registro de despliegues automatizados y reversión \\
RT-04.07 Obligatorio & La estrategia de despliegue permitirá liberar sin interrupción del servicio: azul-verde, canario o despliegue progresivo. Se declarará cuál se emplea y se demostrará en Preproducción antes de cada paso a producción. & Cumple & Despliegue azul-verde en la nube, por extracción en el Nodo de Faena y por anillos en terreno, ensayado en Preproducción & Subdoc. 6, 6.2.2.3; Subdoc. 4, 4.2.3.1; Subdoc. 2, Anexos (RNF-DIS-008); Formulario T-14, paquete 1.4.1.2. Evidencia: Despliegue demostrado en Preproducción antes de cada paso a producción (RNF-DIS-008) \\
RT-04.08 Obligatorio & Toda configuración estará externalizada del artefacto y gestionada por ambiente. Un mismo artefacto deberá poder promoverse de QA a Preproducción y a Producción sin recompilación. & Cumple & Flujo de entrega continua con imagen firmada, configuración externalizada por ambiente y promoción sin recompilación entre QA, Preproducción y Producción & Subdoc. 6, 6.2.2.3 (configuración y secretos fuera del artefacto); Formulario T-14, paquete 1.4.1.2. Evidencia: Promoción de un mismo artefacto entre ambientes sin recompilación. \\
RT-04.09 Obligatorio & Los secretos residirán en un gestor de secretos con rotación automática y auditoría de acceso. Queda prohibida toda credencial embebida en código, imágenes o archivos de configuración. & Cumple & Almacén de secretos con rotación automática y auditoría; Azure Key Vault por región; sin credenciales en código ni imágenes & Subdoc. 6, 6.2.2.3; Subdoc. 4, 4.2.2.2; Formulario T-14, paquete 1.4.1.3. Evidencia: Configuración del gestor de secretos y registro de rotación \\
RT-04.10 Obligatorio & Las migraciones de esquema de base de datos serán versionadas, reversibles y ejecutadas de forma automatizada, con estrategia de compatibilidad que permita convivir dos versiones de la aplicación durante el despliegue. & Cumple & Migraciones de esquema versionadas y reversibles, con dos versiones de la aplicación conviviendo durante el despliegue & Subdoc. 6, 6.2.2.3 (migraciones de esquema seguras); Subdoc. 5, 5.3; Formulario T-14, paquete 1.4.1.2. Evidencia: Scripts de migración versionados y reversibles. \\
RT-04.11 Obligatorio & La cobertura de pruebas automatizadas del código de lógica de negocio será de al menos 70 \%, con umbral bloqueante en el flujo de integración continua. & Cumple & Puertas de calidad en integración continua con umbral bloqueante de cobertura de líneas de 70 \% sobre la lógica de negocio & Subdoc. 9, 9.1.4 y 9.1.4.1 (Tabla \refext{S9}{tab:umbrales-bloqueantes-calidad-codigo}); Subdoc. 3, 3.2.2.4; Formulario T-13; Formulario T-14, paquetes 1.1.4.2 y 1.5.1.1. Evidencia: Reporte de cobertura de al menos 70 \% en cada integración (RNF-TRV-04.11) \\
RT-04.12 Obligatorio & El Proponente declarará su frecuencia de despliegue objetivo, su tiempo desde el compromiso de código hasta producción, su tasa de cambios fallidos y su tiempo de restauración, y los medirá durante la Operación. & Cumple & Informe mensual de nivel de servicio & Subdoc. 3, 3.2.2.5; Subdoc. 6, 6.2.2.6 (Tabla \refext{T10}{tab:metas-flujo-entrega} del Formulario T-10: frecuencia de despliegue, tiempo desde el cambio de código hasta producción y tasa de cambios fallidos); Formulario T-14, paquetes 1.4.1.2 y 1.8.7.1. Evidencia: Métricas medidas durante la Operación (1.4.1.2 y 1.8.7.1). \\
RT-04.13 Deseable & Los ambientes no productivos se apagarán o reducirán fuera del horario de uso, con el ahorro reflejado en la estructura de costos. & Cumple parcialmente & Desarrollo, QA y Preproducción se apagan o reducen fuera del horario de uso; ahorro no reflejado en la estructura de costos & Subdoc. 4, 4.2.3.1 (Desarrollo: servicios que se apagan fuera del horario de uso); Formulario T-14, paquete 1.6.1.1. Falta: ahorro reflejado en la estructura de costos (Oferta Económica). Evidencia: Programa de apagado de ambientes no productivos (1.6.1.1). \\
RT-04.14 Deseable & Se valorará la existencia de ambientes efímeros por rama o por incidencia, creados y destruidos automáticamente. & Cumple parcialmente & Ambiente efímero de pruebas por rama, creado desde OpenTofu, con infraestructura reproducible & Subdoc. 6, 6.2.2.3 (infraestructura reproducible); Formulario T-14, paquete 1.4.1.2. Falta: destrucción automática del ambiente al cerrar la rama. Evidencia: Ambiente efímero de pruebas por rama, creado desde OpenTofu. \\
\end{tabla}
\begin{tabla}{I{0.132}I{0.263}I{0.109}I{0.183}I{0.313}}{Capítulo 05. Datos, integración e interoperabilidad.}{tab:capitulo-05-datos-integracion-interoperabilidad}
  {\cab{ID requerimiento} & \cab{Descripción} & \cab{Estado} & \cab{Componente que lo satisface} & \cab{Sección de la propuesta y evidencia}}
  {\emph{Fuente:} elaboración propia a partir de las Bases técnicas transversales.}
RT-05.01 Obligatorio & El Proponente entregará el modelo de datos documentado y un diccionario de datos con el nombre, el tipo, el dominio de valores, la obligatoriedad, el propietario y la sensibilidad de cada atributo. & Cumple & Modelo de datos por nueve dominios y diccionario de datos con atributo, tipo, dominio de valores y nulidad, propietario y sensibilidad & Subdoc. 5, 5.1; Subdoc. 5, 5.1.6; Subdoc. 5, Anexos (A.3 Atributos por entidad); Formulario T-14, paquetes 1.3.2.1 y 1.3.2.2. Evidencia: Modelo de datos y diccionario de datos \\
RT-05.02 Obligatorio & El Proponente justificará la selección del paradigma y del motor de persistencia: relacional o no relacional, garantías transaccionales, y la posición escogida entre consistencia y disponibilidad conforme al teorema CAP, para cada dominio de datos. & Cumple & Motor y paradigma de persistencia por dominio (PostgreSQL, series de tiempo, analítica, documental) con posición CAP por dominio y sitio & Subdoc. 5, 5.2.1 (Tabla \refext{S5}{tab:naturaleza-dato-motor-razon-eleccion} Tabla \refext{S5}{tab:posicion-cap-dominio}); Subdoc. 4, 4.1.2.3 (ADR-08); Formulario T-14, paquetes 1.3.1.1 y 1.3.2.1. Evidencia: Justificación de motor y consistencia por dominio \\
RT-05.03 Obligatorio & Toda operación de negocio será trazable: la solución permitirá reconstruir quién, qué, cuándo, desde qué dispositivo y con qué valores anteriores y posteriores, para cualquier registro y en cualquier momento del período de retención. & Cumple & Módulo Administración y Servicios Comunes (auditoría inalterable) & Subdoc. 3, 3.4.3 y 3.3.6 (autor de cada registro); Formulario T-14, paquete 1.4.3.1. Evidencia: Auditoría sobre muestra de registros: quién, cuándo, dispositivo, valores antes y después (RF-ADM-005) \\
RT-05.04 Obligatorio & La calidad de datos se gestionará conforme a ISO/IEC 25012, con validación en el punto de captura, indicadores de completitud, exactitud y consistencia, y tablero de calidad disponible para el Cliente. & Cumple & Validación de eventos en el dispositivo, indicadores ISO/IEC 25012 y tablero de calidad de datos de acceso propio del Cliente & Subdoc. 5, 5.2.3 (Tabla \refext{S5}{tab:indicadores-calidad-datos-metas}); Subdoc. 5, 5.2.2 (Tabla \refext{S5}{tab:vias-explotacion-informacion}); Formulario T-14, paquetes 1.4.3.6 y 1.4.4.6. Evidencia: Tablero de calidad de datos con metas 95/98/98 \% (RF-DAT-001) \\
RT-05.05 Obligatorio & El almacenamiento transaccional y el analítico estarán separados. Ninguna consulta analítica podrá degradar el desempeño de la operación. & Cumple & Almacenes transaccionales por módulo y modelo analítico (ClickHouse) separados, con prueba de carga analítica concurrente & Subdoc. 5, 5.2.2; Subdoc. 5, 5.4.6; Subdoc. 4, 4.1.2.3 (ADR-08); Formulario T-14, paquetes 1.3.2.1 y 1.5.3.1. Evidencia: Prueba de carga concurrente de operación y analítica (RNF-DES-017) \\
RT-05.06 Obligatorio & La solución permitirá exportar la totalidad de la información del Cliente en formatos abiertos y documentados, en cualquier momento del Contrato, sin costo adicional y sin intervención del Adjudicatario. & Cumple & Vía de exportación total del Cliente en formatos abiertos y documentados, sin costo y sin intervención del adjudicatario & Subdoc. 5, 5.2.2 (Tabla \refext{S5}{tab:vias-explotacion-informacion}); Subdoc. 5, 5.2.5; Subdoc. 2, Anexos (RNF-POR-002); Formulario T-14, paquetes 1.1.9.1, 1.4.3.1 y 1.8.10.1. Evidencia: Prueba de exportación y lectura por un tercero (RNF-POR-002) \\
RT-05.07 Obligatorio & El Proponente declarará la política de retención, archivado y eliminación por cada dominio de datos, coherente con la normativa aplicable al caso, e implementará un procedimiento verificable de eliminación segura. & Cumple & Política de retención por categoría de registro y procedimiento de eliminación segura en cinco pasos con certificado & Subdoc. 5, 5.2.4 (Tabla \refext{S5}{tab:retencion-categoria-entidad}); Subdoc. 5, 5.2.5 (Tabla \refext{S5}{tab:procedimiento-eliminacion-segura-evidencia-cada-paso}); Formulario T-14, paquete 1.3.2.3. Evidencia: Política de retención por dominio y procedimiento de eliminación segura \\
RT-05.08 Obligatorio & Los datos personales se tratarán conforme al Artículo 85° de las Bases administrativas, con seudonimización o cifrado a nivel de campo para las categorías sensibles que el caso identifique. & Cumple & Cifrado a nivel de campo con clave por categoría, seudonimización en la capa analítica y enmascarado irreversible fuera de producción & Subdoc. 5, 5.2.6; Subdoc. 4, 4.1.8.4 (Tabla \refext{S4}{tab:mecanismo-cifrado-custodia-claves-ambito}); Formulario T-14, paquetes 1.3.2.2, 1.3.3.1 y 1.3.6.2. Evidencia: Revisión con el área legal del Cliente (RNF-CUM-001) \\
RT-05.09 Obligatorio & El Proponente presentará una estrategia de gestión de datos maestros que evite la duplicación de entidades compartidas entre módulos y con sistemas externos. & Cumple & Regla de precedencia de datos maestros: sistema que manda por entidad y operaciones prohibidas, sin herramienta de datos maestros & Subdoc. 5, 5.1.4 (Tabla \refext{S5}{tab:sistema-maestro-operacion-prohibida-entidad}); Formulario T-14, paquete 1.3.2.1. Evidencia: Documento de gestión de datos maestros \\
RT-05.10 Deseable & Se valorará la implementación de un catálogo de datos con linaje automatizado, que permita rastrear el origen de cada indicador de negocio hasta su fuente. & Cumple & Catálogo de interfaces más diccionario con linaje automático desde el indicador hasta el sistema de origen & Subdoc. 5, 5.2.3; Subdoc. 4, 4.1.7.4; Formulario T-14, paquete 1.3.2.2. Evidencia: Catálogo de datos con linaje \\
RT-05.11 Obligatorio & El Proponente presentará un plan de migración con alcance, origen, volumen, reglas de transformación, criterios de calidad, estrategia de ejecución y plan de reversión. & Cumple & Plan de migración en dos etapas por origen, con volumen, reglas, criterios de aceptación y reversa & Subdoc. 5, 5.3 y 5.3.1 (Tabla \refext{S5}{tab:volumen-migracion-origen}; cálculo en la Tabla C.1 del Anexo C de OVERFORE-Subdocumento5-Anexos); Subdoc. 5, 5.3.2; Subdoc. 5, 5.3.4; Formulario T-14, paquetes 1.4.6.1, 1.4.6.2 y 1.4.6.4. Evidencia: Plan de migración de datos (TRA-MIG-001) \\
RT-05.12 Obligatorio & La migración incluirá una etapa de perfilado y saneamiento previo, con informe de los defectos detectados en los datos de origen y la decisión adoptada sobre cada uno. & Cumple & Perfilado y saneamiento previos de los datos de origen, con informe de defectos y decisión por cada uno & Subdoc. 5, 5.3.3; Formulario T-14, paquete 1.4.6.1. Evidencia: Informe de perfilado y saneamiento (TRA-MIG-002) \\
RT-05.13 Obligatorio & Se ejecutarán al menos dos ensayos completos de migración sobre Preproducción antes de la migración definitiva, con medición del tiempo total y del resultado de la conciliación. & Cumple & Al menos dos ensayos completos de migración por etapa en Preproducción con datos enmascarados, antes de la carga definitiva & Subdoc. 5, 5.3.3; Subdoc. 9, 9.3.1; Formulario T-14, paquetes 1.4.6.2 y 1.4.6.4. Evidencia: Bitácoras de dos ensayos completos en Preproducción (TRA-MIG-003) \\
RT-05.14 Obligatorio & La conciliación posterior a la migración será cuantitativa y verificable: recuentos, sumas de control y muestreo dirigido. Toda diferencia deberá quedar explicada. & Cumple & Conciliación cuantitativa por entidad y período con sumas de control sobre el 100 \% y muestreo dirigido de al menos 200 registros, con acta firmada & Subdoc. 5, 5.3.4; Formulario T-14, paquetes 1.4.6.2 y 1.4.6.4. Evidencia: Recuentos, sumas de control y muestreo dirigido (TRA-MIG-004) \\
RT-05.15 Según caso & Los datos históricos que no se migren quedarán accesibles en un repositorio de consulta durante el período de retención que fije el caso. & Cumple & Repositorio de consulta histórica de solo lectura por módulo, con origen, auditoría y control de acceso, habilitado en la Etapa 1 y completado en la Etapa 2 antes del cierre del mes 20 & Subdoc. 5, 5.3.5 y 5.2.1 (Tabla \refext{S5}{tab:naturaleza-dato-motor-razon-eleccion}); Formulario T-14, paquete 1.4.6.5 (meses 10 a 20). Evidencia: Prueba de consulta de registros anteriores a la ventana de migración (RF-ADM-001) \\
RT-05.16 Obligatorio & Los servicios síncronos se documentarán en OpenAPI 3.1 y los flujos dirigidos por eventos en AsyncAPI 2.6 o superior. La documentación se generará desde el código y se mantendrá actualizada automáticamente. & Cumple & Documentación OpenAPI 3.1 y AsyncAPI 3.0 generada desde el código (springdoc-openapi y Springwolf) y comparada en cada compilación & Subdoc. 4, 4.1.7.4; Subdoc. 4, 4.1.1.2 (Tabla \refext{S4}{tab:tecnologias-capas-4-5-alternativa-evaluada}); Formulario T-14, paquete 1.4.1.4. Evidencia: Documentación OpenAPI 3.1 y AsyncAPI 2.6 o superior generada desde el código (RNF-INT-006) \\
RT-05.17 Obligatorio & Los contratos de interfaz se versionarán semánticamente, con compatibilidad hacia atrás y política de obsolescencia con preaviso mínimo de seis meses. & Cumple & Versionado semántico de contratos con compatibilidad hacia atrás, retiro tras seis meses de aviso y gobierno por el Comité de Arquitectura & Subdoc. 4, 4.1.7.5; Subdoc. 4, 4.1.7.6; Subdoc. 4, 4.1.7.3 (Tabla \refext{S4}{tab:mecanismo-cifrado-custodia-claves-ambito}); Formulario T-14, paquete 1.4.1.4. Evidencia: Política de versionado semántico y obsolescencia con aviso de seis meses \\
RT-05.18 Obligatorio & La autenticación entre sistemas empleará OAuth 2.1 con credenciales de cliente o autenticación mutua TLS. Queda prohibida la autenticación por clave estática en la ruta de la dirección web. & Cumple & Puerta de exposición con OAuth 2.1 (credenciales de cliente o TLS mutuo); certificado del dispositivo en la puerta operacional; sin claves fijas en la dirección web & Subdoc. 4, 4.1.4.3 (Tabla \refext{S4}{tab:comparacion-aplicacion-web-progresiva-nativa-hibrida}); Subdoc. 4, 4.1.8.2 (Tabla \refext{S4}{tab:mecanismo-identidad-cada-poblacion-usuaria}); Subdoc. 4, 4.1.7.4; Subdoc. 2, Anexos (RNF-SEG-007); Formulario T-14, paquetes 1.4.2.1 y 1.6.1.1. Evidencia: Revisión de cada integración: OAuth 2.1 o TLS mutuo (RNF-SEG-007) \\
RT-05.19 Obligatorio & Toda integración registrará la transacción de entrada y de salida, con identificador de correlación común que permita seguir una operación de negocio a través de todos los sistemas involucrados. & Cumple & Integrador de Sistemas (un conector por sistema); identificador común de mensaje & Subdoc. 3, 3.4.3; Formulario T-14, paquete 1.4.2.1. Evidencia: Seguimiento de una operación de punta a punta con el identificador único (RNF-OPE-006) \\
RT-05.20 Obligatorio & Las integraciones con sistemas heredados o de terceros se aislarán mediante una capa anticorrupción, de modo que un cambio en el sistema externo no propague su modelo al núcleo de la solución. & Cumple & Integrador de Sistemas (un conector por sistema) (aísla a los módulos de los cambios externos) & Subdoc. 3, 3.3.2.2 y 3.4.3; Formulario T-14, paquete 1.4.2.1. Evidencia: Revisión funcional de conectores y prueba de cambio de fuente (RF-EQU-003) \\
RT-05.21 Obligatorio & El Proponente declarará, por cada integración, el modo (síncrono o asíncrono), el volumen esperado, la ventana de disponibilidad del sistema contraparte y el comportamiento de la solución cuando ese sistema no responde. & Cumple & Integrador de Sistemas (un conector por sistema); sistemas existentes & Subdoc. 3, 3.3.2.2, Tabla \refext{S3}{tab:sistemas-existentes-relacion-sitco-aranda}; Subdoc. 4, 4.1.7.1, Tablas \refext{S4}{tab:conectores-integrador-sistemas-etapa-lugar-ejecucion} y \refext{S4}{tab:modo-ventana-comportamiento-ante-falla-integraciones}; Formulario T-14, paquete 1.2.1.2. Evidencia: Matriz de integraciones con modo, dirección, volumen, frecuencia y ventana (1.2.1.2). \\
RT-05.22 Obligatorio & La solución soportará la carga y descarga masiva de información en formatos abiertos, con validación previa, informe de errores por registro y procesamiento parcial. & Cumple parcialmente & Módulo Administración y Servicios Comunes (trabajo de importación o exportación, RF-DAT-003) & Subdoc. 5, Anexo A.3 (registro Trabajo de importación o exportación, con errores por fila y procesamiento parcial); Subdoc. 2, Anexo A (RF-DAT-003); Formulario T-14, paquete 1.4.3.1 (validación previa y exportación en formatos abiertos). Falta: formato abierto de la carga masiva. Evidencia: Prueba de carga masiva con informe de errores por registro (RF-DAT-003). \\
RT-05.23 Según caso & Se emplearán los estándares sectoriales de intercambio que las Bases técnicas del caso Minería identifiquen. & Cumple & OPC UA (historiador), ISA-95, ISO 15143-3 (telemetría) y Protocolo GHG (huella) & Subdoc. 3, 3.3.4, 3.4.3, 3.4.2.4 y 3.4.2.7; Formulario T-14, paquetes 1.2.4.1, 1.3.1.3, 1.4.2.1 y 1.4.4.3. Evidencia: Revisión del modelo de integración y mapeo de señales por fabricante (RNF-INT-001, RNF-INT-005) \\
RT-05.24 Deseable & Se valorará la publicación de un portal de servicios para desarrolladores con documentación navegable, ambiente de pruebas y credenciales de prueba autoservidas. & Cumple parcialmente & Portal para desarrolladores de la puerta de exposición (Azure API Management) con ambiente de pruebas de datos sintéticos & Subdoc. 4, 4.1.7.4; Subdoc. 4, Anexo A (Puerta de exposición); Formulario T-14, paquete 1.6.1.1. Falta: credenciales de prueba autoservidas. Evidencia: Demostración del portal de desarrolladores con documentación navegable y ambiente de pruebas (Subdoc. 4, 4.1.7.4). \\
RT-05.25 Obligatorio & La solución proveerá una capa analítica con tableros operacionales y de gestión, construidos sobre los indicadores que las Bases técnicas del caso Minería definan. & Cumple & Módulo Indicadores y Centro Integrado (tableros de turno y de gestión) & Subdoc. 3, 3.4.3 y 3.2.2.2; Formulario T-14, paquetes 1.4.3.6 y 1.4.4.6. Evidencia: Tableros disponibles en la marcha blanca de la Etapa 1 y de la Etapa 2 (RF-ANA-001) \\
RT-05.26 Obligatorio & Los tableros permitirán filtrar por período, unidad organizacional y dimensiones propias del caso, y profundizar desde el indicador agregado hasta la transacción de origen. & Cumple & Módulo Indicadores y Centro Integrado (bajada hasta la transacción) & Subdoc. 3, 3.4.3; Formulario T-14, paquetes 1.4.3.6 y 1.4.4.6. Evidencia: Prueba de navegación del indicador a la transacción de origen (RF-ANA-001, RF-CON-014) \\
RT-05.27 Obligatorio & El Cliente podrá construir sus propios informes sin intervención del Adjudicatario, mediante una herramienta de autoservicio con modelo semántico documentado. & Cumple & Módulo Indicadores y Centro Integrado: autoservicio de informes sobre un modelo semántico documentado, en la Etapa 2 & Subdoc. 3, 3.2.2.3 y 3.4.3; Formulario T-14, paquete 1.4.4.6 (meses 15 a 17). Evidencia: Construcción de un informe por el Cliente sin apoyo del Adjudicatario (RF-ANA-002) \\
RT-05.28 Obligatorio & Todo informe será exportable en formatos abiertos y programable para envío automático por calendario. & Cumple & Módulo Indicadores y Centro Integrado (informes programados y autoservicio de informes) & Subdoc. 5, 5.2.2, Tabla \refext{S5}{tab:vias-explotacion-informacion}; Subdoc. 5, Anexo A.3 (registro Preferencia e informe programado, con calendario); Subdoc. 3, 3.4.3; Subdoc. 2, Anexo A (RF-ANA-003); Formulario T-14, paquetes 1.4.3.6 y 1.4.4.6. Evidencia: Prueba de exportación y programación (RF-ANA-003) \\
RT-05.29 Según caso & La latencia máxima entre la ocurrencia de una transacción y su disponibilidad en la capa analítica será la que fije el caso y, en su defecto, no superará las 4 horas. & Cumple parcialmente & Módulo Indicadores y Centro Integrado (capa analítica) & Subdoc. 5, 5.2.2, Tablas \refext{S5}{tab:vias-explotacion-informacion} y \refext{S5}{tab:indicadores-calidad-datos-metas} (15 minutos en tableros de turno y 4 horas en gestión); Subdoc. 3, 3.4.4, Tabla \refext{S3}{tab:correspondencia-entre-problema-solucion}; Subdoc. 2, Anexo A (RNF-DES-010); Formulario T-14, paquetes 1.4.3.2, 1.4.3.6, 1.4.4.1 y 1.4.4.6. Falta: latencia de 60 segundos para las señales. Evidencia: Comparación de marca de tiempo del evento y de su disponibilidad en la capa analítica (RNF-DES-010). \\
RT-05.30 Deseable & Se valorará la incorporación de analítica predictiva pertinente al proceso del caso, con el modelo, sus variables, su métrica de desempeño y su plan de reentrenamiento documentados. & Cumple & Clasificador asistido de detenciones (innovación 2 del Subdocumento 13), uno de los dos componentes de inteligencia artificial, junto con la verificación facial de portería, con modelo, variables, métrica y plan de reentrenamiento documentados & Subdoc. 13, 13.2.2 y 13.2.3; OVERFORE-Subdocumento13-Anexos, Anexo B, B.4 a B.6; Formulario T-14, paquetes 1.7.2.1, 1.7.2.2 y 1.7.2.3. Evidencia: Ficha de la innovación 2 (Formulario T-19) e informe de verificación \\
\end{tabla}
\begin{tabla}{I{0.132}I{0.266}I{0.109}I{0.171}I{0.323}}{Capítulo 06. Site principal on-premise.}{tab:capitulo-06-site-principal-on-premise}
  {\cab{ID requerimiento} & \cab{Descripción} & \cab{Estado} & \cab{Componente que lo satisface} & \cab{Sección de la propuesta y evidencia}}
  {\emph{Fuente:} elaboración propia a partir de las Bases técnicas transversales.}
RT-06.01 Obligatorio & El espacio asignado será de uso exclusivo de la solución y estará aislado de otras dependencias del Cliente, con acceso independiente. & Cumple & Sala técnica nueva prefabricada en faena & Subdoc. 4, 4.3.1 y 4.3.1.1 (Figura \refext{S4}{fig:planta-red-control-lectura-solo-sentido}); Subdoc. 4, Anexo A (sala técnica principal); Subdoc. 3, 3.3.5; Formulario T-14, paquete 1.6.1.2. Evidencia: Especificación del recinto y planos (Formulario T-11) \\
RT-06.02 Obligatorio & Los muros no estructurales del recinto contarán con blindaje perimetral; el Proponente especificará el material y la resistencia. & Cumple & Muros del módulo de paneles de acero con núcleo de lana mineral y 120 min de resistencia al fuego & Subdoc. 4, 4.3.1.1 (material y resistencia de muros); Formulario T-11 (T11-29); Formulario T-14, paquete 1.6.1.2. Evidencia: Especificación del recinto prefabricado con material y resistencia de muros (Formulario T-11, T11-29) \\
RT-06.03 Obligatorio & El Proponente entregará el plano de distribución interna del recinto, con la separación de las zonas de generadores, baterías, climatización, servidores, comunicaciones, trabajo y respaldo. & Cumple & Plano a escala 1:50 con esclusa, zona de trabajo técnico, sala de equipos (racks de comunicaciones y de servidores separados, climatización en fila) y zona de energía (UPS y baterías) tras muro cortafuego; generadores y tablero en base exterior; respaldo en la sala de 2014 & Subdoc. 4, 4.3.1.1 (Figura \refext{S4}{fig:planta-red-control-lectura-solo-sentido}; zonas de esclusa, trabajo técnico, equipos y energía; generadores en base exterior; respaldo en la sala de 2014); Formulario T-11 (T11-29); Formulario T-14, paquete 1.6.1.2. Evidencia: Plano de distribución interna del recinto (Subdoc. 4, Figura \refext{S4}{fig:planta-red-control-lectura-solo-sentido}) \\
RT-06.04 Obligatorio & El piso técnico, la canalización, el cableado estructurado y el etiquetado se ejecutarán conforme a norma, con documentación de la certificación de cada enlace. & Cumple & Sala técnica nueva prefabricada (piso técnico, canalización y cableado estructurado) & Subdoc. 4, 4.3.1.1 (piso técnico elevado de 30 cm, bandejas separadas de energía y datos, cableado certificado enlace por enlace con informe al cliente); Subdoc. 4, Anexo A (fibra certificada y etiquetada en ambos extremos); Formulario T-14, paquete 1.6.1.2. Evidencia: cableado estructurado conforme a ANSI/TIA-568 e ISO/IEC 11801, certificado enlace por enlace, con informe entregado al Cliente (Subdoc. 4, 4.3.1.1; Formulario T-11, T11-29). \\
RT-06.05 Obligatorio & Los racks de servidores serán independientes de los racks de equipos de comunicación. Se declarará la ocupación proyectada de cada rack y su margen de crecimiento. & Cumple & Sala técnica nueva prefabricada (racks RK-1 de comunicaciones y RK-2 de servidores) & Subdoc. 4, 4.3.1.1 (ocupación de 36 \% y 24 \%, tercera posición reservada); Subdoc. 4, Anexo A (racks RK-1 y RK-2); Formulario T-14, paquete 1.6.1.2. Evidencia: Plano de la sala y ocupación por rack (Formulario T-11) \\
RT-06.06 Obligatorio & La obra civil de separación de las instalaciones es de cargo del Cliente; su especificación técnica y su coordinación son de cargo del Proponente. & Cumple & Especificación y coordinación de la obra civil por Overfore; ejecución por el Cliente & Subdoc. 3, 3.2.3.1 y 3.3.5; Formulario T-14, paquetes 1.1.6.2 y 1.6.1.2. Evidencia: Especificación técnica de la obra civil y acta de coordinación \\
RT-06.07 Obligatorio & El suministro eléctrico de los equipos será ininterrumpido, con sistema de alimentación ininterrumpida dimensionado para una autonomía mínima de 30 minutos a plena carga. & Cumple & Sala técnica nueva prefabricada en faena (respaldo eléctrico: dos UPS de 10 kVA) & Subdoc. 4, 4.3.1.3; Subdoc. 4, Anexo A (UPS de doble conversión); Subdoc. 4, Anexo D.3; Formulario T-14, paquete 1.6.1.2. Evidencia: Especificación eléctrica (Formulario T-11) \\
RT-06.08 Obligatorio & La capacidad de generación autónoma asegurará un rango mínimo de 24 horas continuas de operación, con estanque de combustible dimensionado y contrato de reabastecimiento declarado. & Cumple & Sala técnica nueva prefabricada en faena (dos generadores de 20 kVA en N+1 y estanque de 1.000 L) & Subdoc. 4, 4.3.1.3; Subdoc. 4, Anexo A (generadores y estanque); Subdoc. 4, Anexo C (contrato de reabastecimiento a cargo del Cliente); Subdoc. 4, Anexo D.3 (24 h = 53 L); Formulario T-14, paquete 1.6.1.2. Evidencia: Especificación del generador y contrato de combustible \\
RT-06.09 Obligatorio & La instalación eléctrica del recinto será independiente de la del resto del edificio y cumplirá la normativa eléctrica chilena vigente, incluida la NCh Elec. 2777 sobre sistemas de puesta a tierra. & Cumple & Sala técnica nueva prefabricada en faena (tablero y puesta a tierra independientes) & Subdoc. 4, 4.3.1.3; Subdoc. 4, Anexo A (tablero de distribución, puesta a tierra según NCh Elec. 2777); Formulario T-14, paquete 1.6.1.2. Evidencia: Especificación eléctrica y protocolo de puesta a tierra (Formulario T-11) \\
RT-06.10 Obligatorio & Se efectuará revisión y medición semestral de las instalaciones eléctricas del recinto, con informe entregable al Cliente. & Cumple & Sala técnica nueva prefabricada (instalación eléctrica con revisión semestral) & Subdoc. 4, 4.3.1.3; Subdoc. 4, Anexo A (tablero, medición semestral); Subdoc. 4, Anexo C (prueba mensual bajo carga y medición semestral); Formulario T-14, paquetes 1.6.1.2 y 1.8.4.2. Evidencia: Informe semestral de revisión eléctrica \\
RT-06.11 Obligatorio & El Proponente declarará la carga eléctrica proyectada en kW, el factor de potencia y la eficiencia en el uso de la energía (PUE) estimada del recinto. & Cumple & Sala técnica nueva prefabricada (dimensionamiento eléctrico: 3,65 kW de TI, 5,84 kW totales) & Subdoc. 4, 4.3.1.2 y 4.3.1.3; Subdoc. 4, Anexo D.3 (carga, factor de potencia 0,9 y PUE 1,6); Formulario T-14, paquete 1.6.1.2. Evidencia: Memoria de dimensionamiento eléctrico (Subdoc. 4, Anexo D.3) \\
RT-06.12 Deseable & Se valorará la redundancia de alimentación en configuración 2N o N+1 con doble acometida y transferencia automática. & Cumple parcialmente & Sala técnica nueva prefabricada (generadores N+1 y tablero de transferencia automática) & Subdoc. 4, 4.3.1.3; Subdoc. 4, Anexo A (generadores y tablero de transferencia); Formulario T-14, paquete 1.6.1.2 (generador de 24 horas). Falta: doble acometida de la red eléctrica. Evidencia: Prueba mensual de generadores y de transferencia bajo carga (Subdoc. 4, 4.3.1.3). \\
RT-06.13 Obligatorio & El recinto contará con climatización de precisión para operación continua, redundante en configuración N+1, con control de temperatura y de humedad relativa dentro de los rangos que recomienda el fabricante del equipamiento. & Cumple & Sala técnica nueva prefabricada en faena (dos unidades de precisión de 7 kW en N+1) & Subdoc. 4, 4.3.1.2; Subdoc. 4, Anexo A (climatización de precisión en fila); Formulario T-14, paquete 1.6.1.2. Evidencia: Especificación de climatización (Formulario T-11). \\
RT-06.14 Obligatorio & Se monitorearán en línea la temperatura, la humedad y la presencia de agua, con alertamiento integrado a la plataforma de observabilidad. & Cumple & Sensores ambientales de la sala técnica integrados a la consola operativa única & Subdoc. 4, 4.3.1.2; Subdoc. 4, 4.2.3.4; Subdoc. 4, Anexo A (sensores de temperatura, humedad y agua); Formulario T-14, paquete 1.6.1.2. Evidencia: Prueba de alerta de los sensores en la consola operativa única \\
RT-06.15 Deseable & El Proponente declarará la estrategia de contención de pasillo frío o caliente y su efecto en la eficiencia energética. & Cumple & Sala técnica nueva prefabricada (pasillo frío confinado) & Subdoc. 4, 4.3.1.2; Subdoc. 4, Anexo A (climatización con pasillo frío confinado); Formulario T-14, paquete 1.6.1.2. Evidencia: Memoria de climatización con PUE estimado de 1,6 (Subdoc. 4, 4.3.1.2) \\
RT-06.16 Obligatorio & El recinto contará con detección temprana por aspiración de aire con tecnología láser, tipo AnaLASER o equivalente de prestaciones iguales o superiores. & Cumple & Sala técnica nueva prefabricada en faena (detección automática de incendios por aspiración láser) & Subdoc. 4, 4.3.1.4; Subdoc. 4, Anexo A (detección temprana y extinción); Formulario T-14, paquete 1.6.1.2. Evidencia: Especificación del sistema de detección (Formulario T-11) \\
RT-06.17 Obligatorio & La extinción será automática mediante agente limpio tipo FM-200 o equivalente, con aprobación UL e instalación conforme a norma NFPA. & Cumple & Sala técnica nueva prefabricada en faena (extinción automática de incendios) & Subdoc. 4, 4.3.1.4; OVERFORE-Subdocumento4-Anexos, Anexo A (T11-39); Formulario T-14, paquete 1.6.1.2. Evidencia: Especificación del sistema de extinción (Formulario T-11, T11-39). \\
RT-06.18 Obligatorio & Se proveerá un sistema secundario de extintores portátiles habilitados, con mantención y certificación vigentes. & Cumple parcialmente & Extintores portátiles de CO₂ como sistema secundario & Subdoc. 4, 4.3.1.4; OVERFORE-Subdocumento4-Anexos, Anexo A (T11-39); Formulario T-14, paquete 1.8.4.2 (desde el mes 21). Falta: mantención y certificación de los extintores desde su instalación hasta el mes 21. Evidencia: Ficha de extinción con extintores de CO₂ (Formulario T-11, T11-39). \\
RT-06.19 Obligatorio & El sistema de detección y extinción se integrará al monitoreo en línea y notificará al NOC y a la contraparte del Cliente. & Cumple & Detección y extinción de incendios de la sala técnica integrada a la consola operativa única & Subdoc. 4, 4.3.1.4; Subdoc. 4, Anexo A (detección temprana y extinción); Formulario T-14, paquetes 1.6.1.2 y 1.8.2.1. Evidencia: Prueba semestral del sistema con notificación a la consola operativa y al Cliente (Subdoc. 4, 4.3.1.4) \\
RT-06.20 Obligatorio & El ingreso al recinto se controlará mediante seguridad física y control de acceso biométrico basado principalmente en biometría facial, con AFIS como respaldo. Se admite proponer sistemas de mayor seguridad. & Cumple & Control de acceso biométrico a la sala técnica, con biometría facial uno a uno y respaldo AFIS, y esclusa de una persona a la vez (RF-ACR-017) & Subdoc. 2, Anexo A.1 (RF-ACR-017) y Subdoc. 3, 3.3.5; Formulario T-14, paquetes 1.4.3.5 (control de acceso a la sala técnica, meses 4 a 10) y 1.3.6.2 (base legal y consentimiento declarados para el tratamiento biométrico). Evidencia: Prueba de aceptación del resultado esperado de RF-ACR-017 (Anexo A.1) y especificación del control de acceso (Formulario T-11) \\
RT-06.21 Obligatorio & Todo ingreso y egreso quedará registrado en una bitácora auditable, con identificación de la persona, fecha, hora y motivo, conservada por el período de retención declarado. & Cumple & Bitácora inalterable de ingresos y egresos a la sala técnica con identidad, fecha, hora y motivo (RF-ACR-018), conservada 12 meses en línea y 24 en archivo recuperable (RNF-SEG-011) & Subdoc. 2, Anexo A.1 (RF-ACR-018) y Anexo A.2 (RNF-SEG-011); Formulario T-14, paquete 1.4.3.5. Evidencia: Prueba de aceptación del resultado esperado de RF-ACR-018 (Anexo A.1) y revisión de la retención configurada (RNF-SEG-011) \\
RT-06.22 Obligatorio & Entre el acceso principal y el término del pasillo de la zona de control se dispondrá un espacio para la atención de personas en proceso de enrolamiento. Se evaluará mejor la existencia de una estación de enrolamiento fuera de las instalaciones del recinto técnico. & Cumple parcialmente & Sala de 2014 como punto de enrolamiento; esclusa biométrica de la sala técnica & Subdoc. 4, Anexo A (esclusa biométrica; enrolamiento en la sala de 2014); Subdoc. 4, 4.3.1.1; Formulario T-14, paquete 1.6.1.2 (sala de 2014 como punto de enrolamiento). Falta: declarar el espacio de atención para personas en proceso de enrolamiento entre el acceso principal y el término del pasillo de la zona de control. Evidencia: Plano de la sala con el espacio de enrolamiento (Formulario T-11) \\
RT-06.23 Obligatorio & Al término del pasillo se instalará un acceso que impida el paso de más de una persona a la vez, con nueva verificación de identidad previa al ingreso. & Cumple & Sala técnica nueva prefabricada en faena (esclusa biométrica de una persona) & Subdoc. 3, 3.3.5; Subdoc. 4, 4.3.1.1; Subdoc. 4, Anexo A (esclusa biométrica y videovigilancia); Formulario T-14, paquete 1.4.3.5. Evidencia: Especificación del control de acceso (Formulario T-11) \\
RT-06.24 Obligatorio & El recinto contará con videovigilancia y monitoreo IP, con imágenes en línea y disponibles para visualización de al menos los últimos 30 días. Las grabaciones anteriores se respaldarán en un medio secundario recuperable y auditable. & Cumple parcialmente & Videovigilancia IP de 30 días integrada a la consola única (esclusa de la sala técnica) & Subdoc. 4, 4.3.1.1; OVERFORE-Subdocumento4-Anexos, Anexo A (T11-40) y Anexo E (E.3, control 7.4); Formulario T-14, paquete 1.6.1.2. Falta: respaldo auditable de las grabaciones anteriores a 30 días. Evidencia: Especificación de la esclusa biométrica y videovigilancia (Formulario T-11, T11-40) \\
RT-06.25 Obligatorio & El Proponente declarará el procedimiento de acceso de terceros —fabricantes, mantenedores, auditores— con acompañamiento obligatorio y registro. & Cumple parcialmente & Sala técnica nueva prefabricada (bitácora de ingreso y acompañamiento obligatorio) & Subdoc. 4, 4.3.1.1; Subdoc. 4, Anexo A (esclusa biométrica y videovigilancia); Formulario T-14, paquetes 1.6.1.2 y 1.4.3.5. Falta: pasos del procedimiento de acceso de terceros. Evidencia: Procedimiento de acceso de terceros y bitácora de ingresos \\
RT-06.26 Obligatorio & Se habilitará un servicio de custodia de medios de respaldo para el sitio primario, en un medio físico transportable a otro lugar cuando el Cliente lo determine. Se admite proponer una solución más segura y eficiente, debidamente justificada. & Cumple & Sala de 2014 reconvertida en recinto de custodia de respaldos (cartuchos cifrados transportables) & Subdoc. 4, 4.2.1.2; Subdoc. 4, 4.2.3.5; Subdoc. 4, Anexo A (medios removibles cifrados); Formulario T-14, paquete 1.6.1.2. Evidencia: Especificación del servicio de custodia de medios \\
RT-06.27 Obligatorio & El recinto de custodia cumplirá exigencias de luminosidad, humedad, ventilación y cualquier otro factor que pueda afectar la calidad y la disponibilidad de los medios. & Por cumplir & Sala de 2014 como recinto de custodia, sin condiciones ambientales declaradas & Subdoc. 4, 4.2.1.2 y 4.2.3.5 (la sala de 2014 como recinto de custodia, con inventario y verificación de legibilidad); Formulario T-14, paquete 1.6.1.2. Falta: condiciones de luminosidad, humedad, ventilación y demás factores que afecten la calidad y la disponibilidad de los medios en el recinto de custodia. Evidencia: Especificación ambiental del recinto de custodia (Formulario T-11) \\
RT-06.28 Obligatorio & Se llevará un inventario de medios con rotación, verificación periódica de legibilidad y registro de todo movimiento de entrada y de salida. & Cumple & Recinto de custodia en la sala de 2014 (inventario y rotación de medios removibles) & Subdoc. 4, 4.2.3.5; Subdoc. 4, Anexo A (medios removibles cifrados); Formulario T-14, paquete 1.6.1.2. Evidencia: Inventario de medios y registro de entradas y salidas (Subdoc. 4, Anexo C) \\
RT-06.29 Obligatorio & El Proponente habilitará el espacio físico necesario para el personal encargado de la operación y administración de la plataforma, con estaciones de trabajo, telefonía, conexión a Internet y todo elemento que permita realizar la labor en condiciones adecuadas. & Cumple & Sala de 2014 reconvertida en espacio de trabajo del soporte en faena (cuatro estaciones de administración) & Subdoc. 4, 4.2.1.2; Subdoc. 4, 4.3 y 4.3.1.1; Subdoc. 4, Anexo A (estaciones con dos monitores); Formulario T-14, paquetes 1.6.1.2 y 1.6.2.3. Evidencia: Especificación de puestos de operación (Formulario T-11) \\
RT-06.30 Obligatorio & El espacio de operación estará separado de la sala de equipos y no requerirá el ingreso al recinto técnico para las labores habituales de operación. & Cumple & Sala de 2014 como espacio de trabajo, separada de la sala técnica nueva & Subdoc. 3, 3.3.5; Formulario T-14, paquete 1.6.1.2. Evidencia: Plano de distribución de la sala técnica y del espacio de operación (Subdoc. 4, 4.2) \\
RT-06.31 Obligatorio & Las instalaciones sanitarias, las zonas de seguridad ante emergencia y las áreas exteriores existentes en el edificio del Cliente podrán utilizarse y no deben implementarse nuevamente. El Proponente declarará de cuáles hará uso. & Cumple parcialmente & Uso de las instalaciones sanitarias, zonas de seguridad y áreas exteriores del edificio del Cliente, sin implementarlas de nuevo & Formulario T-14, paquete 1.6.1.2 (uso de las instalaciones del edificio del Cliente, sin implementarlas de nuevo). Falta: enumerar las instalaciones sanitarias, zonas de seguridad ante emergencia y áreas exteriores del edificio del Cliente que se usarán. Evidencia: Declaración de las instalaciones del edificio del Cliente que se usarán \\
RT-06.32 Obligatorio & El acceso a las redes de comunicaciones estará provisto a través de rutas físicas distintas, con ingreso al edificio por puntos separados. & Cumple & Segunda fibra por acometida y canalización separadas & Subdoc. 3, 3.3.4 y 3.3.5; Formulario T-14, paquete 1.6.1.3. Evidencia: Plano de acometidas y canalizaciones (Subdoc. 4, 4.2) \\
RT-06.33 Obligatorio & El Proponente proveerá toda la conectividad, la seguridad y las canalizaciones, cañerías o ductos requeridos para cumplir los niveles de servicio comprometidos. & Cumple parcialmente & Segundo enlace de fibra y SD-WAN; ruta ejecutada por el proveedor de telecomunicaciones & Subdoc. 3, 3.2.3.1 y 3.3.4; Subdoc. 4, 4.2.3.2 y 4.3.1.1; OVERFORE-Subdocumento4-Anexos, Anexo A (T11-64 y T11-29); Formulario T-14, paquetes 1.6.1.3 (ductos de fibra hasta el punto de conexión de cada sitio) y 1.1.6.2 (segunda fibra con proveedor de telecomunicaciones distinto del actual). Falta: canalizaciones de energía de la sala técnica, que son obra civil del Cliente (SUP-34). Evidencia: Especificación de conectividad y niveles de servicio (Formulario T-11, T11-64) \\
RT-06.34 Deseable & Se privilegiará al Proponente que ofrezca o provea especificaciones nuevas o mejores que las aquí establecidas, debidamente fundamentadas. & Cumple & Especificaciones superiores y fundamentadas: meta interna de disponibilidad de 99,95 \%, estanque de 1.000 L dimensionado para siete días con 50 \% de reserva y crecimiento a 600 TPS & Subdoc. 4, 4.2.3.4 (meta de 99,95 \%), 4.2.4.2 (crecimiento a 600 TPS, Tabla \refext{S4}{tab:declaracion-uso-ia-seccion}) y 4.3.1.3 (estanque dimensionado por la distancia); OVERFORE-Subdocumento4-Anexos, Anexo D (D.3.6); Formulario T-14, paquete 1.3.1.2. Evidencia: Fundamentos de dimensionamiento de la Tabla \refext{S4}{tab:volumetria-referencia-capacidad-diseno} (Subdoc. 4, 4.2.4.2) y memoria del Anexo D \\
\end{tabla}
\begin{tabla}{I{0.132}I{0.268}I{0.115}I{0.176}I{0.309}}{Capítulo 07. Site secundario y recuperación ante desastres.}{tab:capitulo-07-site-secundario-recuperacion-ante}
  {\cab{ID requerimiento} & \cab{Descripción} & \cab{Estado} & \cab{Componente que lo satisface} & \cab{Sección de la propuesta y evidencia}}
  {\emph{Fuente:} elaboración propia a partir de las Bases técnicas transversales.}
RT-07.01 Obligatorio & El Proponente declarará la modalidad escogida —activo-activo o activo-pasivo— y la justificará frente al costo, al RTO comprometido y a la complejidad operacional que introduce. & Cumple & Región secundaria Azure Brazil South en modalidad activo-pasivo con réplica continua & Subdoc. 4, 4.3.2 y 4.3.2.1; Subdoc. 4, Anexo A (región de recuperación); Formulario T-14, paquete 1.3.1.4. Evidencia: Documento de recuperación ante desastres (Subdoc. 4, 4.3.2) \\
RT-07.02 Obligatorio & El sitio secundario estará emplazado a una distancia suficiente del principal para no verse afectado por el mismo evento de fuerza mayor. El Proponente declarará la distancia y el análisis de amenazas comunes considerado. & Cumple & Regiones Santiago (Chile Central) y São Paulo (Brazil South) a unos 2.600 km y separadas por la cordillera & Subdoc. 3, 3.3.5; Subdoc. 4, 4.3.2.1; Formulario T-14, paquete 1.3.1.4. Evidencia: Análisis de amenazas comunes entre sitios (Subdoc. 4, 4.3.2.1) \\
RT-07.03 Obligatorio & La replicación de datos será continua, con medición y alertamiento del retraso de replicación. & Cumple & Réplica continua de datos a la región secundaria con monitor de retraso & Subdoc. 4, 4.3.2.1; Subdoc. 5, 5.2.1, Tabla \refext{S5}{tab:objetivos-recuperacion-mecanismo-datos-sostiene}; Subdoc. 2, Anexo A (RNF-TRV-07.03); Formulario T-14, paquete 1.3.1.4. Evidencia: Prueba de retraso inducido con alerta emitida (RNF-TRV-07.03) \\
RT-07.04 Obligatorio & El objetivo de tiempo de recuperación (RTO) no superará 4 horas y el objetivo de punto de recuperación (RPO) no superará 15 minutos para los servicios críticos, salvo exigencia superior del caso. & Cumple & Región secundaria São Paulo; replicación continua & Subdoc. 3, 3.2.2.6; Formulario T-14, paquetes 1.3.1.4 y 1.5.5.1. Evidencia: Prueba semestral de conmutación con medición de RTO y RPO (RNF-OPE-005) \\
RT-07.05 Obligatorio & El procedimiento de conmutación estará documentado, automatizado en la mayor medida posible y ejecutable por el personal del Cliente tras la transferencia de conocimiento. & Cumple & Procedimiento de conmutación a Brazil South, documentado, automatizado y ejecutable por el Cliente & Subdoc. 4, 4.3.2.2; Subdoc. 2, Anexo A (RP de transferencia de conocimiento); Formulario T-14, paquetes 1.3.1.4 y 1.5.5.1. Evidencia: Procedimiento de conmutación documentado y ejecutable por el Cliente \\
RT-07.06 Obligatorio & Existirá un procedimiento de retorno al sitio principal igualmente documentado y probado, con reconciliación de los datos generados durante la contingencia. & Cumple & Procedimiento de retorno a Chile Central documentado y probado, con reconciliación por reglas deterministas & Subdoc. 4, 4.3.2.2; Formulario T-14, paquetes 1.3.1.4, 1.5.5.1 y 1.8.8.1. Evidencia: Procedimiento de retorno al sitio principal probado (Subdoc. 4, 4.3.2.2) \\
RT-07.07 Obligatorio & El plan de recuperación ante desastres se probará al menos dos veces al año mediante conmutación real, con informe de resultados, medición del RTO y del RPO efectivamente alcanzados y plan de corrección de las brechas detectadas. & Cumple & Prueba semestral de conmutación real a la región secundaria & Subdoc. 3, 3.2.2.6; Formulario T-14, paquetes 1.5.5.1 y 1.8.8.1. Evidencia: Informe de resultados con RTO y RPO efectivos y plan de corrección de brechas \\
RT-07.08 Deseable & Se valorará que la conmutación sea automática ante la detección de indisponibilidad, con criterio de disparo declarado y protección contra conmutación innecesaria. & Cumple parcialmente & Procedimiento de conmutación con criterio de disparo y decisión del jefe de incidente & Subdoc. 4, 4.3.2.2; Formulario T-14, paquete 1.3.1.4. Falta: conmutación automática ante la detección de indisponibilidad; hoy la decisión la confirma el jefe del incidente. Evidencia: Prueba de conmutación real dos veces al año, en las detenciones de marzo y septiembre \\
RT-07.09 Obligatorio & La política de respaldo seguirá el esquema 3-2-1-1-0: tres copias, en dos medios distintos, una fuera de sitio, una inmutable o fuera de línea y cero errores de verificación de restauración. & Cumple & Política de respaldo 3-2-1-1-0 en la nube, el Nodo de Faena y la sala de 2014 & Subdoc. 4, 4.2.3.5; Subdoc. 4, Anexo A (respaldos); Formulario T-14, paquete 1.3.1.2. Evidencia: Política de respaldo 3-2-1-1-0 (RNF-OPE-003) \\
RT-07.10 Obligatorio & Los respaldos estarán cifrados en reposo y en tránsito, con clave gestionada de forma independiente de la infraestructura respaldada. & Cumple & Respaldos cifrados en reposo y en tránsito con claves en bóveda separada de lo respaldado & Subdoc. 4, 4.2.3.5; Subdoc. 4, 4.1.8.4 (Tabla \refext{S4}{tab:amenazas-mayor-consecuencia-componentes-expuestos-control}); OVERFORE-Subdocumento4-Anexos, Anexo A (bóveda de claves en HSM); Formulario T-14, paquetes 1.3.1.2 y 1.8.8.5. Evidencia: Revisión de cifrado y de gestión independiente de claves (RNF-OPE-003) \\
RT-07.11 Obligatorio & Las copias inmutables estarán protegidas contra borrado y contra modificación durante su período de retención, incluso frente a credenciales administrativas comprometidas. & Cumple & Copia inmutable en Brazil South y repositorio local inmutable de la sala de 2014 & Subdoc. 4, 4.2.3.5; Subdoc. 4, Anexo A (almacén de respaldo inmutable y repositorio local); Formulario T-14, paquetes 1.3.1.2 y 1.8.8.5. Evidencia: Prueba de intento de borrado de copias inmutables \\
RT-07.12 Obligatorio & Se ejecutará y documentará una prueba de restauración al menos mensual, sobre una muestra representativa, con medición del tiempo efectivo de restauración. & Cumple & Prueba mensual de restauración en Preproducción sobre una muestra de cada dominio & Subdoc. 4, 4.2.3.5; Subdoc. 4, 4.3.2.2; Subdoc. 4, Anexo A (prueba de restauración mensual); Formulario T-14, paquetes 1.3.1.2 y 1.8.8.5. Evidencia: Acta mensual de prueba de restauración (RNF-OPE-004) \\
RT-07.13 Obligatorio & El Proponente declarará, por cada dominio de datos, la frecuencia de respaldo, el período de retención y el tiempo estimado de restauración completa. & Cumple parcialmente & Tabla de respaldo por dominio de datos (transaccional, series de tiempo, documentos y analítica) & Subdoc. 4, 4.2.3.5; Formulario T-14, paquete 1.3.1.2. Evidencia: Tabla de frecuencia, retención y tiempo de restauración por dominio de datos \\
RT-07.14 Deseable & Los respaldos permitirán la restauración granular: un registro, una tabla, un módulo o el sistema completo. & Cumple & Restauración granular desde el respaldo (registro, tabla, módulo o sistema completo) & Subdoc. 4, 4.2.3.5; Subdoc. 4, Anexo A (respaldos); Formulario T-14, paquetes 1.3.1.2 y 1.8.8.5. Evidencia: Prueba de restauración de un registro, una tabla y un módulo \\
\end{tabla}
\begin{tabla}{I{0.132}I{0.269}I{0.115}I{0.176}I{0.308}}{Capítulo 08. Hardware, puestos de trabajo y terreno.}{tab:capitulo-08-hardware-puestos-trabajo-terreno}
  {\cab{ID requerimiento} & \cab{Descripción} & \cab{Estado} & \cab{Componente que lo satisface} & \cab{Sección de la propuesta y evidencia}}
  {\emph{Fuente:} elaboración propia a partir de las Bases técnicas transversales.}
RT-08.01 Obligatorio & El Proponente especificará el equipamiento de cómputo con su marca, modelo de referencia, procesador, memoria, almacenamiento local, interfaces y consumo, junto con el cálculo de dimensionamiento que lo sustenta. & Cumple parcialmente & Servidores del Nodo de Faena (Cisco UCS C240 M8), banco de pruebas y servidores de borde & Subdoc. 4, 4.2.1.2 y 4.2.4.3; OVERFORE-Subdocumento4-Anexos, Anexo A (T11-31, T11-32, T11-53) y Anexo D (D.3.6, Tabla D.9: 0,6 kW por servidor); Formulario T-14, paquete 1.3.1.2. Falta: modelo de procesador de cada equipo de cómputo. Evidencia: Especificación de equipamiento de cómputo y cálculo de dimensionamiento (Formulario T-11, T11-31 a T11-32) \\
RT-08.02 Obligatorio & El almacenamiento será redundante, con tolerancia declarada a la falla de discos, control de errores y monitoreo predictivo de salud de los medios. & Cumple parcialmente & Nodo de Faena (sala técnica, Etapa 1) (almacenamiento redundante) & Subdoc. 3, 3.3.5; Formulario T-14, paquete 1.6.1.2. Falta: control de errores y monitoreo predictivo de salud de los medios. Evidencia: Especificación de almacenamiento (Formulario T-11) \\
RT-08.03 Obligatorio & Los conmutadores de núcleo, los cortafuegos y los balanceadores de carga estarán en configuración de alta disponibilidad, sin punto único de falla. & Cumple & Nodo de Faena (sala técnica, Etapa 1) con equipos de red y seguridad duplicados & Subdoc. 4, 4.2.3.4; Subdoc. 4, Anexo A (cortafuegos en HA y conmutadores de núcleo cortafuegos industriales en HA); Formulario T-14, paquete 1.3.1.3. Evidencia: Diagrama de red sin punto único de falla (Subdoc. 4, 4.2.3.2) \\
RT-08.04 Obligatorio & Todo el equipamiento contará con fuentes de poder redundantes y conexión a circuitos eléctricos distintos. & Cumple parcialmente & Sala técnica con UPS A y B por líneas separadas y doble fuente en servidores y conmutadores de borde & Subdoc. 4, 4.3.1.3 (Figura \refext{S4}{fig:gabinete-borde-acopio-puerto}) y 4.2.1.2; OVERFORE-Subdocumento4-Anexos, Anexo A (T11-31, T11-37, T11-53 y T11-54) y Anexo E (E.3, control 7.8); Formulario T-14, paquete 1.6.1.2. Falta: circuitos eléctricos distintos en los gabinetes de relaves y de puerto, que el Subdoc. 4 alimenta sólo con su red local, y en las tabletas, botoneras y consolas de terreno. Evidencia: Diagrama de energía de respaldo de la sala técnica y de los gabinetes (Subdoc. 4, Figura \refext{S4}{fig:gabinete-borde-acopio-puerto}) \\
RT-08.05 Obligatorio & El Proponente declarará el margen de crecimiento del dimensionamiento propuesto, expresado como porcentaje sobre la carga proyectada del caso, y el procedimiento de ampliación. & Cumple & Dimensionamiento con factor de diseño de tres veces y umbrales de ampliación & Subdoc. 4, 4.2.4.2, Tabla \refext{S4}{tab:declaracion-uso-ia-seccion}; Subdoc. 4, 4.2.4.6; Subdoc. 4, Anexo D.3 (margen eléctrico de 30 \%); Formulario T-14, paquete 1.3.1.2. Evidencia: Memoria de dimensionamiento y umbrales de ampliación (Subdoc. 4, 4.2.4) \\
RT-08.06 Obligatorio & El equipamiento será nuevo, sin uso previo, con garantía de fábrica vigente desde la recepción conforme. & Cumple & Equipamiento nuevo de faena con soporte 24x7 de cinco años & Subdoc. 4, 4.2.1.4; Subdoc. 4, Anexo C (equipamiento nuevo); Formulario T-14, paquete 1.1.6.3. Evidencia: Inventario con seriales y garantía vigente en la recepción (Subdoc. 4, Anexo C) \\
RT-08.07 Obligatorio & El Proponente especificará las estaciones de trabajo requeridas para la operación y la administración de la plataforma, en la cantidad que determine el dimensionamiento del caso, con monitores duales. & Cumple & Siete estaciones de administración con dos monitores (cuatro en la faena y tres en la mesa de ayuda) & Subdoc. 4, Anexo A (estación de administración con dos monitores) y Anexo D; Formulario T-14, paquete 1.6.2.3 (instaladas en los meses 11 y 12). Evidencia: Especificación de estaciones de operación (Formulario T-11) \\
RT-08.08 Obligatorio & Los puestos de trabajo cumplirán las condiciones ergonómicas de la NCh 2527 y los equipos contarán con certificación de eficiencia energética. & Cumple parcialmente & Estaciones de administración especificadas con NCh 2527 & OVERFORE-Subdocumento4-Anexos, Anexo A (T11-48: estaciones con dos monitores, TPM, EDR y NCh 2527); Formulario T-14, paquete 1.6.2.3. Falta: condiciones ergonómicas NCh 2527 y certificación de eficiencia energética para los puestos y equipos distintos de las siete estaciones. Evidencia: Ficha de estaciones con certificados de ergonomía NCh 2527 y de eficiencia energética (Formulario T-11, T11-48) \\
RT-08.09 Obligatorio & Las estaciones estarán gestionadas de forma centralizada, con cifrado de disco, control de dispositivos extraíbles, antivirus con detección y respuesta y actualización automatizada. & Cumple & Estaciones de administración gestionadas con Microsoft Intune, con TPM y EDR & Subdoc. 4, 4.2.1.1 y Anexo A (estaciones con TPM y EDR); Formulario T-14, paquete 1.6.2.3 (informe de cumplimiento de políticas antes del mes 13). Evidencia: Informe de cumplimiento de políticas de las estaciones en Intune \\
RT-08.10 Obligatorio & El Proponente especificará cada dispositivo de terreno con marca, modelo de referencia, cantidad, características mínimas, accesorios, consumibles y costo unitario estimado, aun cuando su compra sea de cargo del Cliente. & Cumple parcialmente & Cinco dispositivos nuevos: Botonera de Pala, tabletas rugerizadas, PC industriales de la Consola de Portería, lectores faciales y caudalímetros & Subdoc. 3, 3.3.3; Subdoc. 4, 4.2.1.3; OVERFORE-Subdocumento4-Anexos, Anexo A (T11-67 a T11-71: especificación mínima); Formulario T-14, paquetes 1.6.2.1 y 1.6.2.2. Falta: costo unitario estimado de cada dispositivo, que va en la Oferta Económica con el mismo ID T11. Evidencia: Especificación de dispositivos de terreno (Formulario T-11) \\
RT-08.11 Obligatorio & La especificación considerará las condiciones reales de uso del caso: intemperie, humedad, polvo, vibración, temperatura, uso con guantes, luminosidad y autonomía de batería requerida por turno. & Cumple & Botonera de Pala y tabletas rugerizadas (IP65 o IP67, operables con guantes, 1.000 nits) & Subdoc. 4, 4.2.1 y 4.2.1.3; OVERFORE-Subdocumento4-Anexos, Anexo A (tableta y botonera); Formulario T-14, paquetes 1.6.2.1 y 1.6.2.2. Evidencia: Prueba en cabina durante un turno completo de invierno y uno de verano (RNF-USA-003) \\
RT-08.12 Obligatorio & Los dispositivos declararán su grado de protección contra polvo y agua y su resistencia a caídas, coherentes con el entorno de operación. & Cumple & Tabletas rugerizadas (resistentes al polvo y a las caídas) & Subdoc. 3, 3.3.3; Formulario T-14, paquetes 1.6.2.1 y 1.6.2.2. Evidencia: Ficha técnica por dispositivo (Formulario T-11) \\
RT-08.13 Obligatorio & El Proponente indicará el ciclo de vida esperado de cada dispositivo, la disponibilidad de repuestos y el plan de reposición durante los 56 meses del Contrato. & Cumple parcialmente & Stock de dispositivos de reemplazo en el sitio & Subdoc. 3, 3.2.2.5; Formulario T-14, paquete 1.8.4.2. Falta: ciclo de vida esperado de cada dispositivo. Evidencia: Plan de reposición de dispositivos \\
RT-08.14 Obligatorio & Los dispositivos se integrarán a la gestión centralizada de flota exigida en RT-03.18. & Cumple & Administración remota y centralizada de todos los dispositivos & Subdoc. 3, 3.3.3; Formulario T-14, paquetes 1.4.5.1 y 1.6.2.1. Evidencia: Prueba de gestión de flota: inventario, actualización y bloqueo (RF-CNX-004) \\
RT-08.15 Deseable & El Proponente proveerá una unidad de cada tipo de dispositivo especificado para pruebas de aceptación por parte del Cliente, antes de la compra masiva. & Cumple & Entrega de una unidad de cada tipo de dispositivo de terreno para aceptación antes de la compra masiva & Subdoc. 4, 4.2.1.3; Subdoc. 4, Anexo C (hardware de terreno); Formulario T-14, paquete 1.6.2.1. Evidencia: Acta de recepción de unidades de prueba \\
RT-08.16 Obligatorio & El Proponente presentará el plan de ciclo de vida del equipamiento: recepción, puesta en servicio, mantención, actualización, retiro y disposición final. & Cumple & Soporte de cinco años en sitio; reposición por falla y reposición programada sólo de baterías, filtros y medios de respaldo; retiro con certificado de borrado seguro y de disposición & Subdoc. 4, 4.2.1.4 y 4.1.1.5; OVERFORE-Subdocumento4-Anexos, Anexo A (campo Ciclo de vida / reposición de cada ficha T11); Formulario T-14, paquete 1.8.4.2. Evidencia: Plan de ciclo de vida del equipamiento \\
RT-08.17 Obligatorio & Todo medio de almacenamiento que salga de servicio será borrado de forma segura y verificable, con certificado de destrucción o de sanitización entregado al Cliente. & Cumple & Borrado seguro de cada medio que sale de servicio, con certificado de sanitización entregado al Cliente & Formulario T-14, paquete 1.8.4.2. Evidencia: Certificado de sanitización de medios entregado al Cliente. \\
RT-08.18 Obligatorio & La disposición final de equipamiento electrónico se realizará con gestor autorizado, conforme a la normativa de residuos aplicable, con certificado de disposición. & Cumple & Disposición final de equipamiento electrónico con gestor autorizado y certificado de disposición para el Cliente & Formulario T-14, paquete 1.8.4.2. Evidencia: Certificado de disposición con gestor autorizado. \\
RT-08.19 Deseable & Se valorará una estrategia de reacondicionamiento o de extensión de vida útil que reduzca el impacto ambiental, cuantificada en la propuesta. & No aplica & Sin componente comprometido en la propuesta (requisito deseable) & Deseable no comprometido en la propuesta. Evidencia: Deseable no comprometido. \\
\end{tabla}
\begin{tabla}{I{0.134}I{0.227}I{0.111}I{0.183}I{0.345}}{Capítulo 09. Desempeño, capacidad y escalabilidad.}{tab:capitulo-09-desempeno-capacidad-escalabilidad}
  {\cab{ID requerimiento} & \cab{Descripción} & \cab{Estado} & \cab{Componente que lo satisface} & \cab{Sección de la propuesta y evidencia}}
  {\emph{Fuente:} elaboración propia a partir de las Bases técnicas transversales.}
RT-09.01 Obligatorio & El Proponente presentará el cálculo de capacidad que sustenta su dimensionamiento, con los supuestos de usuarios concurrentes, transacciones por segundo, volumen de datos y crecimiento anual, tomados de la volumetría del caso. & Cumple & 245 personas concurrentes en régimen (40 TPS) y 505 en el relevo (135 TPS calculadas), 200 TPS en la punta, series de unos 94 millones de lecturas diarias y factor de diseño de tres veces & Subdoc. 4, 4.2.4.1 y 4.2.4.2 (Tabla \refext{S4}{tab:declaracion-uso-ia-seccion}); Subdoc. 4, Anexo D, D.1; Subdoc. 2, 2.3.2 (Tabla \ref{tab:capitulo-02-modelo-arquitectura-referencia}); Formulario T-14, paquete 1.3.1.2. Evidencia: Cálculo de capacidad con supuestos de la volumetría del caso; base de datos en Subdoc. 2, 2.3.2 (Tabla \ref{tab:capitulo-02-modelo-arquitectura-referencia}) \\
RT-09.02 Según caso & La solución soportará la concurrencia y el volumen de transacciones que fije el caso, y mantendrá los umbrales del numeral 9.1 bajo esa carga. & Cumple & Volumetría del caso (Tabla \refext{S4}{tab:volumetria-referencia-capacidad-diseno}: 40 TPS en régimen, 200 TPS en la punta, prueba a 300 TPS) y umbrales de terreno en el p95 (2 s, 3 s y 1,5 s) & Subdoc. 4, 4.2.4.2 (Tabla \refext{S4}{tab:volumetria-referencia-capacidad-diseno}); Subdoc. 9, 9.2.4.5; Subdoc. 2, 2.3.2; Subdoc. 2, Anexo A.2 (RNF-CAP-007); Formulario T-14, paquetes 1.5.3.1 y 1.5.6.1. Evidencia: Prueba de carga con el perfil de volumetría (RNF-CAP-007) \\
RT-09.03 Obligatorio & La solución soportará, sin rediseño, un crecimiento de al menos tres veces la volumetría inicial del caso en un horizonte de tres años. & Cumple & Factor de diseño de tres veces: crecimiento hasta 600 TPS por escalado, sin cambiar la topología & Subdoc. 4, 4.2.4, 4.2.4.3 y 4.2.4.6, Tabla \refext{S4}{tab:declaracion-uso-ia-seccion}; Subdoc. 2, 2.5.1; Subdoc. 2, Anexo A.2 (RNF-CAP-001 a 005); Formulario T-14, paquetes 1.3.1.2 y 1.5.3.1. Evidencia: Prueba de carga a 1,5 veces el peak y demostración del escalamiento a 3 veces la volumetría (RNF-CAP-001 a 005) \\
RT-09.04 Obligatorio & El escalamiento de las capas de aplicación e integración será horizontal y automático, con tiempo de reacción declarado y sin pérdida de transacciones en curso. & Cumple parcialmente & Grupo de servicios de 4 a 9 nodos que escala 30 minutos antes de cada cambio de turno; ampliación a 6 o más nodos según el rendimiento medido & Subdoc. 4, 4.2.2.2, 4.2.4.3 y 4.2.4.6; OVERFORE-Subdocumento4-Anexos, Anexo A (T11-10); Formulario T-14, paquete 1.6.1.1. Falta: tiempo de reacción del escalamiento automático ante carga no prevista. Evidencia: Prueba de carga creciente en Preproducción (RNF-CAP-008) \\
RT-09.05 Obligatorio & El Proponente identificará el componente que primero se convertirá en cuello de botella al crecer la carga y explicará cómo lo detectará y cómo lo resolverá. & Cumple & Disco del grupo de datos (series de tiempo) como primer cuello de botella; ampliación en línea al proyectar 80 \% de ocupación a seis meses & Subdoc. 4, 4.2.4.4 y 4.2.4.6; OVERFORE-Subdocumento4-Anexos, Anexo A (T11-15: ampliar al proyectar 80 \% de ocupación sostenida para los siguientes seis meses); Formulario T-14, paquete 1.3.1.2. Evidencia: Plan de crecimiento y umbrales de ampliación (Subdoc. 4, 4.2.4.6) \\
RT-09.06 Obligatorio & Se ejecutarán pruebas de carga sobre Preproducción con un volumen equivalente a 1,5 veces el peak declarado, y pruebas de estrés hasta identificar el punto de quiebre de la solución. & Cumple & Campaña de carga sostenida a 300 TPS (1,5 veces el peak de 200) y estrés por escalones hasta el punto de quiebre, en Preproducción & Subdoc. 9, 9.2.4.5; Subdoc. 4, 4.1.3.4 y 4.2.3.1; Formulario T-14, paquetes 1.5.3.1, 1.5.6.1 y 1.8.8.3. Evidencia: Informe de pruebas de carga a 1,5 veces el peak y de estrés hasta el punto de quiebre \\
RT-09.07 Obligatorio & El informe de pruebas de carga incluirá la curva de tiempo de respuesta frente a carga, el punto de saturación, el consumo de recursos y el comportamiento durante y después del peak. & Cumple & Informe de pruebas de carga por etapa con curva de tiempo de respuesta, punto de saturación, consumo de recursos y comportamiento durante y después del peak & Subdoc. 9, 9.2.4.5; Formulario T-14, paquetes 1.5.3.1 y 1.5.6.1. Evidencia: Informe de pruebas de carga de cada etapa con la curva de tiempo de respuesta, el punto de saturación, el consumo de recursos y el comportamiento durante y después del peak. \\
RT-09.08 Obligatorio & La solución degradará de forma controlada al superarse la capacidad: encolamiento, limitación de tasa y mensaje explícito a la persona usuaria, nunca error genérico ni pérdida silenciosa de transacciones. & Cumple parcialmente & Nivelación de carga con colas (la saturación se encola), límite de 50 llamadas por segundo por integración y rechazo con error explícito de lo que no valida & Subdoc. 4, 4.1.2.4, 4.1.4.3 y 4.2.2.2; Subdoc. 9, 9.2.4.5; Formulario T-14, paquetes 1.5.3.1 y 1.6.1.1. Falta: mensaje explícito en la App SITCO Terreno al superarse la capacidad, sin error genérico. Evidencia: Inyección de fallas y prueba de estrés (RNF-DIS-007) \\
RT-09.09 Obligatorio & Se gestionará la capacidad durante la Operación con proyección trimestral de crecimiento, alertas anticipadas de agotamiento y propuesta de ajuste de dimensionamiento y de costo. & Cumple parcialmente & Centro de operaciones 24x7 (alertas antes de llegar al límite de un recurso); umbrales de ampliación con proyección a seis meses & Subdoc. 3, 3.2.2.6; Subdoc. 4, 4.2.4.6; Formulario T-14, paquete 1.8.2.3. Falta: propuesta de ajuste de costo derivada de la proyección de capacidad. Evidencia: Informe trimestral de capacidad \\
RT-09.10 Deseable & Se valorará la existencia de pruebas de carga automatizadas ejecutadas de forma periódica en el flujo de integración continua, con detección de regresiones de desempeño. & Cumple & Prueba de carga corta automatizada en la puerta 2 del flujo de integración continua, con límite de degradación de 10 \% en el percentil 95 & Subdoc. 9, 9.2.1 (puerta 2) y 9.2.4.4 (automatización); Formulario T-14, paquetes 1.5.1.2 y 1.5.3.1. Evidencia: Prueba de carga corta automatizada en la puerta 2, con límite de degradación de 10 \% en el percentil 95 para detectar regresiones de desempeño. \\
\end{tabla}
\begin{tabla}{I{0.132}I{0.269}I{0.115}I{0.175}I{0.309}}{Capítulo 10. Disponibilidad, continuidad y resiliencia.}{tab:capitulo-10-disponibilidad-continuidad-resiliencia}
  {\cab{ID requerimiento} & \cab{Descripción} & \cab{Estado} & \cab{Componente que lo satisface} & \cab{Sección de la propuesta y evidencia}}
  {\emph{Fuente:} elaboración propia a partir de las Bases técnicas transversales.}
RT-10.01 Obligatorio & La solución alcanzará una disponibilidad mensual mínima de 99,9 \% para los servicios clasificados como críticos, medida sobre la transacción de negocio de extremo a extremo. & Cumple & Servicios críticos con disponibilidad del 99,9 \% mensual & Subdoc. 3, 3.2.2.6; Formulario T-14, paquetes 1.3.1.2, 1.8.2.1 y 1.8.7.1. Evidencia: Informe mensual de nivel de servicio medido sobre la transacción de extremo a extremo (RNF-DIS-002) \\
RT-10.02 Obligatorio & El Proponente clasificará cada servicio de la solución en crítico, alto, medio o bajo, justificando la clasificación con el impacto operacional de su indisponibilidad, y aplicará a cada uno el nivel de servicio correspondiente del Artículo 78° de las Bases administrativas. & Cumple & Clasificación de servicios del Art. 78.2: crítico para las funciones de terreno (99,9 \%) y alto para gestión, analítica y portales (99,5 \%); los cuatro niveles de la sección 3.2.2.6 & Subdoc. 3, 3.2.2.6; Subdoc. 2, Anexo A.2 (RNF-DIS-002); Formulario T-14, paquetes 1.3.1.2 y 1.8.2.1. Evidencia: Clasificación de cada servicio en crítico, alto, medio o bajo conforme al Art. 78° \\
RT-10.03 Obligatorio & El plan de continuidad del negocio se elaborará conforme a ISO 22301, con análisis de impacto en el negocio, escenarios de contingencia, procedimientos manuales de respaldo y criterios de activación. & Cumple & Plan de continuidad del negocio conforme a ISO 22301, con análisis de impacto, escenarios de contingencia y criterios de activación & Formulario T-14, paquetes 1.3.1.4 y 1.8.6.1. Evidencia: Plan de continuidad del negocio conforme a ISO 22301 (análisis de impacto, escenarios, criterios de activación). \\
RT-10.04 Obligatorio & La continuidad TIC se estructurará conforme a ISO/IEC 27031, articulada con el plan de recuperación ante desastres del Capítulo 7. & Cumple & Plan de continuidad TIC conforme a ISO/IEC 27031, con recuperación ante desastres (RTO de 4 h y RPO de 15 min) & Subdoc. 4, 4.3.2 y 4.3.2.2 (recuperación ante desastres con RTO de 4 h y RPO de 15 min); Formulario T-14, paquetes 1.3.1.4 y 1.8.6.2. Evidencia: Plan de continuidad TIC conforme a ISO/IEC 27031 \\
RT-10.05 Obligatorio & Los mantenimientos programados se ejecutarán fuera de la ventana operacional crítica que defina el caso, con aviso previo mínimo de diez días hábiles. & Cumple & Ventana operacional protegida: intervenciones menores de 03:00 a 05:00 y mayores en las detenciones de marzo y septiembre, nunca en el cierre contable ni en las 48 h previas a un embarque & Subdoc. 3, 3.2.3.3; Subdoc. 4, 4.1.3.4, 4.2.3.1 y 4.3.2.2; Subdoc. 6, 6.2.1.4; Formulario T-14, paquete 1.8.4.1. Evidencia: Calendario de intervenciones y registro de avisos con 10 días hábiles \\
RT-10.06 Obligatorio & La solución permitirá desplegar cambios sin interrupción del servicio. Las ventanas de indisponibilidad programada serán excepcionales y deberán justificarse caso a caso. & Cumple & Estrategia de despliegue sin interrupción: azul-verde en la nube, extracción en el Nodo de Faena con versión anterior lista y despliegue por anillos en terreno & Subdoc. 3, 3.2.3.3; Subdoc. 4, 4.1.3.4; Subdoc. 4, 4.2.3.1 (despliegues progresivos); OVERFORE-Subdocumento4-Anexos, Anexo E (8.32); Formulario T-14, paquete 1.4.1.2. Evidencia: Despliegue sin interrupción demostrado en Preproducción (RNF-DIS-008) \\
RT-10.07 Obligatorio & Se ejecutarán pruebas de resiliencia mediante inyección controlada de fallas —caída de instancia, de zona, de dependencia externa, latencia elevada, saturación de disco— antes de cada paso a producción y al menos una vez por semestre durante la Operación. & Cumple & Pruebas de resiliencia por inyección de fallas (instancia, zona, dependencia, latencia y disco) antes de cada paso a producción y semestralmente en la operación & Subdoc. 9, 9.2.4.5, 9.2.4.3 y 9.3.3; Subdoc. 6, 6.2.2.3; Subdoc. 3, 3.2.2.6; Formulario T-14, paquetes 1.5.3.2, 1.5.6.2 y 1.8.8.1. Evidencia: Informe de pruebas de resiliencia por escenario (caída de instancia, zona, dependencia, latencia, disco) \\
RT-10.08 Obligatorio & El Proponente documentará, por cada dependencia externa, el comportamiento de la solución cuando esa dependencia no responde, responde con error o responde con lentitud. & Cumple parcialmente & Tablas \refext{S4}{tab:campos-obligatorios-evento-dominio} y \refext{S4}{tab:llamadas-sincronas-entre-componentes-umbral}: comportamiento ante falla de las 14 integraciones; reglas comunes de tiempo de espera, reintento, cortacircuito y mamparo para error y lentitud & Subdoc. 4, 4.1.7.1 (Tablas \refext{S4}{tab:campos-obligatorios-evento-dominio} y \refext{S4}{tab:llamadas-sincronas-entre-componentes-umbral}; reglas de resiliencia comunes); Subdoc. 4, 4.1.7.2 (Tabla \refext{S4}{tab:campos-obligatorios-evento-dominio}); Subdoc. 3, Figura \refext{S3}{fig:funcionamiento-sitco-aranda-ante-corte-enlace}; Formulario T-14, paquetes 1.4.2.1, 1.3.1.4 y 1.5.3.2. Falta: comportamiento ante error y lentitud de las integraciones de la Etapa 2 (1.4.2.3). Evidencia: Catálogo de dependencias externas y comportamiento (Subdoc. 4) \\
RT-10.09 Deseable & Se declarará un presupuesto de error por servicio crítico y su vinculación con el ritmo de despliegue de cambios. & Por cumplir & Informe mensual de niveles de servicio con porcentaje de cambios fallidos de los despliegues del mes & Falta: presupuesto de error por servicio crítico y su vinculación con el ritmo de despliegue de cambios. Evidencia: Informe mensual de niveles de servicio con el porcentaje de cambios fallidos de los despliegues del mes (Formulario T-14, paquete 1.8.7.1). \\
\end{tabla}
\begin{tabla}{I{0.132}I{0.241}I{0.109}I{0.214}I{0.305}}{Capítulo 11. Seguridad de la información.}{tab:capitulo-11-seguridad-informacion}
  {\cab{ID requerimiento} & \cab{Descripción} & \cab{Estado} & \cab{Componente que lo satisface} & \cab{Sección de la propuesta y evidencia}}
  {\emph{Fuente:} elaboración propia a partir de las Bases técnicas transversales.}
RT-11.01 Obligatorio & La arquitectura de seguridad se basará en el modelo Zero Trust conforme a NIST SP 800-207: verificación explícita de cada solicitud, privilegio mínimo y presunción de compromiso. & Cumple & Modelo de confianza cero (verificación explícita, privilegio mínimo y presunción de compromiso), con una excepción acotada del Nodo de Faena durante un corte & Subdoc. 4, 4.1.8.1 y 4.1.4.2; Formulario T-14, paquetes 1.3.3.1 y 1.3.3.3. Evidencia: Documento de arquitectura de seguridad (Subdoc. 4) \\
RT-11.02 Obligatorio & El Proponente entregará un modelado de amenazas documentado por cada componente y por cada integración externa, con metodología declarada (STRIDE u otra), y lo actualizará ante cada cambio arquitectónico relevante. & Cumple & Modelado de amenazas con STRIDE sobre canales, módulos, plataformas, servicios técnicos y las catorce integraciones, versionado y revisado ante cada cambio & Subdoc. 4, 4.1.8.5 (Tabla \refext{S4}{tab:amenazas-mayor-consecuencia-componentes-expuestos-control}); Subdoc. 6, 6.2.2.2; Formulario T-14, paquete 1.3.3.2. Evidencia: Modelado de amenazas por componente e integración (STRIDE u otra) \\
RT-11.03 Obligatorio & La información del Cliente se clasificará por nivel de sensibilidad, con controles diferenciados por nivel documentados en una matriz. & Cumple & Módulo Seguridad y Salud Ocupacional (datos de salud con cifrado y acceso restringido) & Subdoc. 4, 4.1.8.4 (Tabla \refext{S4}{tab:mecanismo-cifrado-custodia-claves-ambito}); OVERFORE-Subdocumento4-Anexos, Anexo E (5.12 y 5.13); Subdoc. 3, 3.4.2.6; Formulario T-14, paquete 1.3.2.2. Evidencia: Matriz de clasificación y controles por nivel \\
RT-11.04 Obligatorio & Existirá un programa de gestión de vulnerabilidades con escaneo continuo y plazos máximos de remediación de 7 días corridos para vulnerabilidades críticas, 15 días para altas y 30 días para medias, contados desde su publicación o detección. & Cumple & Plazos de corrección de 7, 15 y 30 días; análisis de composición en cada integración; informe mensual de nivel de servicio & Subdoc. 6, 6.2.2.3; Subdoc. 9, 9.2.3; Subdoc. 3, 3.2.2.6; Formulario T-14, paquetes 1.5.4.3 y 1.8.2.2. Evidencia: Informe mensual del escáner con fecha de detección y cierre (RNF-TRV-11.04) \\
RT-11.05 Obligatorio & El Proponente mantendrá una matriz de controles de seguridad trazable a ISO/IEC 27001 e ISO/IEC 27002, indicando el control, su implementación concreta en la solución y la evidencia que lo acredita. & Cumple & Matriz de controles trazable a ISO/IEC 27001:2022 y 27002:2022 (Anexo E), con implementación, ficha T-11 y evidencia por control & Subdoc. 4, 4.1.8.5 (Tabla \refext{S4}{tab:amenazas-mayor-consecuencia-componentes-expuestos-control}); Subdoc. 4, Anexo E, E.1 a E.4; certificación de la empresa en Subdoc. 1, sección 1.4; Formulario T-14, paquete 1.3.3.2. Evidencia: Matriz de controles de seguridad (Subdoc. 4, Anexo E) \\
RT-11.06 Obligatorio & Se aplicarán los controles de ISO/IEC 27017 para los servicios en nube y de ISO/IEC 27018 para el tratamiento de datos personales en nube. & Cumple & Matriz de controles de ISO/IEC 27017 (siete controles de nube) y de ISO/IEC 27018 (diez principios de privacidad), con certificación institucional declarada & Subdoc. 4, 4.1.8.5; Subdoc. 4, Anexo E, E.5 y E.6; Subdoc. 1, sección 1.4 (Tabla \refext{S1}{tab:certificaciones-institucionales}); Formulario T-14, paquete 1.3.3.2 Evidencia: Matriz de controles ISO 27017 y 27018 (Subdoc. 4) \\
RT-11.07 Obligatorio & La publicación de servicios se realizará exclusivamente a través de la capa de borde, con red de distribución de contenidos, cortafuegos de aplicaciones web con reglas gestionadas y personalizadas, y protección contra denegación de servicio distribuida en capas 3, 4 y 7. & Cumple parcialmente & Capa de borde con Azure Front Door Premium: red de distribución de contenidos, WAF en prevención con reglas gestionadas y propias, DDoS en capas 3, 4 y 7 y TLS 1.3 & Subdoc. 4, 4.1.8.2 (Tabla \refext{S4}{tab:mecanismo-identidad-cada-poblacion-usuaria}) y 4.1.4.2; Subdoc. 4, Anexo A (T11-03); Formulario T-14, paquete 1.6.1.1. Evidencia: Revisión de la superficie publicada por la capa de borde (RNF-TRV-11.07) \\
RT-11.08 Obligatorio & El cifrado en tránsito empleará TLS 1.3, con prohibición expresa de TLS 1.0 y 1.1, conjuntos de cifrado modernos, HSTS con precarga y gestión automatizada de certificados con rotación y alerta anticipada de vencimiento. & Cumple & TLS 1.3 con TLS 1.0 y 1.1 rechazados, HSTS con precarga en la capa de borde, TLS mutuo entre servicios y certificados renovados en forma automática & Subdoc. 4, 4.1.8.4 (Tabla \refext{S4}{tab:amenazas-mayor-consecuencia-componentes-expuestos-control}); Subdoc. 4, Anexo E, 8.24; Formulario T-14, paquetes 1.3.3.1 y 1.6.1.1. \\
RT-11.09 Obligatorio & La totalidad de los datos en reposo estará cifrada, con claves gestionadas en un servicio de gestión de claves o en un módulo de seguridad de hardware, política de rotación declarada y separación de funciones en la custodia de claves. & Cumple & Cifrado del 100 \% de los almacenes en nube, faena y respaldos con claves en HSM (Key Vault Premium por región); dispositivos de terreno con registros cifrados y claves en TPM & Subdoc. 4, 4.1.8.4 (Tabla \refext{S4}{tab:mecanismo-cifrado-custodia-claves-ambito}) y 4.2.2.2; Subdoc. 4, Anexo A (T11-21); Subdoc. 3, 3.4.2.5; Formulario T-14, paquete 1.3.3.1. Evidencia: Revisión de configuración de almacenamiento y política de claves (RNF-SEG-009) \\
RT-11.10 Según caso & Los datos de categoría sensible que el caso identifique se cifrarán adicionalmente a nivel de campo, de modo que el acceso a la base de datos no revele su contenido. & Cumple & Cifrado de campo con clave separada por categoría para salud, exámenes preocupacionales, fatiga, alcohol y drogas, biometría, documento de identidad e identidad del operador & Subdoc. 4, 4.1.8.4 (Tabla \refext{S4}{tab:amenazas-mayor-consecuencia-componentes-expuestos-control}); Subdoc. 3, 3.4.2.6; Formulario T-14, paquetes 1.4.3.5 y 1.4.4.2. Evidencia: Prueba de lectura directa de la base de datos (RNF-SEG-001) \\
RT-11.11 Obligatorio & La puerta de enlace de servicios aplicará autenticación, autorización, cuotas, límites de tasa, validación de esquema e inspección de carga útil. & Cumple parcialmente & Dos puertas de enlace (exposición en la nube con API Management Premium y operacional en el Nodo) con autenticación, autorización, cuotas, límites de tasa, validación de esquema e inspección de carga útil & Subdoc. 4, 4.1.4.3, 4.1.8.2 (Tabla \refext{S4}{tab:servicios-alcanzables-desde-fuera-red-compania}) y 4.2.2.2; Formulario T-14, paquetes 1.4.2.2 y 1.6.1.1. Evidencia: Revisión de la puerta de enlace de servicios (RNF-TRV-11.11) \\
RT-11.12 Obligatorio & Los puntos de entrada públicos contarán con protección contra bots y abuso automatizado, con reto progresivo que no degrade la accesibilidad ni bloquee a personas usuarias legítimas. & Cumple & Protección contra bots en Azure Front Door Premium con reto progresivo compatible con WCAG 2.2 AA en los puntos de entrada públicos & Subdoc. 4, 4.1.4.2 y 4.2.2.2; Subdoc. 4, Anexo A (T11-03); Formulario T-14, paquete 1.6.1.1. Evidencia: Prueba de abuso automatizado en portales (RNF-SEG-010) \\
RT-11.13 Obligatorio & El Proponente declarará la superficie de exposición completa de la solución: cada nombre de dominio, puerto y servicio alcanzable desde fuera de la red del Cliente. & Cumple & Tabla \refext{S4}{tab:servicios-alcanzables-desde-fuera-red-compania}: cuatro servicios alcanzables desde fuera (Portal de Contratistas, Portal para Compradores, Portal Web SITCO y puerta de exposición), con acceso, autenticación y control de borde & Subdoc. 4, 4.1.8.2 (Tabla \refext{S4}{tab:mecanismo-identidad-cada-poblacion-usuaria}); Formulario T-14, paquete 1.3.3.1. Evidencia: Declaración de dominios, puertos y servicios expuestos \\
RT-11.14 Obligatorio & Los eventos de seguridad se registrarán de forma centralizada e inalterable, con retención mínima de doce meses en línea y veinticuatro meses adicionales en archivo recuperable. & Cumple & Registro centralizado e inalterable de eventos de seguridad (Log Analytics y Microsoft Sentinel en Chile Central) con 12 meses en línea y 24 en archivo; bitácora inalterable de la sala técnica & Subdoc. 4, 4.1.8.5 y 4.2.2.2; Subdoc. 4, Anexo A (T11-24); Subdoc. 3, 3.4.2.5; Subdoc. 2, Anexo A.1 (RF-ACR-018); Formulario T-14, paquetes 1.6.1.1 y 1.8.2.2. Evidencia: Revisión de retención de 12 meses en línea y 24 en archivo (RNF-SEG-011) \\
RT-11.15 Obligatorio & Los eventos se correlacionarán en una plataforma SIEM, con casos de uso de detección definidos específicamente para el proceso de negocio del caso, y no sólo genéricos de infraestructura. & Cumple & Microsoft Sentinel como SIEM con seis casos de detección propios de la operación de faena (credencial en dos sitios, exportación masiva, consulta de salud fuera de rol, cambio sin segunda aprobación, sello de tiempo, escritura al historiador) & Subdoc. 4, 4.1.8.5 (casos 1 a 6); Subdoc. 4, Tabla \refext{S4}{tab:tecnologias-capas-7-8-alternativa-evaluada} (Microsoft Sentinel); OVERFORE-Subdocumento4-Anexos, Anexo E (5.24 y 5.25); Formulario T-11 (T11-24); Formulario T-14, paquete 1.8.2.2. Evidencia: Casos de uso de detección para el proceso del caso \\
RT-11.16 Obligatorio & Se implementará detección y respuesta en puntos finales y en cargas de trabajo, tanto en nube como on-premise. & Cumple & EDR en las estaciones de administración (T11-48); tabletas y consolas gestionadas con Intune y servidores de faena endurecidos con CIS & OVERFORE-Subdocumento4-Anexos, Anexo E (8.7); OVERFORE-Subdocumento4-Anexos, Anexo A (T11-48); Formulario T-14, paquete 1.8.2.2. Evidencia: Cobertura de detección y respuesta en puntos finales y cargas de trabajo, nube y on-premise \\
RT-11.17 Obligatorio & El Proponente dispondrá de un centro de operaciones de seguridad con cobertura 24×7, propio o subcontratado, y declarará su ubicación, dotación y procedimientos. & Cumple & SOC propio 24x7 en el centro de operaciones de Santiago con cinco analistas de seguridad & Subdoc. 1, 1.2; Subdoc. 1, Anexo F del Subdocumento 1 (Tabla \refext{A1}{tab:dotacion-apoyo-overfore-asignada-proyecto} de los Anexos del Subdocumento 1); Formulario T-14, paquete 1.8.2.2. Evidencia: Ubicación, dotación y procedimientos del centro de operaciones de seguridad \\
RT-11.18 Obligatorio & El plan de respuesta a incidentes de seguridad definirá clasificación, cadena de escalamiento, plazos, responsables y protocolo de comunicación al Cliente dentro de las dos horas de detectado un incidente de severidad crítica. & Cumple & Contención en arquitectura (bloqueo de dispositivo en 15 min, revocación en 24 h, cuentas de emergencia), alertas al responsable de turno y notificación a la autoridad (Ley 21.663) & OVERFORE-Subdocumento4-Anexos, Anexo E (5.5, 5.24 a 5.26 y 6.8); Subdoc. 4, 4.1.8.1; Formulario T-14, paquetes 1.3.3.4 y 1.5.4.2. Evidencia: Plan de respuesta a incidentes de seguridad \\
RT-11.19 Obligatorio & Toda brecha de seguridad o de datos personales se notificará al Cliente dentro de las 24 horas de su detección, con informe preliminar, y el análisis de causa raíz se entregará dentro de los cinco días hábiles siguientes. & Cumple & Procedimiento de notificación de brechas al Cliente dentro de 24 horas, con informe preliminar y análisis de causa raíz en cinco días hábiles & Formulario T-14, paquetes 1.3.3.4 y 1.8.3.2. Evidencia: Procedimiento de notificación de brechas en 24 horas. \\
RT-11.20 Obligatorio & Se ejecutarán pruebas de intrusión por un tercero independiente del Adjudicatario, anualmente y antes de cada paso a producción, con entrega íntegra del informe al Cliente y plan de remediación con plazos. & Cumple & Pruebas de intrusión de un tercero independiente antes de cada paso a producción y una vez al año, con informe íntegro al cliente y plan de remediación con plazos & Subdoc. 9, 9.2.3 y 9.3.3; Subdoc. 6, 6.2.2.3; Formulario T-14, paquetes 1.5.4.1 y 1.8.8.2. Evidencia: Informe de pruebas de intrusión de un tercero independiente \\
RT-11.21 Deseable & Se realizarán ejercicios de simulación de incidente con participación del Cliente al menos una vez al año durante la Operación. & Cumple parcialmente & Ejercicio de respuesta a un incidente de seguridad realizado con el Cliente antes de la Operación & Formulario T-14, paquete 1.5.4.2. Falta: ejercicio anual de simulación de incidente con el Cliente durante la Operación. Evidencia: Ejercicio de respuesta a un incidente de seguridad realizado con el Cliente antes de la Operación. \\
RT-11.22 Obligatorio & El flujo de integración continua incorporará análisis estático de código, análisis de composición de software, análisis dinámico y escaneo de imágenes de contenedor, con criterios de bloqueo automático del despliegue ante hallazgos críticos. & Cumple & Integración continua con análisis estático, de composición, escaneo de imágenes y de infraestructura como código, y análisis dinámico en QA; bloqueo ante hallazgos críticos o altos & Subdoc. 6, 6.2.2.2 y 6.2.2.3; Subdoc. 9, 9.2.3; Formulario T-14, paquetes 1.4.1.1 y 1.5.1.3. Evidencia: Configuración del flujo de integración continua con análisis estático, de composición, dinámico y de imágenes \\
RT-11.23 Obligatorio & Cada versión liberada se acompañará de su inventario de componentes de software en formato CycloneDX o SPDX, entregado al Cliente. & Cumple & Lista de materiales de cada versión generada por el análisis de composición y entregada en CycloneDX o SPDX, con procedencia SLSA nivel 3 & Subdoc. 6, 6.2.2.2; Formulario T-14, paquetes 1.4.1.3 y 1.4.1.4. Evidencia: Inventario de componentes CycloneDX o SPDX por versión \\
RT-11.24 Obligatorio & Los artefactos se firmarán y su procedencia se verificará conforme a SLSA nivel 3 o superior. & Cumple & Firma de artefactos desplegados con lista de componentes CycloneDX o SPDX y verificación de procedencia SLSA nivel 3 & Subdoc. 6, 6.2.2.2 (firma de la imagen) y 6.2.2.3 (verificación de la firma); Formulario T-14, paquete 1.4.1.3. Evidencia: Firma de artefactos y verificación de procedencia SLSA 3. \\
RT-11.25 Obligatorio & Queda prohibido el uso de datos productivos reales en ambientes no productivos sin anonimización o seudonimización verificable. & Cumple & Datos sintéticos en Desarrollo y QA; copias con campos sensibles enmascarados de forma irreversible en Preproducción; borrado al cerrar cada campaña & Subdoc. 9, 9.2.4.2; Subdoc. 4, 4.1.3.4 y 4.2.2.2; Formulario T-14, paquete 1.4.6.2. Evidencia: Revisión de datos de ambientes no productivos \\
RT-11.26 Obligatorio & El Proponente declarará el proceso de aprobación de nuevas dependencias de terceros, incluyendo criterios de licencia, mantención activa y ausencia de vulnerabilidades conocidas. & Cumple & Proceso de aprobación de cada dependencia nueva de terceros por licencia, mantención y vulnerabilidades & Subdoc. 6, 6.2.2.2; Formulario T-14, paquete 1.4.1.1. Evidencia: Proceso de aprobación de dependencias de terceros por licencia, mantención y vulnerabilidades. \\
RT-11.27 Obligatorio & Las personas desarrolladoras no tendrán acceso interactivo directo al ambiente de producción. Todo acceso excepcional será temporal, aprobado, registrado y con sesión grabada. & Cumple & Producción con acceso restringido y auditado; promoción sólo con imagen firmada por GitOps; elevación temporal de administración con aprobación previa y sesión grabada & Subdoc. 9, 9.2.4.1; Subdoc. 4, 4.1.8.3 y 4.1.8.5 (Tabla \refext{S4}{tab:modulo-plataforma-sostiene-cada-funcion-durante}); Subdoc. 4, 4.1.3.4 y 4.2.3.1; OVERFORE-Subdocumento4-Anexos, Anexo E (8.2 y 8.18); Formulario T-14, paquetes 1.3.3.3 y 1.4.1.1. Evidencia: Registro de accesos excepcionales a producción con sesión grabada \\
RT-11.28 Deseable & Se aplicará el marco OWASP SAMM o equivalente para medir y mejorar la madurez del proceso de desarrollo seguro, con evaluación inicial y reevaluación anual. & Cumple & Evaluación inicial de madurez OWASP SAMM en las quince prácticas, con reevaluación anual & Subdoc. 9, 9.1.3; Subdoc. 6, 6.2; Formulario T-14, paquetes 1.1.4.4 y 1.8.8.6. Evidencia: Evaluación de madurez OWASP SAMM inicial y reevaluada cada año. \\
\end{tabla}
\begin{tabla}{I{0.134}I{0.233}I{0.110}I{0.194}I{0.329}}{Capítulo 12. Identidad, acceso y sesiones.}{tab:capitulo-12-identidad-acceso-sesiones}
  {\cab{ID requerimiento} & \cab{Descripción} & \cab{Estado} & \cab{Componente que lo satisface} & \cab{Sección de la propuesta y evidencia}}
  {\emph{Fuente:} elaboración propia a partir de las Bases técnicas transversales.}
RT-12.01 Obligatorio & La gestión de identidad será centralizada, con federación mediante OpenID Connect y OAuth 2.1, o SAML 2.0 cuando la integración con el Cliente lo requiera, e integración con el directorio corporativo del Cliente por LDAP o su equivalente en la nube. & Cumple & Identidad federada con el directorio corporativo (Entra ID con doble factor) y plano externo separado (Entra External ID); OAuth 2.1 o TLS mutuo en las integraciones & Subdoc. 4, 4.1.8.3 (Tabla \refext{S4}{tab:mecanismo-identidad-cada-poblacion-usuaria}) y 4.2.2.2; Subdoc. 3, 3.4.3; Formulario T-14, paquete 1.3.3.3. Evidencia: Modelo de identidad (Subdoc. 4) \\
RT-12.02 Obligatorio & Existirá inicio de sesión único para todos los módulos de la solución, con cierre de sesión propagado a todos ellos. & Cumple & Módulo Administración y Servicios Comunes (identidad compartida por todos los módulos); Portal Web SITCO con sesión entregada tras verificación & Subdoc. 3, 3.4.3; Subdoc. 4, 4.1.4.2; Formulario T-14, paquete 1.3.3.3. Evidencia: Prueba de navegación entre módulos y cierre de sesión (RNF-SEG-012) \\
RT-12.03 Obligatorio & La autenticación multifactor será obligatoria para personas usuarias administradoras, para todo acceso privilegiado y para todo acceso originado fuera de la red corporativa. & Cumple & Identidad federada con doble factor; en terreno credencial más PIN & Subdoc. 3, 3.2.3.3 y 3.4.3; Formulario T-14, paquetes 1.3.3.3 y 1.4.5.2. Evidencia: Revisión de accesos administrativos, privilegiados y externos con doble factor (RNF-SEG-003) \\
RT-12.04 Deseable & Se soportarán factores resistentes a la suplantación de identidad, tipo FIDO2 o claves de acceso, al menos para los perfiles administradores. & Por cumplir & Doble factor obligatorio para las personas administradoras, dentro del modelo de identidad y acceso con federación & Subdoc. 4, 4.1.8.3; Formulario T-14, paquete 1.3.3.3. Falta: factor resistente a la suplantación de identidad, tipo FIDO2 o claves de acceso, al menos para los perfiles administradores. Evidencia: Doble factor obligatorio para las personas administradoras. \\
RT-12.05 Obligatorio & El control de acceso será basado en roles, complementado con control basado en atributos donde el proceso lo exija, con matriz de segregación de funciones documentada y verificable. & Cumple parcialmente & Autorización por rol y por atributos de sitio, empresa y área; Módulo Administración y Servicios Comunes (usuarios y roles) & Subdoc. 4, 4.1.8.1; OVERFORE-Subdocumento4-Anexos, Anexo E (5.3); Subdoc. 3, 3.4.3; Subdoc. 2, Anexo A.1 (RF-ACR-016); Formulario T-14, paquete 1.3.3.3. Falta: que la matriz de segregación de funciones sea verificable. Evidencia: Matriz de roles y segregación de funciones (RF-ACR-016) \\
RT-12.06 Obligatorio & Los accesos privilegiados se gestionarán con elevación temporal a demanda, aprobación previa y grabación de sesión para las operaciones de mayor riesgo. & Cumple & Elevación de privilegios de administración temporal, con aprobación previa y sesión grabada; zona de gestión fuera de banda solo con doble factor y sesión grabada & Subdoc. 4, 4.1.8.3, Tabla \refext{S4}{tab:modulo-plataforma-sostiene-cada-funcion-durante} y \refext{S4}{tab:tecnologias-capas-4-5-alternativa-evaluada}.3.3; Subdoc. 4, Anexo E, 8.18; Formulario T-14, paquete 1.3.3.3. Evidencia: Revisión de cuentas privilegiadas y de sesiones grabadas (RNF-SEG-013) \\
RT-12.07 Obligatorio & La política de sesión declarará duración máxima, caducidad por inactividad, renovación de la credencial de sesión tras la autenticación, revocación inmediata y control de sesiones concurrentes. & Cumple & Política de sesión con duración máxima, caducidad por inactividad, renovación tras la autenticación, revocación inmediata y control de sesiones concurrentes & Subdoc. 4, 4.1.8.3; Formulario T-14, paquete 1.3.3.3. Evidencia: Prueba de caducidad y revocación de sesión (RNF-SEG-014) \\
RT-12.08 Obligatorio & Las credenciales de sesión serán firmadas y de vida breve, con credencial de refresco rotatoria. Queda prohibido transportar identificadores de sesión en la ruta de la dirección web. & Cumple & Credenciales de sesión firmadas y de vida breve, con refresco rotatorio y sin identificadores de sesión en la dirección web & Subdoc. 4, 4.1.8.3; Formulario T-14, paquete 1.3.3.3. Evidencia: Prueba de manipulación de credenciales de sesión (RNF-SEG-014) \\
RT-12.09 Obligatorio & Se registrará auditoría completa del ciclo de vida de la identidad: creación, modificación, elevación, bloqueo y baja de cuentas, con no repudio y retención declarada. & Cumple & Auditoría inalterable con autor, fecha, motivo y valores anterior y posterior; registro de cuenta, rol y sesión con evidencia de auditoría de 10 años & Subdoc. 4, 4.1.8.5 (Tabla \refext{S4}{tab:modulo-plataforma-sostiene-cada-funcion-durante}) y 4.1.8.3; Subdoc. 3, 3.4.3; Subdoc. 5, Anexo A.3 (asiento de auditoría, 10 años); OVERFORE-Subdocumento4-Anexos, Anexo E (5.26, 6.5 y 8.2); Formulario T-14, paquete 1.3.3.3. Evidencia: Revisión del registro de ciclo de vida de cuentas (RF-ADM-007) \\
RT-12.10 Obligatorio & El aprovisionamiento y el desaprovisionamiento estarán automatizados y ligados al ciclo de vida laboral, con baja efectiva en un plazo no superior a 24 horas desde la desvinculación. & Cumple parcialmente & Módulo Acreditación y Acceso (habilitación por contrato, plantilla biométrica eliminada al perder el acceso) e identidad federada; revocación propagada en 24 h o menos con enlace & Subdoc. 4, 4.1.8.1; OVERFORE-Subdocumento4-Anexos, Anexo E (6.5); Subdoc. 3, 3.4.2.5; Subdoc. 2, Anexo A.1 (RF-ACR-013); Formulario T-14, paquetes 1.3.3.3 (baja en 24 horas o menos desde la desvinculación) y 1.4.2.2 (copia local de identidades con vigencia de 72 horas). Falta: durante un corte prolongado la baja no se propaga al Nodo; el área de seguridad del Cliente la registra en la consola local y el Nodo rechaza la credencial al vencer su copia de 72 horas, con lo admitido marcado en la reconciliación. El plazo de 24 horas no se garantiza en faena. Evidencia: Prueba de revocación cronometrada sobre 20 casos, incluidos los nodos de borde (RNF-SEG-006) \\
RT-12.11 Según caso & El mecanismo de autenticación se adecuará al perfil operacional real descrito en el caso: entornos de terreno, uso con guantes, baja alfabetización digital, dispositivos compartidos por turno y ausencia de correo electrónico personal. & Cumple & Botonera de Pala (toque de credencial); tabletas (credencial y PIN numérico) & Subdoc. 3, 3.2.3.3, 3.3.3; Formulario T-14, paquetes 1.3.3.3 y 1.4.5.1. Evidencia: Prueba en turno de noche de invierno con guantes puestos (RNF-DES-004, RNF-USA-001) \\
RT-12.12 Según caso & Las personas usuarias externas del Cliente —clientes, proveedores, pacientes, productores, según el caso— dispondrán de un mecanismo de registro, verificación de identidad y recuperación de acceso autoservido y seguro. & Cumple & Portal de Contratistas y Portal para Compradores con identidad externa (Entra External ID): compradores, operador portuario y laboratorios con registro y recuperación autoservidos; contratistas con doble factor y administración delegada & Subdoc. 4, 4.1.8.3 (Tabla \refext{S4}{tab:mecanismo-cifrado-custodia-claves-ambito}); Subdoc. 3, 3.4.2.5; Formulario T-14, paquetes 1.4.4.5 y 1.4.5.2. Evidencia: Prueba de autoatención y recuperación de acceso (RF-ACR-015) \\
RT-12.13 Obligatorio & El Proponente declarará el procedimiento de acceso de emergencia (cuenta de último recurso), su custodia, su control y su auditoría. & Cumple & Cuentas de emergencia custodiadas para el caso de falla del proveedor de identidad & Subdoc. 4, 4.2.3.4; Subdoc. 4, Anexo E, E.1 (5.26); Formulario T-14, paquete 1.3.3.3. Evidencia: Procedimiento de cuenta de último recurso con custodia y auditoría \\
\end{tabla}
\begin{tabla}{I{0.132}I{0.278}I{0.101}I{0.179}I{0.310}}{Capítulo 13. Usabilidad, accesibilidad y experiencia de usuario.}{tab:capitulo-13-usabilidad-accesibilidad-experiencia-usuario}
  {\cab{ID requerimiento} & \cab{Descripción} & \cab{Estado} & \cab{Componente que lo satisface} & \cab{Sección de la propuesta y evidencia}}
  {\emph{Fuente:} elaboración propia a partir de las Bases técnicas transversales.}
RT-13.01 Obligatorio & Todas las interfaces destinadas a personas usuarias cumplirán WCAG 2.2 nivel AA, verificado con herramientas automatizadas y con pruebas manuales, con informe de conformidad entregable. & Cumple & Capa de presentación conforme a WCAG 2.2 AA, verificada en QA con herramienta automatizada y en Preproducción con pruebas manuales & Subdoc. 9, 9.2.4.3; Subdoc. 4, 4.1.4.1; Subdoc. 9, Anexo C, Tabla C.1; Subdoc. 2, Anexo A.2 (RNF-USA-010); Formulario T-14, paquetes 1.5.1.3 y 1.5.2.1. Evidencia: Informe de conformidad WCAG 2.2 AA con herramientas automatizadas y pruebas manuales (RNF-USA-010) \\
RT-13.02 Obligatorio & El diseño será responsivo y funcionará correctamente en escritorio, tableta y teléfono, con puntos de quiebre coherentes y adaptación inteligente del contenido, no simple reducción. & Cumple parcialmente & Diseño responsivo de las interfaces de usuario (sobre el sistema de diseño documentado, paquete 1.3.4.2 del Formulario T-14) & Subdoc. 2, Anexo A.2 (RNF-USA-011); Subdoc. 9, Anexos (verificación de RNF-USA-011); Formulario T-14, paquetes 1.3.4.2 y 1.4.5.2. Evidencia: Revisión de cada pantalla en 1920x1080, 1366x768, 768x1024 (vertical y horizontal) y 375x812 (RNF-USA-011) \\
RT-13.03 Obligatorio & El Proponente ejecutará investigación con personas usuarias reales del Cliente, prototipado y pruebas de usabilidad antes de la construcción definitiva, y documentará los hallazgos y los cambios de diseño que produjeron. & Cumple & Participación de quienes operan en el diseño; investigación con usuarios, prototipos navegables y pruebas de usabilidad con informe de hallazgos y cambios de diseño & Subdoc. 9, 9.2.4.3; Subdoc. 3, 3.4.5; Formulario T-14, paquetes 1.3.4.1 y 1.3.4.3. Evidencia: Informe de investigación con usuarios, prototipos validados con al menos cinco usuarios del perfil e informe de hallazgos y cambios de diseño \\
RT-13.04 Obligatorio & Se comprometerán indicadores de usabilidad medibles: tiempo máximo de la transacción operacional crítica, número máximo de pasos, tasa de error tolerada y tiempo de aprendizaje esperado por perfil. & Cumple & Umbrales de tiempo de la transacción crítica (2 s, 3 s y 1,5 s), 3 interacciones, tasa de error y tiempo de aprendizaje & Subdoc. 9, 9.2.4.3; Formulario T-13; Subdoc. 3, 3.2.4.3; Formulario T-14, paquetes 1.5.2.1 y 1.5.2.2. Evidencia: Medición de usabilidad por perfil: tiempo, pasos, tasa de error (5 \% con prototipo, 2 \% en marcha blanca) y aprendizaje \\
RT-13.05 Obligatorio & Ninguna funcionalidad principal requerirá más de tres interacciones desde la pantalla de inicio del perfil correspondiente. & Cumple & Interfaces de terreno y portales con tres interacciones o menos desde la pantalla de inicio & Subdoc. 9, 9.2.4.3; Subdoc. 4, 4.1.5.1; Subdoc. 2, Anexo A.2 (RNF-USA-012); Formulario T-14, paquetes 1.3.4.2, 1.3.4.3 y 1.4.5.1. Evidencia: Recorrido de cada funcionalidad principal por perfil contando interacciones (RNF-USA-012) \\
RT-13.06 Obligatorio & La solución entregará retroalimentación visual clara ante cada acción, manejará los errores de forma comprensible —indicando qué ocurrió y qué hacer— y evitará mensajes técnicos dirigidos a la persona usuaria final. & Cumple & Confirmación visual de cada acción en App SITCO Terreno y Consola de Portería; mensajes de error con causa y acción sugerida, según RNF-USA-013 (Subdoc. 2, Anexo A.2) & Subdoc. 4, 4.1.5.1 y 4.1.4.1; Subdoc. 2, Anexo A.2 (RNF-USA-013); Formulario T-14, paquete 1.3.4.2. Evidencia: Revisión del catálogo de mensajes y prueba de usabilidad con personas usuarias del Cliente (RNF-USA-013) \\
RT-13.07 Obligatorio & La solución soportará a personas usuarias con baja alfabetización digital: alto contraste, objetivos táctiles de al menos 44 × 44 píxeles, iconografía acompañada de texto, y flujos guiados paso a paso. & Cumple & Botonera de Pala (cuatro botones grandes, sin pantalla ni formularios) y tabletas con alto contraste, objetivos táctiles de 44 x 44 px e íconos con texto & Subdoc. 3, 3.3.3; Subdoc. 4, 4.1.5.1; Subdoc. 2, Anexo A.2 (RNF-USA-014); Formulario T-14, paquete 1.3.4.2. Evidencia: Revisión de interfaces de terreno y prueba con personas operadoras (RNF-USA-014) \\
RT-13.08 Según caso & Cuando el caso lo requiera, las interfaces de terreno operarán con guantes, a la intemperie, con luminosidad variable y sin conexión. & Cumple & Botonera de Pala y tabletas rugerizadas (IP65, legibles al sol, operables con guantes, de -6 a 34 C) y App SITCO Terreno con operación sin conexión & Subdoc. 3, 3.3.3 y 3.3.6; Subdoc. 4, 4.1.5.1 (Tabla \refext{S4}{tab:interfaces-sincronas-eventos-publicados-datos-consumidos}) y 4.2.1.3; Formulario T-14, paquetes 1.4.5.1, 1.5.2.1, 1.5.2.2 y 1.6.2.1. Evidencia: Prueba en cabina de pala y de camión durante un turno completo de invierno y uno de verano (RNF-USA-001, RNF-USA-002, RNF-USA-003) \\
RT-13.09 Obligatorio & Existirá un sistema de diseño documentado con paleta de colores acotada, jerarquía tipográfica de a lo más dos familias, iconografía coherente, retícula y componentes reutilizables. & Cumple parcialmente & Sistema de diseño documentado (biblioteca de componentes) a cargo del Diseñador UX/UI & Formulario T-14, paquete 1.3.4.2. Evidencia: Sistema de diseño documentado y biblioteca publicada y usada por los prototipos \\
RT-13.10 Obligatorio & El Proponente declarará la matriz de navegadores y de versiones soportadas y su política de actualización. & Cumple & Matriz de navegadores y versiones soportadas con su política de actualización, dentro del prototipo UX/UI & Subdoc. 4, 4.1.1 (tecnologías) y prototipo UX/UI; Formulario T-14, paquete 1.3.4.2. Evidencia: Matriz de navegadores y versiones soportadas y su política de actualización. \\
RT-13.11 Obligatorio & La solución será navegable íntegramente por teclado, con orden de foco lógico y atajos para las operaciones frecuentes. & Cumple & Interfaces de escritorio navegables por teclado, verificadas con pruebas manuales de accesibilidad & Subdoc. 9, 9.2.4.3; Subdoc. 2, Anexo A.2 (RNF-USA-015); Formulario T-14, paquetes 1.3.4.2 y 1.5.1.3. Evidencia: Prueba manual de navegación por teclado en cada pantalla de escritorio (RNF-USA-015) \\
RT-13.12 Deseable & Se valorará el soporte de modo oscuro, de personalización de la interfaz por persona usuaria y de múltiples idiomas cuando el caso lo justifique. & Por cumplir & Prototipo interactivo UX/UI (Cap. 25) sobre el sistema de diseño documentado y su biblioteca de componentes reutilizables & Prototipo interactivo UX/UI (Cap. 25) y Subdoc. 9, 9.2 (verificación); Formulario T-14, paquetes 1.3.4.2 y 1.3.4.3. Falta: modo oscuro, personalización de la interfaz por persona usuaria y soporte de múltiples idiomas. Evidencia: Revisión del prototipo según Cap. 25. \\
\end{tabla}
\begin{tabla}{I{0.132}I{0.273}I{0.115}I{0.178}I{0.303}}{Capítulo 14. Observabilidad y gestión del servicio.}{tab:capitulo-14-observabilidad-gestion-servicio}
  {\cab{ID requerimiento} & \cab{Descripción} & \cab{Estado} & \cab{Componente que lo satisface} & \cab{Sección de la propuesta y evidencia}}
  {\emph{Fuente:} elaboración propia a partir de las Bases técnicas transversales.}
RT-14.01 Obligatorio & La observabilidad será unificada para nube y on-premise, con métricas, registros y trazas distribuidas correlacionadas por un identificador único de transacción, instrumentadas conforme a OpenTelemetry. & Cumple & Capa 8 de observabilidad transversal (OpenTelemetry, Azure Monitor, consola operativa única) e Integrador de Sistemas (identificador común de mensaje) & Subdoc. 4, 4.1.4.8; Subdoc. 3, 3.4.3; Subdoc. 2, Anexo A.2 (RNF-OPE-006); Formulario T-14, paquete 1.6.1.1. Evidencia: Consola única y traza de una operación de punta a punta (RNF-OPE-001, RNF-OPE-006) \\
RT-14.02 Obligatorio & El Cliente dispondrá de acceso propio y permanente a los tableros operacionales y de negocio, con datos en tiempo real y capacidad de exportación. & Cumple parcialmente & Módulo Indicadores y Centro Integrado (tableros de turno y de gestión) y vías de explotación del modelo analítico & Subdoc. 5, 5.2.2 (Tabla \refext{S5}{tab:vias-explotacion-informacion}); Subdoc. 3, 3.4.3; Subdoc. 2, Anexo A.1 (RF-ANA-001); Formulario T-14, paquetes 1.4.3.6 y 1.8.2.1. Falta: datos en tiempo real en los tableros de turno y de gestión; el catálogo define tiempo real en 60 segundos y la latencia comprometida es de 15 minutos en turno y 4 horas en gestión. Evidencia: Acceso del Cliente a los tableros con exportación y medición de la latencia (RF-ANA-001, RNF-DES-010) \\
RT-14.03 Obligatorio & Los indicadores de nivel de servicio se medirán sobre la experiencia real de la persona usuaria y no sobre pruebas sintéticas, sin perjuicio de que estas se empleen como complemento. & Cumple & Informe mensual de nivel de servicio con indicadores medidos sobre la experiencia real de la persona usuaria & Subdoc. 3, 3.2.2.6; Formulario T-14, paquete 1.8.7.1; Subdoc. 2, Anexo A.2 (RNF-OPE-006). Evidencia: Informe mensual de nivel de servicio, con revisión de la definición y del origen de datos de cada indicador (RNF-OPE-006) \\
RT-14.04 Obligatorio & El alertamiento se basará en síntomas de negocio y no sólo en umbrales de infraestructura, con supresión de ruido, agrupación, escalamiento automático y turnos de disponibilidad declarados. & Cumple & Alertas por síntoma con anticipación y escalamiento automático al responsable de turno (4.1.4.8 y 4.1.6.3) & Subdoc. 4, 4.1.4.8; Subdoc. 4, 4.1.6.1.1 (alerta a las 72 h); Subdoc. 4, 4.1.6.3 (escalamiento de fatiga); Subdoc. 3, 3.2.2.6; Subdoc. 2, Anexo A.2 (RNF-OPE-003); Formulario T-14, paquete 1.8.2.1. Evidencia: Catálogo de alertas por síntomas de negocio \\
RT-14.05 Obligatorio & Existirá un libro de operación y una guía de resolución documentados para cada escenario de falla previsible, con automatización progresiva de las tareas repetitivas. & Cumple & Libros de operación y guías de resolución por escenario de falla, con matriz de escalamiento de incidentes & Formulario T-14, paquete 1.1.10.2. Evidencia: Libros de operación y guías de resolución por escenario de falla, con matriz de escalamiento (Art. 91). \\
RT-14.06 Obligatorio & Todo incidente crítico dará lugar a un análisis de causa raíz obligatorio, con informe entregable dentro de cinco días hábiles y seguimiento de las acciones correctivas hasta su cierre. & Cumple & Informe de causa raíz de incidentes críticos en 5 días hábiles & Subdoc. 3, 3.2.2.6; Formulario T-14, paquete 1.8.3.2. Evidencia: Informe de causa raíz y plan de acción con plazo y responsable \\
RT-14.07 Obligatorio & Los registros de la solución no contendrán datos personales sensibles ni credenciales, y su acceso estará controlado y auditado. & Cumple parcialmente & Capa 8 de observabilidad que excluye datos sensibles y credenciales de los registros técnicos (retención de 90 días) & Subdoc. 4, 4.1.4.8; Subdoc. 5, 5.2.4 (Tabla \refext{S5}{tab:retencion-categoria-entidad}) y Anexos A.4 (Tabla A.15); Subdoc. 2, Anexo A.2 (RNF-SEG-015); Formulario T-14, paquete 1.6.1.1. Falta: control de acceso y auditoría de los registros técnicos del Nodo de Faena y de los gabinetes. Evidencia: Revisión automatizada de muestras de registros técnicos (RNF-SEG-015) \\
RT-14.08 Obligatorio & El Proponente declarará la retención de métricas, registros y trazas, y su costo asociado, distinguiendo el almacenamiento en línea del archivado. & Cumple parcialmente & Observabilidad con métricas, registros y trazas correlacionadas según OpenTelemetry, y retención en línea y en archivo declarada & Formulario T-14, paquete 1.6.1.1. Falta: costo asociado a la retención de métricas, registros y trazas. Evidencia: Retención en línea y en archivo declarada; registros de seguridad con 12 meses en línea y 24 en archivo (RNF-SEG-011). \\
RT-14.09 Deseable & Se valorará la detección proactiva de anomalías mediante análisis del comportamiento histórico, con alerta antes de que el incidente afecte a la operación. & Por cumplir & Centro de operaciones de red 24x7 con monitoreo continuo, alertas por síntomas de negocio y supresión de ruido & Formulario T-14, paquete 1.8.2.1. Falta: detección proactiva de anomalías mediante análisis del comportamiento histórico, con alerta antes de que el incidente afecte a la operación. Evidencia: Monitoreo continuo con alertas por síntomas de negocio, supresión de ruido y escalamiento automático (1.8.2.1). \\
\end{tabla}
\begin{tabla}{I{0.136}I{0.202}I{0.111}I{0.198}I{0.353}}{Capítulo 15. Sostenibilidad, eficiencia y certificaciones.}{tab:capitulo-15-sostenibilidad-eficiencia-certificaciones}
  {\cab{ID requerimiento} & \cab{Descripción} & \cab{Estado} & \cab{Componente que lo satisface} & \cab{Sección de la propuesta y evidencia}}
  {\emph{Fuente:} elaboración propia a partir de las Bases técnicas transversales.}
RT-15.01 Obligatorio & El Proponente dimensionará la infraestructura ajustada a la demanda real, evitando capacidad ociosa permanente, y declarará el factor de utilización proyectado. & Cumple & Dimensionamiento desde la volumetría (Tabla \refext{S4}{tab:declaracion-uso-ia-seccion}) con escalado de servicios de 4 a 9 nodos y clúster mínimo de recuperación que escala; utilización declarada por componente (memoria del Nodo de 37 \% a 56 \% con tres veces la carga) & Subdoc. 4, 4.2.4 (Tabla \refext{S4}{tab:volumetria-referencia-capacidad-diseno}; 4.2.4.3 y 4.2.4.4); Subdoc. 4, 4.2.2.2 (escalado de AKS); Subdoc. 4, 4.2.4.6; OVERFORE-Subdocumento4-Anexos, Anexo D (D.3.3); Formulario T-14, paquete 1.3.1.2. Evidencia: Factor de utilización proyectado \\
RT-15.02 Obligatorio & Los ambientes no productivos se apagarán o reducirán fuera del horario de uso. & Cumple & Ambientes no productivos con instancias de menor nivel que se apagan fuera del horario de uso & Subdoc. 4, 4.2.3.1 y 4.2.2.2; Formulario T-13; Formulario T-14, paquete 1.6.1.1. Evidencia: Configuración de apagado de ambientes no productivos \\
RT-15.03 Obligatorio & El Proponente estimará la huella de carbono anual de la operación de la solución y declarará la metodología empleada. & Cumple & Informe anual de consumo y huella de carbono de la solución, con método de cálculo declarado según el Protocolo de Gases de Efecto Invernadero & Formulario T-14, paquete 1.8.7.3. Evidencia: Estimación de huella de carbono anual de la operación de la solución y metodología. \\
RT-15.04 Obligatorio & Se declarará la eficiencia en el uso de la energía (PUE) del recinto on-premise y la intensidad de carbono de la región de nube escogida. & Cumple & Sala técnica del Nodo de Faena con pasillo frío confinado (PUE estimado 1,6) & Subdoc. 4, 4.3.1.2 y 4.3.1.3; Subdoc. 4, Anexos (cálculo eléctrico); Formulario T-14, paquete 1.3.1.2. Evidencia: PUE del recinto on-premise e intensidad de carbono de la región de nube \\
RT-15.05 Deseable & Se valorará la elección de regiones de nube con menor intensidad de carbono cuando la latencia y la regulación lo permitan, con el análisis comparativo correspondiente. & Cumple parcialmente & Declaración de la intensidad de carbono de las regiones Chile Central y Brazil South, con su fuente y su fecha & Subdoc. 4, 4.2.2.1; Formulario T-14, paquete 1.3.1.2. Falta: compromiso de elegir la región de menor intensidad de carbono cuando la latencia y la regulación lo permitan, con su análisis comparativo. Evidencia: Intensidad de carbono de Chile Central y Brazil South, con su fuente y su fecha. \\
RT-15.06 Deseable & Se valorará la definición de metas de reducción del consumo durante la Operación, con medición y reporte anual. & Cumple & Informe anual de consumo y huella de la solución (paquete 1.8.7.3) & Formulario T-14, paquete 1.8.7.3. Evidencia: Informe anual de consumo y huella de la solución con método de cálculo declarado \\
RT-15.07 Obligatorio & Las certificaciones se acreditarán con copia del certificado vigente y con el código de verificación del organismo emisor cuando exista. & Cumple & Certificaciones institucionales ISO 9001, ISO/IEC 27001, ISO 22301, ISO/IEC 20000-1, ISO/IEC 27017 y 27018, ISO 45001 e ISO 14001; ISO/IEC 42001 en plan & Subdoc. 1, 1.4 (Tabla \refext{S1}{tab:certificaciones-institucionales}); Subdoc. 1, Anexo E del Subdocumento 1; Formulario T-14, paquete 1.8.7.2. Evidencia: Copias de certificados vigentes adjuntas en el Sobre N.º 1 e informe anual de certificaciones vigentes \\
RT-15.08 Obligatorio & Las personas certificadas formarán parte del equipo efectivamente asignado al PROYECTO, con dedicación declarada. No se aceptará acreditar personal que no participe. & Cumple & Equipo clave nominado de nueve personas con certificaciones por integrante & Subdoc. 1, Anexo E del Subdocumento 1 (certificaciones y competencias de cada integrante del equipo clave); Formulario T-14, paquete 1.1.7.1 y Tabla \refext{T14}{tab:roles-equipo-clave-dedicacion-paquetes-cargo} del Formulario T-14. Evidencia: Certificados individuales y dedicación por persona \\
RT-15.09 Obligatorio & El Adjudicatario mantendrá vigentes estas certificaciones durante todo el período contractual y las repondrá ante la salida de una persona certificada, conforme al Artículo 76° de las Bases administrativas. & Cumple parcialmente & Reposición del personal certificado del equipo clave con igual certificación y traslape de quince días hábiles (Art. 76) & Formulario T-14, paquete 1.1.7.1 (reposición con igual certificación y traslape de quince días hábiles, Art. 76.2). Falta: mantener vigentes las certificaciones del equipo clave durante todo el período contractual, con renovación antes de su vencimiento. Evidencia: Registro de certificaciones del equipo clave y compromiso de reposición (Art. 76) \\
\end{tabla}
\begin{tabla}{I{0.134}I{0.222}I{0.110}I{0.186}I{0.348}}{Capítulo 16. Módulos transversales obligatorios.}{tab:capitulo-16-modulos-transversales-obligatorios}
  {\cab{ID requerimiento} & \cab{Descripción} & \cab{Estado} & \cab{Componente que lo satisface} & \cab{Sección de la propuesta y evidencia}}
  {\emph{Fuente:} elaboración propia a partir de las Bases técnicas transversales.}
RT-16.01 Obligatorio & La solución dispondrá de un módulo de administración que permita al Cliente, sin intervención del Adjudicatario, gestionar personas usuarias, roles, permisos, unidades organizacionales y sus jerarquías. & Cumple & Módulo Administración y Servicios Comunes (usuarios, roles y jerarquías) & Subdoc. 3, 3.4.3; Formulario T-14, paquete 1.4.3.1. Evidencia: Prueba de administración de usuarios y roles por el Cliente sin apoyo del Adjudicatario (RF-ADM-002) \\
RT-16.02 Obligatorio & Las reglas de negocio parametrizables —umbrales, plazos, montos, tolerancias, catálogos, listas de valores, textos de notificación— serán configurables desde la interfaz de administración, con control de versiones y registro de quién cambió qué y cuándo. & Cumple & Módulo Administración y Servicios Comunes (parametrización con control de versiones) & Subdoc. 3, 3.4.3; Formulario T-14, paquete 1.4.3.1. Evidencia: Catálogo de parámetros y prueba de cambio con registro de quién, cuándo y qué (RF-ADM-003) \\
RT-16.03 Obligatorio & Todo cambio de parámetro con impacto operacional requerirá aprobación de un segundo perfil y quedará registrado con su justificación. & Cumple & Módulo Administración y Servicios Comunes (parametrización versionada con segundo aprobador) & Subdoc. 5, 5.1.6; Subdoc. 5, Anexos A.3 (Tabla A.12, Parámetro / Regla); Subdoc. 4, 4.1.8.5 (Tabla \refext{S4}{tab:amenazas-mayor-consecuencia-componentes-expuestos-control}, fila de alteración de un parámetro); Subdoc. 2, Anexo A.1 (RF-ADM-004); Formulario T-14, paquete 1.4.3.1. Evidencia: Prueba de cambio de parámetro con doble aprobación (RF-ADM-004) \\
RT-16.04 Obligatorio & El Proponente declarará expresamente qué elementos son parametrizables y cuáles requieren desarrollo. Presentar como parametrizable lo que exige desarrollo se evaluará como observación grave. & Cumple & Módulo Administración y Servicios Comunes (catálogos y reglas administrables como Parámetro / Regla versionado) & Subdoc. 5, Anexos A.4 (Tabla A.14); Subdoc. 2, Anexo A.1 (RF-ADM-003); Formulario T-14, paquete 1.4.3.1. Evidencia: Catálogo de parámetros: qué es parametrizable y qué requiere desarrollo (RF-ADM-003) \\
RT-16.05 Deseable & Existirá un ambiente de simulación que permita probar el efecto de un cambio de parámetro antes de aplicarlo a producción. & Cumple & Ambiente de Preproducción con la configuración de producción & Subdoc. 4, 4.1.3.4 y 4.2.3.1 (Figura \refext{S4}{fig:nube-principal-chile-central-acceso-produccion}); Formulario T-14, paquetes 1.4.3.1 y 1.6.1.1. Evidencia: Prueba de un cambio de parámetro ensayado en Preproducción antes de su aplicación en Producción \\
RT-16.06 Obligatorio & Toda operación que cree, modifique o elimine información quedará registrada con identificación de la persona o del sistema que la ejecutó, fecha y hora con zona horaria, origen, valores anteriores y valores posteriores. & Cumple & Módulo Administración y Servicios Comunes (auditoría inalterable) & Subdoc. 3, 3.4.3; Formulario T-14, paquete 1.4.3.1. Evidencia: Auditoría sobre muestra de operaciones (RF-ADM-005) \\
RT-16.07 Obligatorio & El registro de auditoría será inalterable y no podrá ser modificado ni eliminado por ningún perfil, incluido el administrador de la plataforma. & Cumple & Módulo Administración y Servicios Comunes (auditoría inalterable) & Subdoc. 3, 3.4.3; Formulario T-14, paquete 1.4.3.1. Evidencia: Prueba de intento de modificación o borrado, incluido el administrador (RF-ADM-005) \\
RT-16.08 Obligatorio & El Cliente podrá consultar y exportar la auditoría desde la interfaz, con filtros por persona, período, entidad y tipo de operación, sin requerir acceso a la base de datos. & Cumple parcialmente & Módulo Administración y Servicios Comunes (auditoría inalterable con copia de consulta) & Subdoc. 5, 5.2.1 y Anexos A.3 (Asiento de auditoría) y A.4 (Tabla A.15); Subdoc. 2, Anexo A.1 (RF-ADM-006); Formulario T-14, paquete 1.4.3.1. Falta: filtros por entidad y tipo de operación, y consulta sin acceso a la base de datos. Evidencia: Prueba de consulta y exportación de auditoría (RF-ADM-006) \\
RT-16.09 Según caso & Las consultas a información sensible quedarán registradas, no sólo las modificaciones. & Cumple & Módulo Seguridad y Salud Ocupacional (registro de cada consulta) & Subdoc. 3, 3.4.2.6; Formulario T-14, paquetes 1.4.3.1 y 1.4.4.2. Evidencia: Auditoría de una muestra de accesos a información sensible (RNF-SEG-002) \\
RT-16.10 Según caso & El período de retención de la auditoría será el que fije el caso y, en su defecto, no inferior a cinco años. & Cumple & Política de retención por categoría de registro (auditoría y trazabilidad 10 años, operacionales 5 años, tributarios 7 años en SAP) & Subdoc. 5, 5.2.4 (Tabla \refext{S5}{tab:retencion-categoria-entidad}); Subdoc. 2, Anexo A.2 (RNF-SEG-005); Subdoc. 4, Anexos (matriz de controles); Formulario T-14, paquetes 1.3.2.3 y 1.4.3.1. Evidencia: Prueba de recuperación de registros en el extremo de su plazo (RNF-SEG-005) \\
RT-16.11 Obligatorio & La solución soportará flujos de trabajo con estados, transiciones, responsables, plazos, escalamiento automático por vencimiento y delegación por ausencia. & Cumple & Módulo Administración y Servicios Comunes (flujos de trabajo con plazos y escalamiento) & Subdoc. 3, 3.4.3; Formulario T-14, paquete 1.4.3.1. Evidencia: Prueba de flujo con escalamiento por vencimiento (RF-ADM-008) \\
RT-16.12 Obligatorio & Los flujos serán configurables por el Cliente sin desarrollo, al menos en lo relativo a responsables, plazos y niveles de aprobación. & Cumple & Módulo Administración y Servicios Comunes (definición de flujo como parámetro versionado, con casos, tareas y transiciones) & Subdoc. 5, Anexos A.4 (Tabla A.14, Definición de flujo) y A.3 (Caso de flujo, Tarea de flujo y Transición); Subdoc. 3, 3.4.3; Subdoc. 2, Anexo A.1 (RF-ADM-008); Formulario T-14, paquete 1.4.3.1. Evidencia: Prueba de configuración de un flujo por el Cliente (RF-ADM-008) \\
RT-16.13 Obligatorio & Toda solicitud pendiente será visible para su responsable en una bandeja de tareas unificada, con priorización y alerta de vencimiento. & Cumple & Módulo Administración y Servicios Comunes (tareas de flujo asignadas a cuentas, con fecha límite) & Subdoc. 5, Anexos A.3 (Tarea de flujo); Subdoc. 2, Anexo A.1 (RF-ADM-009); Formulario T-14, paquete 1.4.3.1. Evidencia: Prueba de bandeja de tareas (RF-ADM-009) \\
RT-16.14 Deseable & El motor de reglas permitirá definir condiciones de negocio evaluables sin recopilación, con trazabilidad de qué regla se aplicó a cada transacción. & Cumple & Reglas de cálculo y de flujo versionadas como Parámetro / Regla; motor de reglas deterministas del clasificador de detenciones & Subdoc. 5, 5.2.1 y Anexos A.4 (Tabla A.14); Subdoc. 13, 13.2.2; Formulario T-14, paquete 1.4.3.1. Evidencia: Prueba de cambio de una regla sin recompilar y consulta de la regla aplicada a una transacción \\
RT-16.15 Obligatorio & La solución gestionará documentos con versionado, metadatos, control de acceso, búsqueda por contenido y por metadato, y previsualización sin descarga. & Cumple & Módulo Administración y Servicios Comunes (gestión documental) & Subdoc. 3, 3.4.3; Formulario T-14, paquete 1.4.3.1. Evidencia: Prueba de versionado, búsqueda y vista previa (RF-ADM-010) \\
RT-16.16 Obligatorio & Los documentos se almacenarán cifrados, con verificación de integridad y retención conforme a la política del caso. & Cumple & Gestión documental sobre almacén de objetos con cifrado en reposo, hash de contenido, inmutabilidad y retención por categoría & Subdoc. 4, 4.1.8.4 (Tabla \refext{S4}{tab:amenazas-mayor-consecuencia-componentes-expuestos-control}) y 4.2.2.2; Subdoc. 5, 5.2.4 (Tabla \refext{S5}{tab:retencion-categoria-entidad}) y Anexos A.3 (Documento); Subdoc. 4, Anexos (matriz de controles); Subdoc. 2, Anexo A.1 (RF-ADM-010); Formulario T-14, paquetes 1.3.2.3 y 1.4.3.1. Evidencia: Revisión de configuración de almacenamiento y de retención \\
RT-16.17 Según caso & La solución soportará firma electrónica conforme a la Ley N° 19.799, con firma avanzada para los actos que el caso lo requiera, y verificación de validez del certificado al momento de la firma. & Cumple & Firma electrónica avanzada de un prestador acreditado (Ley N.º 19.799) en actas de embarque, registros de seguridad, reportes ambientales y gestión documental & Subdoc. 4, 4.1.1.3 y 4.1.7.1 (Tabla \refext{S4}{tab:llamadas-sincronas-entre-componentes-umbral}); Subdoc. 3, 3.4.2.6 y 3.4.3; Subdoc. 5, Anexos A.3 (Documento); Subdoc. 2, Anexo A.1 (RF-TRZ-017); Formulario T-14, paquete 1.4.3.1. Evidencia: Verificación de firma y de validez del certificado al emitirse (RNF-CUM-002) \\
RT-16.18 Según caso & Se generará el sello de tiempo y se conservará la evidencia de firma que permita verificar el documento con posterioridad al vencimiento del certificado. & Cumple & Firma electrónica avanzada con sello de tiempo del prestador acreditado y evidencia de firma inmutable & Subdoc. 5, Anexos A.3 (Documento: sello\_tiempo y evidencia\_firma); Subdoc. 4, 4.1.7.1 (Tabla \refext{S4}{tab:llamadas-sincronas-entre-componentes-umbral}); Subdoc. 2, Anexo A.2 (RNF-CUM-002); Formulario T-14, paquete 1.4.3.1. Evidencia: Verificación de documentos con certificado vencido (RNF-CUM-002) \\
RT-16.19 Obligatorio & La solución generará documentos a partir de plantillas administrables por el Cliente, con datos de la transacción y salida en formato abierto. & Cumple & Gestión documental con plantillas versionadas y variables administrables por el Cliente & Subdoc. 5, Anexos A.4 (Tabla A.15, Plantilla documental o de notificación) y A.3 (Documento); Subdoc. 2, Anexo A.1 (RF-NOT-001); Formulario T-14, paquete 1.4.3.1. Evidencia: Prueba de generación desde plantilla con salida en formato abierto (RF-ADM-011) \\
RT-16.20 Obligatorio & La solución enviará notificaciones por al menos tres canales: correo electrónico, notificación en la aplicación y mensajería instantánea o SMS, según lo que el caso requiera. & Cumple & Módulo Administración y Servicios Comunes (notificaciones multicanal, incluida radio y megafonía) & Subdoc. 3, 3.4.3; Formulario T-14, paquete 1.4.3.1. Evidencia: Prueba de envío por cada canal comprometido (RF-NOT-001) \\
RT-16.21 Obligatorio & Las plantillas de notificación serán administrables por el Cliente, con variables de la transacción, y versionadas. & Cumple & Módulo Administración y Servicios Comunes (plantillas de notificación versionadas con variables) & Subdoc. 5, Anexos A.4 (Tabla A.15); Subdoc. 13, Anexos (traspaso del bloque 2); Subdoc. 2, Anexo A.1 (RF-NOT-003); Formulario T-14, paquete 1.4.3.1. Evidencia: Prueba de edición de plantillas con variables (RF-NOT-003) \\
RT-16.22 Obligatorio & Cada persona usuaria podrá configurar sus preferencias de canal y de frecuencia, respetando las notificaciones que el Cliente defina como obligatorias. & Cumple & Módulo Administración y Servicios Comunes (preferencias de canal y frecuencia, con obligaciones no cancelables) & Subdoc. 5, Anexos A.4 (Tabla A.15, Preferencia e informe programado); Subdoc. 2, Anexo A.1 (RF-NOT-002); Formulario T-14, paquete 1.4.3.1. Evidencia: Prueba de preferencias con notificaciones obligatorias (RF-NOT-002) \\
RT-16.23 Obligatorio & El envío será asíncrono, con reintento ante falla, control de duplicados y registro de entrega, apertura y error por cada mensaje. & Cumple & Módulo Administración y Servicios Comunes (envío asíncrono con reintento, deduplicación e historial de entregas) & Subdoc. 4, 4.1.7.1 (Tabla \refext{S4}{tab:llamadas-sincronas-entre-componentes-umbral}); Subdoc. 5, Anexos A.4 (Tabla A.15, Notificación y entrega); Subdoc. 2, Anexo A.2 (RNF-DES-018); Formulario T-14, paquete 1.4.3.1. Evidencia: Prueba de envío con caída simulada del proveedor de un canal (RNF-DES-018) \\
RT-16.24 Obligatorio & El Proponente declarará el proveedor de cada canal, su costo unitario, su volumen proyectado y el tratamiento del costo variable en la Oferta Económica. & Cumple parcialmente & Canales de terreno y portales con marca y modelo de referencia en el T-11; la Botonera de Pala se fabrica según especificación de Overfore & Subdoc. 4, 4.1.4.1 (Tabla \refext{S4}{tab:diferencias-entre-puerta-exposicion-puerta-operacional}) y 4.2.1; Formulario T-11 (T11-67 a T11-71); OVERFORE-Subdocumento4-Anexos, Anexo C (Tabla C.2); Oferta Económica (costo unitario, volumen proyectado y tratamiento del costo variable de cada canal). Falta: fabricante de la Botonera de Pala. Evidencia: Declaración de proveedor, costo unitario y volumen por canal. \\
RT-16.25 Obligatorio & Las notificaciones respetarán la normativa de comunicaciones comerciales y permitirán la baja cuando corresponda. & Cumple & Módulo Administración y Servicios Comunes (baja de notificaciones informativas sin suprimir las obligatorias) & Subdoc. 5, Anexos A.4 (Tabla A.15); Subdoc. 2, Anexo A.1 (RF-NOT-002); Formulario T-14, paquete 1.4.3.1. Evidencia: Prueba de baja de notificaciones informativas conforme al art. 28 B de la Ley N.º 19.496 y la Ley N.º 21.719 \\
RT-16.26 Deseable & Se valorará la integración con canales conversacionales que permitan a la persona usuaria responder y ejecutar acciones desde el propio canal. & Por cumplir & Módulo de Administración y Servicios Comunes con notificaciones salientes por correo, aplicación, mensajería y mensaje de texto & Subdoc. 4, 4.1; Formulario T-14, paquete 1.4.3.1. Falta: integración con canales conversacionales que permitan responder y ejecutar acciones desde el propio canal. Evidencia: Notificaciones salientes por correo, aplicación, mensajería y mensaje de texto (1.4.3.1). \\
RT-16.27 Obligatorio & La solución dispondrá de búsqueda global con indexación de texto completo, tolerancia a errores de escritura, filtros facetados y respeto del control de acceso de la persona que busca. & Cumple parcialmente & Índice de búsqueda global OpenSearch en el grupo de datos, filtrado por rol & Subdoc. 4, 4.1.1.2 (Tabla \refext{S4}{tab:tecnologias-capa-6-alternativa-evaluada}, OpenSearch 3); OVERFORE-Subdocumento4-Anexos, Anexo A (T11-18); Subdoc. 5, 5.4.2 (Tabla \refext{S5}{tab:indices-tipo-consulta-sostienen}); Subdoc. 2, Anexo A.1 (RF-ADM-012); Formulario T-14, paquete 1.4.3.1. Falta: filtros facetados. Evidencia: Prueba de búsqueda con errores de tipeo y control de acceso (RF-ADM-012) \\
RT-16.28 Obligatorio & Los listados serán ordenables, filtrables, paginados y exportables en formatos abiertos, con el filtro aplicado reflejado en la exportación. & Cumple parcialmente & Módulo Administración y Servicios Comunes (exportación en formatos abiertos y trabajos de exportación con filtro) & Subdoc. 5, 5.2.2 (Tabla \refext{S5}{tab:vias-explotacion-informacion}) y Anexos A.4 (Tabla A.15, Trabajo de exportación); Subdoc. 2, Anexo A.1 (RF-ADM-013); Formulario T-14, paquete 1.4.3.1. Falta: ordenación, filtrado y paginación de listados, y reflejo del filtro en la exportación. Evidencia: Prueba de exportación con el filtro aplicado (RF-ADM-013) \\
RT-16.29 Obligatorio & Las exportaciones de gran volumen se procesarán de forma asíncrona, con notificación al completarse y sin bloquear la sesión. & Cumple parcialmente & Trabajos de exportación procesados por Administración con estado, archivo, errores y resultado parcial & Subdoc. 5, Anexos A.4 (Tabla A.15, Trabajo de importación o exportación); Subdoc. 2, Anexo A.1 (RF-ADM-013); Formulario T-14, paquete 1.4.3.1. Falta: aviso al completarse la exportación. Evidencia: Prueba de exportación de más de 10.000 registros sin bloquear la sesión (RF-ADM-013) \\
RT-16.30 Obligatorio & Toda exportación de información sensible quedará registrada en la auditoría, con identificación de quién exportó qué y cuándo. & Cumple & Módulo Administración y Servicios Comunes (auditoría que registra consultas y exportaciones sensibles) & Subdoc. 5, Anexos A.4 (Tabla A.15, Auditoría de origen) y A.3 (Asiento de auditoría); Subdoc. 4, 4.1.8.5; Subdoc. 2, Anexo A.1 (RF-ADM-006); Formulario T-14, paquetes 1.4.3.1 y 1.4.4.2. Evidencia: Revisión del registro de exportaciones de información sensible (RF-ADM-006) \\
RT-16.31 Según caso & La solución dispondrá de un portal público con la información que el caso determine, accesible sin autenticación, con los mismos estándares de accesibilidad y desempeño que el resto de la plataforma. & Cumple & Portal de Contratistas y Portal para Compradores (autenticados; sin portal abierto a la ciudadanía) & Subdoc. 3, 3.1.1, 3.4.2.5 y 3.4.2.2; Formulario T-14, paquete 1.4.5.2. Evidencia: Prueba de acceso autenticado a ambos portales en la marcha blanca (RF-ACR-005, RF-TRZ-016) \\
RT-16.32 Obligatorio & Las personas usuarias externas dispondrán de autoatención para las consultas de mayor frecuencia, evitando el contacto telefónico para operaciones simples. & Cumple & Portal de Contratistas y Portal para Compradores (consulta de estado, vencimientos, lotes y recuperación de acceso por autoservicio) & Subdoc. 3, 3.4.2.5; Subdoc. 4, 4.1.7.7; Subdoc. 2, Anexo A.1 (RF-ACR-015); Formulario T-14, paquete 1.4.5.2. Evidencia: Prueba de autoatención (RF-ACR-015) \\
RT-16.33 Obligatorio & El Proponente estimará la reducción esperada del volumen de atención asistida por efecto de la autoatención y comprometerá el indicador. & Cumple & Mesa de ayuda dimensionada con Erlang C y base de conocimiento, con estimación de reducción de atención asistida en el informe mensual & Formulario T-14, paquete 1.8.1.1. Evidencia: Estimación de reducción de atención asistida y compromiso de indicador en el informe mensual. \\
RT-16.34 Obligatorio & El portal público resistirá picos de tráfico sin degradar los servicios transaccionales internos, mediante aislamiento de recursos y caché. & Cumple & Capa de borde con rutas separadas de la ruta de terreno, mitigación en el borde y caché de consultas de los portales & Subdoc. 4, 4.1.4.2 y 4.1.8.5 (Tabla \refext{S4}{tab:amenazas-mayor-consecuencia-componentes-expuestos-control}); Subdoc. 5, 5.4.4 (Tabla \refext{S5}{tab:caches-vigencia-origen-vigencia}); Subdoc. 9, 9.2.4.5; Subdoc. 2, Anexo A.2 (RNF-DIS-009); Formulario T-14, paquetes 1.4.4.5 y 1.4.5.2. Evidencia: Prueba de carga de portales a 1,5 veces el peak, midiendo los tiempos de terreno (RNF-DIS-009) \\
\end{tabla}
\begin{tabla}{I{0.132}I{0.272}I{0.115}I{0.178}I{0.303}}{Capítulo 17. Canales digitales y movilidad.}{tab:capitulo-17-canales-digitales-movilidad}
  {\cab{ID requerimiento} & \cab{Descripción} & \cab{Estado} & \cab{Componente que lo satisface} & \cab{Sección de la propuesta y evidencia}}
  {\emph{Fuente:} elaboración propia a partir de las Bases técnicas transversales.}
RT-17.01 Según caso & La solución proveerá una aplicación móvil para los perfiles operacionales que el caso identifique, con funcionamiento sin conexión y sincronización diferida cuando la operación en terreno lo requiera. & Cumple & App SITCO Terreno, Botonera de Pala y tabletas (operación sin conexión y sincronización diferida) & Subdoc. 3, 3.3.3, 3.3.6 y 3.4.2.6; Formulario T-14, paquetes 1.4.2.2 y 1.4.5.1. Evidencia: Prueba de registro sin señal y sincronización en zona de sombra (RNF-DIS-003) \\
RT-17.02 Obligatorio & El Proponente declarará si la aplicación es nativa, híbrida o web progresiva, y justificará la decisión frente al requisito de operación desconectada, al acceso a periféricos y al costo de mantención. & Cumple & App SITCO Terreno híbrida (núcleo web en contenedor.NET con WebView2) & Subdoc. 4, 4.1.2.3 (ADR-06, Tabla \refext{S4}{tab:registro-decision-adr-06-canales-terreno}), 4.1.5 (Tabla \refext{S4}{tab:comparacion-aplicacion-web-progresiva-nativa-hibrida}) y 4.1.5.1; Formulario T-14, paquete 1.4.5.1. Evidencia: Declaración de tipo de aplicación (nativa, híbrida o web progresiva) y justificación \\
RT-17.03 Obligatorio & La aplicación soportará las versiones de sistema operativo móvil vigentes y las dos anteriores, con política de actualización declarada. & Cumple & App SITCO Terreno y Consola de Portería sobre Windows 11 IoT Enterprise LTSC 2024 administrado por Intune & Subdoc. 4, 4.1.1.4, 4.1.1.5 (Tabla \refext{S4}{tab:version-fin-soporte-plan-actualizacion-componentes}) y 4.1.2.3 (ADR-06); Formulario T-14, paquete 1.4.5.1. Evidencia: Prueba de la aplicación en la versión vigente de sistema operativo y las dos anteriores (RNF-POR-004) \\
RT-17.04 Obligatorio & La aplicación se distribuirá por las tiendas oficiales o mediante gestión de flota corporativa, con firma de la aplicación y verificación de integridad. & Cumple parcialmente & Distribución firmada de la aplicación mediante gestión de dispositivos (Intune) con verificación de firma. Falta: rechazo automático de las versiones no firmadas. & Subdoc. 4, 4.1.1.3 y 4.1.5.1; Subdoc. 4, Anexos (distribución firmada por Intune); Formulario T-14, paquete 1.4.5.1. Evidencia: Revisión del proceso de distribución e intento de instalar una versión no firmada (RNF-POR-003) \\
RT-17.05 Obligatorio & La información almacenada en el dispositivo estará cifrada, con borrado remoto y bloqueo ante pérdida o desvinculación de la persona usuaria. & Cumple & Dispositivos con registro local cifrado (SQLCipher, TPM) y gestión remota con bloqueo y borrado & Subdoc. 4, 4.1.8.4 (Tabla \refext{S4}{tab:mecanismo-cifrado-custodia-claves-ambito}) y 4.1.5.1; Subdoc. 4, 4.1.1.3; Subdoc. 2, Anexo A.1 (RF-CNX-004); Formulario T-14, paquete 1.4.5.1. Evidencia: Prueba de extracción de almacenamiento y de borrado remoto (RNF-SEG-016) \\
RT-17.06 Según caso & La aplicación integrará los periféricos que el caso requiera: cámara, lector de códigos, NFC, GPS, impresora de etiquetas, balanza o báscula. & Cumple & Tabletas con código QR y credencial de proximidad; integraciones con báscula, fatiga, telemetría e instrumentación & Subdoc. 3, 3.3.3, 3.3.2.2 y 3.4.2.7; Formulario T-14, paquete 1.4.5.1. Evidencia: Revisión de la integración de cada periférico exigido por el caso \\
RT-17.07 Obligatorio & El consumo de datos móviles y de batería se optimizará y se declarará el consumo estimado por turno de trabajo. & Cumple parcialmente & Botonera de Pala y App SITCO Terreno, con declaración del consumo de datos y de batería por turno de trabajo & Subdoc. 4, 4.2.1.3; Subdoc. 4, 4.2.1 y Formulario T-11; Formulario T-14, paquete 1.4.5.1. Falta: optimización del consumo de datos móviles y de batería. Evidencia: Declaración de consumo de datos y batería por turno de trabajo (1.4.5.1). \\
RT-17.08 Deseable & Se valorará la disponibilidad de una versión de la interfaz para dispositivos de bajo costo o de generaciones anteriores, ampliando la cobertura de personas usuarias. & Por cumplir & App SITCO Terreno en los dispositivos de terreno especificados, sin versión para dispositivos de bajo costo o de generaciones anteriores & Subdoc. 4, 4.2.1; Formulario T-14, paquete 1.4.5.1. Falta: versión de la interfaz para dispositivos de bajo costo o de generaciones anteriores. Evidencia: Especificación de los dispositivos de terreno (Subdoc. 4, 4.2.1.3). \\
\end{tabla}
\begin{tabla}{I{0.136}I{0.226}I{0.074}I{0.173}I{0.391}}{Capítulo 18. Inteligencia artificial y automatización.}{tab:capitulo-18-inteligencia-artificial-automatizacion}
  {\cab{ID requerimiento} & \cab{Descripción} & \cab{Estado} & \cab{Componente que lo satisface} & \cab{Sección de la propuesta y evidencia}}
  {\emph{Fuente:} elaboración propia a partir de las Bases técnicas transversales.}
RT-18.01 Obligatorio & El Proponente declarará cada componente de inteligencia artificial: propósito, modelo o servicio empleado, proveedor, versión y ubicación de procesamiento de los datos. & Cumple & Dos componentes de inteligencia artificial: clasificador asistido de detenciones (innovación 2 del Subdocumento 13) y verificación facial de portería & Subdoc. 13, 13.2.2 y 13.2.3; OVERFORE-Subdocumento13-Anexos, Anexo B, B.5 (Tabla B.4); Formulario T-14, paquetes 1.4.4.4 (verificación facial en portería: propósito, proveedor, versión y lugar de procesamiento) y 1.7.2.1 (clasificador: modelo LightGBM 4.6.0, proyecto de código abierto, en Chile Central con respaldo en Brazil South). Evidencia: Registro de versiones del modelo \\
RT-18.02 Obligatorio & Se garantizará contractualmente que los datos del Cliente no serán utilizados para entrenar modelos de terceros, salvo autorización expresa y escrita. & Cumple & Clasificador asistido de detenciones y verificación facial de portería (dos componentes de inteligencia artificial) & Subdoc. 13, 13.2.2 y 13.2.3; OVERFORE-Subdocumento13-Anexos, Anexo B, B.5 (Tabla B.4, uso de los datos del Cliente); Formulario T-14, paquetes 1.4.4.4 (contrato con el proveedor del modelo de reconocimiento) y 1.7.2.1 (cláusula de datos del Cliente, incluido el contrato del científico de datos externo). Evidencia: Cláusula del contrato \\
RT-18.03 Obligatorio & Se documentarán los límites de uso, los casos en que el resultado requiere validación humana previa y el procedimiento de supervisión. & Cumple & Clasificador asistido de detenciones y verificación facial de portería (dos componentes de inteligencia artificial) & Subdoc. 13, 13.2.2 y 13.2.3; OVERFORE-Subdocumento13-Anexos, Anexo B, B.5 (Tabla B.4, límites de uso y validación humana previa); Formulario T-14, paquetes 1.4.4.4, 1.7.2.1 y 1.7.2.3. Evidencia: Descripción del flujo y registro de decisiones \\
RT-18.04 Obligatorio & Los riesgos de sesgo, alucinación, fuga de información y uso indebido se evaluarán y mitigarán conforme al NIST AI Risk Management Framework 1.0 y a la norma ISO/IEC 42001. & Cumple & Clasificador asistido de detenciones y verificación facial de portería (dos componentes de inteligencia artificial) & Subdoc. 13, 13.2.2 y 13.2.3 (sesgo, fuga de información y uso indebido); OVERFORE-Subdocumento13-Anexos, Anexo B, B.5 (Tabla B.4); Formulario T-14, paquetes 1.7.2.1 (evaluación según NIST AI RMF 1.0 e ISO/IEC 42001, con controles del clasificador) y 1.7.2.4 (sistema de gestión de la inteligencia artificial con alcance en el clasificador) y 1.4.4.4 (verificación facial de portería). Evidencia: Registro de riesgos del componente \\
RT-18.05 Obligatorio & Toda interacción relevante con un componente de inteligencia artificial quedará registrada para efectos de auditoría y trazabilidad, incluyendo la entrada, la salida y la decisión humana posterior. & Cumple & Clasificador asistido de detenciones y verificación facial de portería (dos componentes de inteligencia artificial) & Subdoc. 13, 13.2.2 y 13.2.3; OVERFORE-Subdocumento13-Anexos, Anexo B, B.5 (Tabla B.4, registro de entrada, causa sugerida, confianza y decisión); Formulario T-14, paquetes 1.4.4.4, 1.7.2.1 y 1.7.2.3. Evidencia: Tabla de auditoría del componente \\
RT-18.06 Obligatorio & El componente podrá desactivarse sin comprometer la operación del resto de la solución, y existirá un procedimiento manual de respaldo para la función que automatiza. & Cumple & Clasificador asistido de detenciones y verificación facial de portería (dos componentes de inteligencia artificial) & Subdoc. 13, 13.2.2 y 13.2.3; OVERFORE-Subdocumento13-Anexos, Anexo B, B.5 (Tabla B.4, desactivación sin afectar al resto; reglas y clasificación manual); Formulario T-14, paquetes 1.4.4.4, 1.7.2.1 y 1.7.2.3. Evidencia: Procedimiento de desactivación \\
RT-18.07 Obligatorio & Cuando el resultado se presente a una persona usuaria, se indicará expresamente que fue generado o sugerido de forma automática, y su nivel de confianza cuando el modelo lo provea. & Cumple & Clasificador asistido de detenciones y verificación facial de portería (dos componentes de inteligencia artificial) & Subdoc. 13, 13.2.2 y 13.2.3; OVERFORE-Subdocumento13-Anexos, Anexo B, B.5 (Tabla B.4, transparencia al usuario con la confianza de la propuesta); Formulario T-14, paquetes 1.4.4.4, 1.7.2.1 y 1.7.2.3. Evidencia: Pantalla de la App SITCO Terreno \\
RT-18.08 Obligatorio & Los modelos predictivos declararán sus variables de entrada, su métrica de desempeño, su línea base y su plan de reentrenamiento y de detección de deriva. & Cumple & Clasificador asistido de detenciones y verificación facial de portería (dos componentes de inteligencia artificial) & Subdoc. 13, 13.2.2, 13.2.3 y 13.2.5; OVERFORE-Subdocumento13-Anexos, Anexo B: B.2 (Tabla B.2, variables de entrada), B.4 (Tabla B.3, métricas, reentrenamiento y deriva) y B.6 (líneas base); Formulario T-14, paquetes 1.4.4.4, 1.7.2.1, 1.7.2.2 y 1.7.2.3. Evidencia: Informe de verificación de la innovación 2 \\
RT-18.09 Obligatorio & La responsabilidad por los resultados de los componentes de inteligencia artificial recae íntegramente en el Adjudicatario, conforme al Artículo 86° de las Bases administrativas. & Cumple & Clasificador asistido de detenciones y verificación facial de portería (dos componentes de inteligencia artificial) & Subdoc. 13, 13.2.2 y 13.2.3; OVERFORE-Subdocumento13-Anexos, Anexo B, B.5 (Tabla B.4, responsabilidad íntegra de Overfore conforme al Art. 86); Formulario T-14, paquetes 1.4.4.4, 1.7.2.1 y 1.7.2.3. Evidencia: Contrato de servicio \\
RT-18.10 Deseable & Se valorará el uso de automatización robótica de procesos o de agentes para tareas repetitivas de back office, con el ahorro de horas cuantificado y reflejado en la Oferta Económica. & No aplica & La solución no incorpora automatización robótica ni agentes & Subdoc. 3, 3.2.3.1 Evidencia: Deseable no comprometido \\
\end{tabla}
\begin{tabla}{I{0.143}I{0.266}I{0.113}I{0.226}I{0.252}}{Capítulo 19. Estructura y gobierno del proyecto.}{tab:capitulo-19-estructura-gobierno-proyecto}
  {\cab{ID requerimiento} & \cab{Descripción} & \cab{Estado} & \cab{Componente que lo satisface} & \cab{Sección de la propuesta y evidencia}}
  {\emph{Fuente:} elaboración propia a partir de las Bases técnicas transversales.}
RT-19.01 Obligatorio & El Adjudicatario constituirá una oficina de gestión del proyecto con metodología declarada, basada en el PMBOK del Project Management Institute e integrando prácticas ágiles donde el Proponente lo justifique. & Cumple & Oficina de gestión del proyecto dirigida por el Jefe de Proyecto, con metodología basada en el PMBOK y enfoque adaptativo en el desarrollo & Subdoc. 6, 6.1 y 6.1.1 y 6.1.2; Subdoc. 6, 6.2.1; Formulario T-14, paquetes 1.1.1.2 y 1.1.1.4. Evidencia: Oficina de gestión del proyecto y metodología declarada \\
RT-19.02 Obligatorio & El Jefe de Proyecto tendrá dedicación exclusiva durante toda la fase de implementación y facultades para comprometer al Adjudicatario en materias de ejecución. & Cumple & Jefe de Proyecto con dedicación exclusiva durante la implementación y facultades para comprometer al adjudicatario & Subdoc. 6, 6.1.2; Formulario T-14, paquete 1.1.1.4 (dedicación exclusiva del Jefe de Proyecto, meses 1 a 21, y facultades para comprometer a Overfore). Evidencia: Declaración de dedicación exclusiva del Jefe de Proyecto \\
RT-19.03 Obligatorio & Existirá un procedimiento formal de control de cambios conforme al Artículo 72° de las Bases administrativas, con registro de cambios, análisis de impacto y aprobación previa a la ejecución. & Cumple & Procedimiento formal de control de cambios con análisis de impacto en cinco días hábiles y aprobación previa del comité ejecutivo & Subdoc. 6, 6.1.3 y 6.1.5.1; Formulario T-14, paquete 1.1.2.3. Evidencia: Procedimiento de control de cambios conforme al Art. 72° \\
RT-19.04 Obligatorio & La gestión del riesgo seguirá la norma ISO 31000, con registro de riesgos vivo, revisión en cada Comité de Proyecto y planes de mitigación con responsable, plazo y disparador. & Cumple & Registro de riesgos revisado en cada sesión del comité de proyecto conforme a ISO 31000 & Subdoc. 6, 6.1.2; Formulario T-14, paquetes 1.1.3.1 y 1.1.3.2. Evidencia: fichas de los 24 riesgos con responsable, plazo y disparador (Formulario T-16, sección 4) y registro de riesgos vivo conforme a ISO 31000. \\
RT-19.05 Obligatorio & El Adjudicatario habilitará un espacio colaborativo accesible al Cliente, con la documentación del proyecto, los entregables, las actas, el registro de riesgos y el registro de cambios siempre actualizados. & Cumple & Espacio colaborativo del proyecto en SharePoint y Microsoft Teams del cliente, con documentos, entregables, actas y registros de riesgos y cambios & Subdoc. 6, 6.1.5.1 y 6.1.5.2 (Tablas \refext{T9}{tab:canales-control-comunicaciones} y \refext{T9}{tab:plan-comunicaciones} del Formulario T-9); Formulario T-14, paquete 1.1.1.3. Evidencia: Espacio colaborativo con documentación, actas, riesgos y cambios \\
RT-19.06 Obligatorio & El Adjudicatario entregará un informe mensual de avance con estado del cronograma, avance físico y financiero, entregables del período, desviaciones, riesgos, incidencias y compromisos del período siguiente. & Cumple parcialmente & Informe mensual de avance enviado al comité ejecutivo y a la contraparte técnica, con valor ganado & Subdoc. 6, 6.1.1, 6.1.3 y 6.1.5.2; Formulario T-14, paquete 1.1.2.2. Falta: avance financiero del período; el valor ganado se mide en horas-hombre. Evidencia: Informe mensual de avance \\
RT-19.07 Obligatorio & El avance se medirá con valor ganado, reportando el índice de desempeño del cronograma y del costo, y no mediante declaración cualitativa de porcentaje de avance. & Cumple & Informe mensual de avance con valor ganado, con índices de desempeño del cronograma y del costo por paquete & Subdoc. 6, 6.1.1; Formulario T-14, paquete 1.1.2.2. Evidencia: Reporte mensual de valor ganado con índices de desempeño de cronograma y costo. \\
RT-19.08 Obligatorio & El Cliente dispondrá de un tablero de estado del proyecto actualizado, accesible en cualquier momento. & Cumple & Tablero de estado en ClickUp publicado como pestaña en Teams, con acceso permanente de solo lectura para la contraparte técnica y los referentes & Subdoc. 6, 6.1.5.1 y 6.1.5.2; Formulario T-9, Tabla \refext{T9}{tab:plan-comunicaciones}; Formulario T-14, paquete 1.1.5.1. Evidencia: Tablero de estado del proyecto \\
RT-19.09 Obligatorio & Las actas de todos los comités del Artículo 71° de las Bases administrativas se levantarán dentro de los dos días hábiles siguientes y registrarán acuerdos, responsables y plazos. & Cumple & Actas de cada instancia de gobierno publicadas en SharePoint dentro de dos días hábiles, con acuerdos, responsables y plazos & Subdoc. 6, 6.1.2, 6.1.5.2 y 6.1.5.4; Formulario T-14, paquetes 1.1.5.1 y 1.8.7.4. Evidencia: Actas de comités dentro de dos días hábiles \\
RT-19.10 Obligatorio & Toda desviación superior al 10 \% en un hito se comunicará dentro de los cinco días hábiles de detectada, con plan de recuperación. & Cumple & Aviso de desviación mayor al 10 \% en un hito al comité ejecutivo dentro de cinco días hábiles desde su detección, con plan de recuperación & Subdoc. 6, 6.1.1; Formulario T-14, paquete 1.1.2.2. Evidencia: Aviso de desviación de hito con plan de recuperación (Formulario T-9, Tabla \refext{T9}{tab:plan-comunicaciones}) \\
\end{tabla}
\begin{tabla}{I{0.143}I{0.263}I{0.076}I{0.247}I{0.271}}{Capítulo 20. Implantación, pruebas y aceptación.}{tab:capitulo-20-implantacion-pruebas-aceptacion}
  {\cab{ID requerimiento} & \cab{Descripción} & \cab{Estado} & \cab{Componente que lo satisface} & \cab{Sección de la propuesta y evidencia}}
  {\emph{Fuente:} elaboración propia a partir de las Bases técnicas transversales.}
RT-20.01 Obligatorio & El Proponente presentará un plan de implantación por sitio y por perfil, coherente con el cronograma obligatorio del Artículo 17° de las Bases administrativas. & Cumple & Implantación en olas por área y sitio: Etapa 1 (portería, mina, planta, concentrado y puerto) y Etapa 2 (seguridad, mantenimiento y cumplimiento ambiental) & Subdoc. 3, 3.2.2.4; Formulario T-14, paquetes 1.1.2.1 y 1.6.4.1. Evidencia: Plan de implantación y de marcha blanca (Formulario T-18) \\
RT-20.02 Obligatorio & La estrategia de paso a producción será gradual y reversible. Se declarará el criterio de avance entre olas y el procedimiento de reversión, con su tiempo de ejecución. & Cumple & Paso a producción por olas con criterio de avance entre olas y reversión a planillas y radio sin perder lo registrado, con tiempo máximo por tipo de reversión & Subdoc. 3, 3.2.2.5 y 3.1.3; Formulario T-14, paquetes 1.6.4.1 y 1.6.4.2. Evidencia: Ensayo de reversión cronometrado (TRA-006) \\
RT-20.03 Obligatorio & Durante la marcha blanca, la solución convivirá con la operación vigente del Cliente, con conciliación diaria y sin doble digitación no declarada. & Cumple & Marcha blanca con planillas activas y comparación diaria con SITCO & Subdoc. 3, 3.1.3 y 3.2.2.4; Formulario T-14, paquetes 1.6.3.1 y 1.6.3.2. Evidencia: Informe de conciliación diaria entre SITCO y el registro vigente, y acta de apagado del método anterior (TRA-008) \\
RT-20.04 Obligatorio & El Proponente definirá los indicadores diarios que se medirán durante la marcha blanca y sus umbrales de avance, conforme al Artículo 17.3 de las Bases administrativas. & Cumple & Tablero diario de la marcha blanca con indicadores y umbrales de cierre por cada condición del Art. 17.3, consultable por la Contraparte Técnica & Subdoc. 3, 3.2.5; Formulario T-14, paquetes 1.6.3.1 y 1.6.3.2. Evidencia: Tablero diario de marcha blanca \\
RT-20.05 Obligatorio & Se dispondrá de acompañamiento en terreno durante las primeras semanas de cada ola, con dotación declarada y decreciente según la curva de adopción. & Cumple & Acompañamiento presencial en faena en ambos turnos, con dotación decreciente declarada por etapa & Subdoc. 3, 3.1.3; Formulario T-14, paquetes 1.6.3.3, 1.6.3.4 y 1.6.4.1. Evidencia: Plan de acompañamiento por ola \\
RT-20.06 Obligatorio & Se establecerá un período de estabilización con atención reforzada tras cada paso a producción, con dotación y duración declaradas y sin costo adicional. & Cumple & Estabilización en faena de 52 días tras el paso de la Etapa 1 (del 10 de abril al 31 de mayo de 2028) y de 50 días tras el paso de la Etapa 2 (del 20 de septiembre al 8 de noviembre de 2028), con atención reforzada y sin costo adicional & Formulario T-14, paquetes 1.6.4.1 (estabilización de la Etapa 1, 52 días, del 10 de abril al 31 de mayo de 2028), 1.6.4.2 (estabilización de la Etapa 2, 50 días, del 20 de septiembre al 8 de noviembre de 2028) y 1.6.4.3 (soporte de la Etapa 1 entre su estabilización y la operación). Evidencia: Plan de estabilización con dotación y duración sin costo adicional (Formulario T-14, paquetes 1.6.4.1 y 1.6.4.2) \\
RT-20.07 Obligatorio & Existirá una definición de terminado acordada con el Cliente, aplicable a cada entregable, que incluya código, pruebas, documentación, seguridad y despliegue. & Cumple & Definición de terminado acordada con el cliente (código, pruebas, documentación, seguridad y despliegue) aplicada a cada entregable & Subdoc. 9, 9.2.6 y Anexo B, Tabla B.1 (Subdoc. 9 Anexos); Subdoc. 6, 6.2.3; Formulario T-14, paquete 1.1.4.2. Evidencia: Definición de terminado acordada con el Cliente \\
RT-20.08 Obligatorio & El protocolo de aceptación de cada hito se formalizará conforme al Formulario T-17, con criterios objetivos y verificables. & Cumple & Protocolo de aceptación por hito conforme al Formulario T-17, con entregable, prueba, valor exigido y consecuencia del incumplimiento & Subdoc. 9, 9.2.6; Formulario T-17; Formulario T-14, paquete 1.1.4.3. Evidencia: Protocolo de aceptación (Formulario T-17) \\
\end{tabla}
\begin{tabla}{I{0.140}I{0.263}I{0.112}I{0.218}I{0.267}}{Capítulo 21. Operación, mantención y soporte.}{tab:capitulo-21-operacion-mantencion-soporte}
  {\cab{ID requerimiento} & \cab{Descripción} & \cab{Estado} & \cab{Componente que lo satisface} & \cab{Sección de la propuesta y evidencia}}
  {\emph{Fuente:} elaboración propia a partir de las Bases técnicas transversales.}
RT-21.01 Obligatorio & El Adjudicatario dispondrá de un centro de operaciones de red con cobertura 24×7×365, propio o subcontratado, y declarará su ubicación, dotación por turno y procedimientos. & Cumple & Centro de operaciones 24x7 propio en Santiago (casa matriz), con un puesto de mesa de ayuda, uno de NOC y uno de SOC por turno de 12 horas, sobre nube, Nodo de Faena, gabinetes y enlaces & Subdoc. 3, 3.2.2.6; Subdoc. 1, 1.2; Subdoc. 1, Anexo F del Subdocumento 1; Formulario T-14, paquetes 1.8.2.1 y 1.8.2.2. Evidencia: Declaración del centro de operaciones \\
RT-21.02 Obligatorio & Se designará un gerente de servicio dedicado como contraparte permanente del Cliente durante la fase de Operación. & Cumple & Gerente de servicio dedicado a la operación (meses 21 a 56) como contraparte permanente del cliente, que dirige la mesa de ayuda y los centros de operaciones & Subdoc. 1, 1.5; Subdoc. 1, Anexo E del Subdocumento 1; Formulario T-14, paquetes 1.1.1.4 y 1.8.7.5. Evidencia: Nominación del gerente de servicio \\
RT-21.03 Obligatorio & El equipo de operación contará con especialistas por tecnología, nominados y con dedicación declarada. & Cumple & Nivel 2 de especialistas por dominio (mina, planta, red y plataforma) & Subdoc. 3, 3.2.2.6; Formulario T-14, paquete 1.8.3.1. Evidencia: Nómina de especialistas de operación \\
RT-21.04 Obligatorio & Los procedimientos de operación estarán documentados y serán ejecutables por el personal del Cliente tras la transferencia de conocimiento. & Cumple parcialmente & Libros de operación por escenario de falla, matriz de escalamiento y transferencia de conocimiento al equipo de TI del Cliente & Subdoc. 3, 3.1.3; Formulario T-14, paquete 1.1.10.2 (libros de operación, matriz de escalamiento y ejecución sin asistencia por el personal de TI del Cliente antes del cierre de cada marcha blanca). Falta: que el personal del Cliente ejecute sin asistencia el procedimiento de conmutación antes del cierre de la marcha blanca (RP-009). Evidencia: Ejecución de procedimientos críticos por el Cliente antes del cierre de la marcha blanca (RP-009) \\
RT-21.05 Obligatorio & El Proponente dispondrá de un centro de atención adecuado para soportar integralmente las consultas, propio o subcontratado con un operador de nivel acreditado. & Cumple & Mesa de ayuda de nivel 1, remota, 24x7 y en español, operada por el centro de operaciones propio de Overfore, con certificación ISO/IEC 20000-1 & Subdoc. 3, 3.2.2.6; Subdoc. 1, 1.2 y 1.4 (Tabla \refext{S1}{tab:certificaciones-institucionales}); Formulario T-14, paquetes 1.8.1.1 y 1.8.1.2. Evidencia: Declaración del centro de atención \\
RT-21.06 Obligatorio & El centro de atención cumplirá un tiempo de atención al usuario final de 80 \% antes de 20 segundos y una resolución al primer contacto de al menos 70 \%. La tasa de abandono no superará el 5 \%. & Cumple & Mesa de ayuda con metas mensuales de 80 \% de llamadas atendidas antes de 20 segundos, resolución al primer contacto de 70 \% o más y abandono de 5 \% o menos & Subdoc. 2 Anexos, RP-007; Subdoc. 3, 3.2.2.6; Formulario T-14, paquete 1.8.1.1. Evidencia: Medición mensual de nivel de atención, resolución al primer contacto y abandono \\
RT-21.07 Según caso & El horario mínimo de atención será de 8:00 a 20:00 en días hábiles, ampliado a 24×7 para los incidentes de severidad crítica y para la ventana operacional que defina el caso. & Cumple & Mesa de ayuda 24x7x365 en español con respuesta presencial en faena para incidentes críticos & Subdoc. 3, 3.2.2.6; Formulario T-14, paquetes 1.8.1.1 y 1.8.1.2 Evidencia: Informe mensual de nivel de servicio por severidad \\
RT-21.08 Obligatorio & La atención cubrirá orientación funcional de distintos grados de complejidad, desde preguntas simples hasta situaciones complejas, y las preguntas técnicas más frecuentes sobre la aplicación y su entorno de uso. & Cumple & Soporte en cuatro niveles (mesa, especialistas, expertos y fabricantes) & Subdoc. 3, 3.2.2.5; Formulario T-14, paquetes 1.6.5.2 y 1.8.1.1. Evidencia: Base de conocimiento y guías de resolución \\
RT-21.09 Obligatorio & Se mantendrá un registro histórico de las actividades de soporte que permita monitorear y gestionar los principales requerimientos, con análisis de tendencia mensual. & Cumple & Informe mensual de nivel de servicio revisado con la Contraparte Técnica & Subdoc. 3, 3.2.2.6; Subdoc. 2 Anexos, RNF-TRV-21.15; Formulario T-14, paquetes 1.8.1.1 y 1.8.3.1. Evidencia: Informe mensual con tendencia \\
RT-21.10 Obligatorio & Existirá un proceso de aprendizaje del centro de soporte que acumule conocimiento sobre las complejidades de la operación y las necesidades de las personas usuarias, reflejado en la base de conocimiento. & Cumple & Mesa de ayuda con base de conocimiento alimentada por cada consulta resuelta y análisis mensual de tendencias de los requerimientos & Formulario T-14, paquete 1.8.1.1. Evidencia: Base de conocimiento alimentada por cada consulta resuelta, con análisis mensual de tendencias de los requerimientos. \\
RT-21.11 Obligatorio & Se proveerán servicios de capacitación en línea en modalidad de autoformación, con registro de avance para efectos de monitoreo. & Cumple & Curso de autoformación en línea por perfil, con registro de avance por persona consultable por el Cliente & Subdoc. 2 Anexos, RNF-TRV-21.11; Formulario T-14, paquete 1.6.5.2. Evidencia: Curso de autoformación en línea por perfil con registro de avance por persona, consultable por el Cliente (RNF-TRV-21.11). \\
RT-21.12 Deseable & Se habilitarán espacios de interacción entre las entidades y personas usuarias del sistema, que permitan compartir experiencias de uso. & Por cumplir & Base de conocimiento alimentada por cada consulta resuelta, con métricas de uso y de utilidad & Formulario T-14, paquetes 1.8.1.1 y 1.8.7.5. Falta: espacios de interacción entre las entidades y las personas usuarias para compartir experiencias de uso. Evidencia: Base de conocimiento alimentada por cada consulta resuelta (1.8.1.1). \\
RT-21.13 Obligatorio & Los indicadores clave del proceso de soporte estarán abiertos al Cliente, y los indicadores específicos serán accesibles según nivel de autorización. & Cumple & Informe mensual de nivel de servicio revisado con la Contraparte Técnica & Subdoc. 3, 3.2.2.6; Subdoc. 2 Anexos, RNF-TRV-21.13; Formulario T-14, paquetes 1.8.7.1 y 1.8.2.1 (tablero de observabilidad disponible para el Cliente en tiempo real). Evidencia: Demostración de visibilidad de indicadores de soporte según nivel de autorización (RT-21.13). \\
RT-21.14 Obligatorio & El Proponente dimensionará la dotación del centro de atención con fundamento cuantitativo, empleando teoría de colas o el modelo Erlang C, a partir del volumen de contactos proyectado del caso. & Cumple & Dimensionamiento de la mesa de ayuda con Erlang C: un analista logra 84 \% antes de 20 segundos y dos en el cambio de turno logran 91 \% & Subdoc. 2 Anexos, RP-007 y SUP-27; Formulario T-14, paquete 1.8.1.1 (1.250 contactos al mes, por encima del mínimo de 900). Evidencia: Cálculo Erlang C sobre 1.250 contactos mensuales, por encima del mínimo de 900 (RP-007). \\
RT-21.15 Obligatorio & Existirá un canal único de registro de incidentes y solicitudes, con número de ticket, clasificación por severidad y seguimiento del ciclo de vida completo hasta el cierre conforme. & Cumple & Mesa de ayuda de nivel 1 que registra y clasifica cada contacto por las cuatro severidades del Art. 78 & Subdoc. 3, 3.2.2.6; Subdoc. 2, Anexos (RNF-TRV-21.15); Formulario T-14, paquetes 1.8.1.1 y 1.8.3.1. Evidencia: Inspección del sistema de tickets (número único, clasificación por severidad y ciclo de vida hasta el cierre conforme). \\
RT-21.16 Según caso & Cuando el caso comprenda sitios alejados, los especialistas de niveles 2 y 3 deberán trasladarse cuando la resolución lo requiera, por el medio más rápido disponible y con disponibilidad de tiempo suficiente, sin alterar la atención normal del resto de los sitios. El costo del traslado está incluido en la oferta. & Cumple & Niveles 2 y 3 trasladables a faena, con costo incluido; personal acreditado y con examen de altura & Subdoc. 3, 3.2.2.6 y 3.2.3.3; Formulario T-14, paquete 1.8.1.2. Evidencia: Revisión de acreditaciones y exámenes de altura vigentes del personal (RP-008) \\
RT-21.17 Obligatorio & El cierre de un ticket requerirá confirmación de la persona usuaria o transcurso del plazo de confirmación automática declarado. & Cumple & Cierre de ticket con confirmación de la persona usuaria o confirmación automática a los 7 días corridos & Subdoc. 2 Anexos, RNF-TRV-21.15; Formulario T-14, paquete 1.8.3.1 (cierre con confirmación del usuario o a los siete días corridos). Evidencia: Inspección del cierre de tickets (RNF-TRV-21.15). \\
RT-21.18 Obligatorio & El Adjudicatario reportará mensualmente el cumplimiento de los niveles de servicio de atención, con el detalle por severidad y el análisis de los incumplimientos. & Cumple & Informe mensual de nivel de servicio dentro de los primeros cinco días hábiles & Subdoc. 3, 3.2.2.5; Formulario T-14, paquete 1.8.7.1. Evidencia: Informe mensual con cumplimiento por severidad y plan de acción \\
RT-21.19 Obligatorio & El Proponente declarará el tamaño de la bolsa anual de horas de mantención evolutiva, su composición por perfil y el procedimiento de solicitud, estimación, aprobación y liquidación. & Cumple parcialmente & Bolsa anual de horas de mantención evolutiva priorizada por el Cliente, con plan anual aprobado y flujo Kanban & Subdoc. 3, 3.2.2.6; Formulario T-14, paquete 1.8.5.1 (bolsa de 960 HH al año y liquidación mensual de consumo y saldo). Falta: composición de la bolsa por perfil. Evidencia: Declaración de la bolsa de horas (960 HH al año). \\
RT-21.20 Obligatorio & La mantención evolutiva se someterá al mismo estándar de calidad, pruebas y seguridad que el desarrollo original. & Cumple & Mantención evolutiva sujeta a cinco puertas de calidad, pruebas unitarias, de integración y de regresión automatizadas, y análisis estático y dinámico de seguridad & Formulario T-14, paquetes 1.1.4.2, 1.5.1.1, 1.5.1.2, 1.5.1.3, 1.8.5.1 y 1.8.8.4. Evidencia: Puertas de calidad, pruebas y análisis de seguridad aplicados a cada cambio de la mantención evolutiva. \\
RT-21.21 Obligatorio & El Proponente mantendrá un registro de deuda técnica cuantificado y destinará una fracción declarada de la capacidad de la Operación a reducirla. & Cumple parcialmente & Registro cuantificado de deuda técnica, con umbral de 5 \% o menos en el código nuevo y reserva de 10 \% del esfuerzo evolutivo mensual (8 HH al mes) & Subdoc. 6, 6.2.1.3; Subdoc. 9, 9.1.4.4; Formulario T-14, paquete 1.8.5.2 (8 HH al mes para reducir la deuda, 9 \% del esfuerzo evolutivo). Falta: fracción de la capacidad total de la Operación (meses 21 a 56) destinada a reducir la deuda. Evidencia: Registro de deuda técnica cuantificado y de su reducción en cada versión. \\
RT-21.22 Obligatorio & Las actualizaciones de versión de los componentes de base se planificarán anualmente, con ventana acordada y plan de reversión. & Cumple & Plan anual de actualización de los componentes de base, con ventana de 03:00 a 05:00 acordada, aviso de diez días hábiles, prueba de carga y reversión & Subdoc. 3, 3.2.2.6; Formulario T-14, paquete 1.8.4.1 (plan anual con ventana de 03:00 a 05:00, aviso de diez días hábiles y reversión). Evidencia: Plan anual de actualización con ventana y reversión (1.8.4.1). \\
\end{tabla}
\begin{tabla}{I{0.139}I{0.236}I{0.112}I{0.215}I{0.299}}{Capítulo 22. Capacitación y transferencia de conocimiento.}{tab:capitulo-22-capacitacion-transferencia-conocimiento}
  {\cab{ID requerimiento} & \cab{Descripción} & \cab{Estado} & \cab{Componente que lo satisface} & \cab{Sección de la propuesta y evidencia}}
  {\emph{Fuente:} elaboración propia a partir de las Bases técnicas transversales.}
RT-22.01 Obligatorio & El plan de capacitación se estructurará por perfil: personas usuarias finales, usuarias avanzadas, administradoras, equipo técnico y soporte de niveles 1 y 2. & Cumple parcialmente & Plan de capacitación por ola, perfil y turno, con dos ciclos del 7x7 antes del inicio de cada ola & Formulario T-14, paquetes 1.6.5.1 y 1.6.5.2 (cinco perfiles, con manuales y cursos por perfil). Falta: duración de la capacitación desglosada por perfil. Evidencia: Plan de capacitación por perfil (TRA-009). \\
RT-22.02 Obligatorio & Las modalidades incluirán capacitación presencial en cada sitio de operación, sesiones en línea sincrónicas, autoformación y acompañamiento en puesto de trabajo durante la marcha blanca. & Cumple & Plan de capacitación con las cuatro modalidades: presencial en cada sitio, en línea sincrónica, autoformación y acompañamiento en puesto & Formulario T-14, paquetes 1.6.5.1 y 1.6.3.3. Evidencia: Plan de capacitación con las cuatro modalidades (TRA-009). \\
RT-22.03 Obligatorio & Todo el material de capacitación se entregará en español, en formato editable y de propiedad del Cliente: manuales por perfil, guías rápidas, preguntas frecuentes, videos tutoriales y base de conocimiento consultable. & Cumple & Manuales por perfil, guías rápidas, preguntas frecuentes, videos tutoriales y base de conocimiento, en español, editables y de propiedad del cliente, publicados antes de la capacitación de cada perfil & Subdoc. 2 Anexos, TRA-009; Formulario T-14, paquete 1.6.5.2. Evidencia: Material de capacitación en español, editable y de propiedad del Cliente (TRA-009) \\
RT-22.04 Según caso & La capacitación no podrá afectar la operación del Cliente: se programará por turnos y en horarios acordados, considerando la estacionalidad y la ventana operacional del caso. & Cumple & Capacitación por turnos (día y noche, al inicio del turno) de al menos dos ciclos del 7x7, con registro por área y turno & Subdoc. 3, 3.1.3; Subdoc. 6, 6.1.5.3; Formulario T-14, paquete 1.6.5.1 (calendario acordado con la Contraparte Técnica del Cliente, fuera de marzo y septiembre; turnos de día y de noche). Evidencia: Plan y registro de capacitación por área y turno (TRA-003, TRA-007). \\
RT-22.05 Obligatorio & Las personas usuarias administradoras y el equipo técnico del Cliente serán evaluados y certificados como condición para el cierre de cada marcha blanca. & Cumple & Evaluación y certificación de administradoras y equipo técnico del Cliente como condición de cierre de cada marcha blanca & Subdoc. 3, 3.2.5; Subdoc. 2 Anexos, TRA-009; Formulario T-14, paquetes 1.6.3.1, 1.6.5.1 y 1.6.5.3. Evidencia: Actas de evaluación y certificación \\
RT-22.06 Obligatorio & Durante la Operación se ejecutarán al menos dos jornadas anuales de actualización y se capacitará al personal nuevo del Cliente sin costo adicional. & Cumple & Seis jornadas semestrales de actualización (dos por año) por perfil y a ambos grupos del 7x7, e inducción del personal nuevo antes de que opere el sistema & Formulario T-14, paquetes 1.8.9.1 (seis jornadas de actualización, dos por año, de M23 a M53) y 1.8.9.2 (inducción del personal nuevo del Cliente antes de operar el sistema, sin cargo adicional); Subdoc. 2, Anexos, RP-010. Evidencia: Registro de jornadas anuales de actualización (RP-010) \\
RT-22.07 Obligatorio & El Adjudicatario proveerá acompañamiento y mentoría al equipo técnico del Cliente durante al menos seis meses posteriores al paso a producción de la Etapa 2. & Cumple & Mentoría de seis meses al equipo de TI del Cliente tras el paso a producción de la Etapa 2, con acta firmada y ejercicio práctico por competencia & Subdoc. 13, 13.4.3; Subdoc. 2 Anexos, RP-010; Formulario T-14, paquete 1.6.5.3 (meses 22 a 27, seis meses posteriores al paso a producción del paquete 1.6.4.2). Evidencia: Plan de mentoría de seis meses tras la Etapa 2 (RP-010). \\
RT-22.08 Obligatorio & La base de conocimiento será mantenida y actualizada durante todo el Contrato, con métricas de uso y de utilidad percibida. & Cumple & Base de conocimiento editable en español, creada antes de la capacitación de cada perfil & Subdoc. 2 Anexos, RP-010; Formulario T-14, paquetes 1.6.5.2 (creación, meses 10 a 18) y 1.8.1.1 (mantención durante todo el Contrato, con métricas de uso y de utilidad). Evidencia: Métricas de uso y de utilidad de la base de conocimiento. \\
RT-22.09 Deseable & Se valorará la existencia de un ambiente permanente de entrenamiento, con datos ficticios, disponible para la práctica sin riesgo. & Cumple & Instancia de capacitación con datos sintéticos en el ambiente de QA & Formulario T-14, paquete 1.6.1.1 (instancia de capacitación con datos sintéticos alojada en el ambiente de QA). Evidencia: Declaración de la instancia de capacitación en QA con datos sintéticos (Formulario T-14, paquete 1.6.1.1) \\
\end{tabla}
\begin{tabla}{I{0.134}I{0.192}I{0.115}I{0.353}I{0.206}}{Capítulo 23. Información corporativa y presencia digital.}{tab:capitulo-23-informacion-corporativa-presencia-digital}
  {\cab{ID requerimiento} & \cab{Descripción} & \cab{Estado} & \cab{Componente que lo satisface} & \cab{Sección de la propuesta y evidencia}}
  {\emph{Fuente:} elaboración propia a partir de las Bases técnicas transversales.}
RT-23.01 Obligatorio & Información institucional: descripción detallada del giro principal, historia y trayectoria de la empresa con al menos tres años de antigüedad comprobable, misión, visión y valores, certificaciones y acreditaciones vigentes, y presencia geográfica y oficinas. & Cumple parcialmente & Sitio web corporativo (https://lucasdzamora.github.io/FEP\_Proyecto/) & Subdoc. 1, 1.1 (Tabla \refext{S1}{tab:datos-institucionales-overfore-s}) y 1.4 (Tabla \refext{S1}{tab:certificaciones-institucionales}); sitio web (Sobre N° 1). Falta: misión, visión y valores, e historia detallada. Evidencia: Sitio web activo con información institucional, certificaciones (Subdoc. 1, Tabla \refext{S1}{tab:certificaciones-institucionales}) y presencia geográfica (Subdoc. 1, Tabla \refext{S1}{tab:datos-institucionales-overfore-s}). \\
RT-23.02 Obligatorio & Experiencia y casos de éxito: portafolio de proyectos similares de los últimos cinco años, casos documentados con métricas verificables, testimonios de clientes con autorización de publicación, industrias atendidas con énfasis en la del caso asignado, y volumen de transacciones o de personas usuarias gestionadas. & Cumple parcialmente & Sitio web corporativo (https://lucasdzamora.github.io/FEP\_Proyecto/) & Subdoc. 1, 1.4 (Tabla \refext{S1}{tab:proyectos-referencia}); Subdoc. 1, Anexos A a D del Subdocumento 1; sitio web (Sobre N° 1). Falta: testimonios de clientes con autorización de publicación. Evidencia: Subdoc. 1, Tabla \refext{S1}{tab:proyectos-referencia}, con cuatro proyectos de referencia, sus métricas de servicio y su volumen; sitio web con portafolio de proyectos. \\
RT-23.03 Obligatorio & Equipo profesional: organigrama del equipo directivo, perfiles de socios y directores, currículos resumidos del equipo técnico clave y certificaciones profesionales del personal. & Cumple parcialmente & Sitio web corporativo (https://lucasdzamora.github.io/FEP\_Proyecto/) & Subdoc. 1, 1.5 (Tabla \refext{S1}{tab:niveles-dotacion-proyecto}); Subdoc. 1, Anexo E del Subdocumento 1 (Tablas \refext{A1}{tab:roles-dedicacion-periodo-responsabilidad-equipo-clave} y \refext{A1}{tab:certificaciones-competencias-equipo-clave} de los Anexos del Subdocumento 1); sitio web (Sobre N° 1). Falta: currículos resumidos del equipo técnico clave. Evidencia: Sitio web con equipo nominado y certificaciones de cada integrante (Tabla \refext{A1}{tab:certificaciones-competencias-equipo-clave} de los Anexos del Subdocumento 1). \\
RT-23.04 Obligatorio & Capacidades técnicas: servicios y soluciones ofrecidas, conjunto tecnológico dominado, alianzas tecnológicas con proveedores de nube y de plataforma, metodologías de trabajo certificadas, e infraestructura y capacidad instalada. & Cumple parcialmente & Sitio web corporativo (https://lucasdzamora.github.io/FEP\_Proyecto/) & Subdoc. 1, 1.1 (Tabla \refext{S1}{tab:capacidades-instaladas-respaldo}), 1.4 (Tabla \refext{S1}{tab:certificaciones-institucionales}) y 1.6 (Tabla \refext{S1}{tab:alianzas-tecnologicas-vigentes}); sitio web (Sobre N° 1). Falta: metodologías de trabajo certificadas; la Tabla \refext{S1}{tab:certificaciones-institucionales} declara certificaciones de sistemas de gestión. Evidencia: Sitio web con capacidades instaladas, certificaciones y alianzas con Azure e Integra OT. \\
RT-23.05 Obligatorio & El sitio será accesible conforme a WCAG 2.2 nivel AA, responsivo y con certificado TLS válido y vigente. & Se acredita en la instancia que fija el T-22 & Sitio web corporativo & Sitio web corporativo (Sobre N° 1). Falta: sitio accesible conforme a WCAG 2.2 nivel AA, responsivo y con certificado TLS válido y vigente. Evidencia: Informe WCAG 2.2 AA y certificado TLS vigente del sitio. \\
RT-23.06 Obligatorio & El sitio se mantendrá activo y disponible durante todo el proceso de licitación, desde el registro del participante hasta la adjudicación. & Se acredita en la instancia que fija el T-22 & Sitio web corporativo & Sitio web corporativo (Sobre N° 1). Falta: compromiso de mantener el sitio activo y disponible desde el registro del participante hasta la adjudicación, en la declaración jurada del Cap. 23.3. Evidencia: Monitoreo de disponibilidad del sitio hasta la adjudicación. \\
RT-23.07 Deseable & Material educativo digital: seminarios en línea, artículos técnicos o centro de recursos y documentación. & No aplica & Sitio web corporativo (Deseable no comprometido) & Sitio web; el sitio revisado no publica este contenido Evidencia: Deseable; su ausencia no penaliza \\
RT-23.08 Deseable & Demostración en línea, recorrido virtual de la solución o portal de soporte para clientes. & No aplica & Sitio web corporativo (Deseable no comprometido) & Sitio web; el sitio revisado no publica este contenido Evidencia: Deseable; su ausencia no penaliza \\
RT-23.09 Deseable & Métricas de disponibilidad y desempeño de los servicios del proponente publicadas en tiempo real. & No aplica & Sitio web corporativo (Deseable no comprometido) & Sitio web; el sitio revisado no publica este contenido Evidencia: Deseable; su ausencia no penaliza \\
RT-23.10 Deseable & Calculadora de retorno de la inversión o herramienta de estimación pertinente a la industria del caso. & No aplica & Sitio web corporativo (Deseable no comprometido) & Sitio web; el sitio revisado no publica este contenido Evidencia: Deseable; su ausencia no penaliza \\
\end{tabla}
\begin{tabla}{I{0.132}I{0.278}I{0.101}I{0.179}I{0.310}}{Capítulo 24. Video de presentación de la propuesta.}{tab:capitulo-24-video-presentacion-propuesta}
  {\cab{ID requerimiento} & \cab{Descripción} & \cab{Estado} & \cab{Componente que lo satisface} & \cab{Sección de la propuesta y evidencia}}
  {\emph{Fuente:} elaboración propia a partir de las Bases técnicas transversales.}
RT-24.01 Obligatorio & Todos los integrantes clave del equipo deberán aparecer individualmente presentándose, en toma de medio cuerpo, mirando directamente a la cámara, con vestimenta formal o de negocio informal, y con una duración mínima de aparición de 10 segundos por persona. & Se acredita en la instancia que fija el T-22 & Video de presentación de la propuesta (Cap. 24), que se entrega con la propuesta final & Video de presentación de la propuesta (entrega con la propuesta final). Falta: presentación individual de cada integrante clave en toma de medio cuerpo, mirando a cámara, con vestimenta formal o de negocio informal y al menos 10 segundos por persona. Evidencia: Revisión del video contra las especificaciones del Cap. 24. \\
RT-24.02 Obligatorio & Cada vez que aparezca un integrante deberá mostrarse en pantalla su nombre completo, su cargo en el proyecto, sus años de experiencia y su especialización relevante. & Se acredita en la instancia que fija el T-22 & Video de presentación de la propuesta (Cap. 24), que se entrega con la propuesta final & Video de presentación de la propuesta (entrega con la propuesta final). Falta: rótulo en pantalla con nombre completo, cargo en el proyecto, años de experiencia y especialización relevante de cada integrante, cada vez que aparece. Evidencia: Revisión del video contra las especificaciones del Cap. 24. \\
RT-24.03 Deseable & Las certificaciones principales de cada integrante podrán mostrarse en pantalla. & Se acredita en la instancia que fija el T-22 & Video de presentación de la propuesta (Cap. 24), que se entrega con la propuesta final & Video de presentación de la propuesta (entrega con la propuesta final). Falta: certificaciones principales de cada integrante en pantalla. Evidencia: Revisión del video contra las especificaciones del Cap. 24. \\
RT-24.04 Obligatorio & El fondo será el mismo para todas las tomas de personas, con iluminación consistente y encuadre similar para todos los participantes. & Se acredita en la instancia que fija el T-22 & Video de presentación de la propuesta (Cap. 24), que se entrega con la propuesta final & Video de presentación de la propuesta (entrega con la propuesta final). Falta: fondo común, iluminación consistente y encuadre similar para todos los participantes. Evidencia: Revisión del video contra las especificaciones del Cap. 24. \\
RT-24.05 Obligatorio & La paleta de colores corporativa, la tipografía y el uso de la marca serán consistentes en todo el video, con logotipo visible en todas las escenas y datos de contacto en el encabezado o pie. & Se acredita en la instancia que fija el T-22 & Video de presentación de la propuesta (Cap. 24), que se entrega con la propuesta final & Video de presentación de la propuesta (entrega con la propuesta final). Falta: paleta corporativa, tipografía y uso de la marca consistentes en todo el video, con logotipo visible en todas las escenas y datos de contacto en el encabezado o pie. Evidencia: Revisión del video final contra RT-24.05. \\
RT-24.06 Obligatorio & Los niveles de volumen estarán normalizados, con la misma calidad de grabación y sin variaciones bruscas de sonido. & Se acredita en la instancia que fija el T-22 & Video de presentación de la propuesta (Cap. 24), que se entrega con la propuesta final & Video de presentación de la propuesta (entrega con la propuesta final). Falta: niveles de volumen normalizados, con la misma calidad de grabación y sin variaciones bruscas de sonido. Evidencia: Revisión del video final contra RT-24.06. \\
\end{tabla}
\begin{tabla}{I{0.140}I{0.249}I{0.123}I{0.190}I{0.299}}{Capítulo 25. Prototipo interactivo y diseño UX/UI.}{tab:capitulo-25-prototipo-interactivo-diseno-ux}
  {\cab{ID requerimiento} & \cab{Descripción} & \cab{Estado} & \cab{Componente que lo satisface} & \cab{Sección de la propuesta y evidencia}}
  {\emph{Fuente:} elaboración propia a partir de las Bases técnicas transversales.}
RT-25.01 Obligatorio & Navegación intuitiva sin necesidad de manual, con un máximo de tres interacciones para alcanzar cualquier funcionalidad principal. & Cumple parcialmente & Prototipos navegables de los seis canales, con tres interacciones como máximo en las tareas críticas de terreno y validación con al menos cinco usuarios por perfil & Subdoc. 4, 4.1.4.1; Formulario T-14, paquetes 1.3.4.3. Falta: navegación sin necesidad de manual, como criterio de aceptación del prototipo. Evidencia: Revisión del prototipo según Cap. 25 y validación con cinco usuarios de cada perfil. \\
RT-25.02 Obligatorio & Retroalimentación visual clara ante cada acción, estados de carga y transición fluidos, y manejo elegante de los errores. & Cumple parcialmente & Sistema de diseño documentado con mensajes de error que dicen qué ocurrió y qué hacer, y respuesta visual a cada acción & Subdoc. 4, 4.1.4.1; Formulario T-14, paquetes 1.3.4.2 y 1.3.4.3. Falta: estados de carga y transiciones fluidas. Evidencia: Revisión del prototipo según Cap. 25. \\
RT-25.03 Obligatorio & Cumplimiento de WCAG 2.2 nivel AA: contraste adecuado, tamaños de fuente legibles, objetivos táctiles de al menos 44 × 44 píxeles, navegación completa por teclado, y textos alternativos y etiquetas de accesibilidad. & Cumple & Interfaces conformes con WCAG 2.2 nivel AA, verificadas con herramientas automatizadas y pruebas manuales de teclado, sobre el sistema de diseño documentado & Subdoc. 9, 9.2.4.3; Formulario T-14, paquete 1.3.4.3. Evidencia: Revisión del prototipo según Cap. 25, con herramienta automatizada y pruebas manuales de teclado. \\
RT-25.04 Obligatorio & Diseño responsivo demostrado en escritorio (1920 × 1080 y 1366 × 768), tableta (768 × 1024 en vertical y horizontal) y teléfono (375 × 812), con puntos de quiebre coherentes. & Cumple parcialmente & Diseño responsivo comprometido en escritorio, tableta y teléfono, a demostrar en el prototipo interactivo UX/UI & Subdoc. 2, Anexos, RNF-USA-011; Formulario T-14, paquete 1.4.5.2 (diseño responsivo probado en 1920×1080, 1366×768, tableta de 768×1024 en vertical y horizontal y teléfono de 375×812). Falta: declarar los puntos de quiebre del diseño responsivo y demostrarlo en el prototipo UX/UI (Cap. 25) en las cuatro resoluciones exigidas. Evidencia: Revisión del prototipo según Cap. 25 \\
RT-25.05 Obligatorio & Sistema de diseño documentado: paleta de a lo más cinco colores principales, a lo más dos familias tipográficas jerarquizadas, iconografía coherente, retícula y espaciado consistentes, y componentes reutilizables. & Cumple parcialmente & Sistema de diseño documentado con paleta de colores acotada, no más de dos familias tipográficas, iconografía, retícula y componentes reutilizables & Formulario T-14, paquete 1.3.4.2. Falta: paleta de a lo más cinco colores principales, jerarquía tipográfica de las dos familias y espaciado consistente. Evidencia: Sistema de diseño documentado y biblioteca publicada y usada por los prototipos. \\
RT-25.06 Obligatorio & El prototipo se entregará como enlace navegable en línea, sin instalación ni complementos, compatible con los navegadores modernos vigentes. & Cumple parcialmente & Prototipos navegables de los seis canales, revisados en los navegadores de la matriz de versiones soportadas & Formulario T-14, paquetes 1.3.4.2 y 1.3.4.3. Falta: entrega como enlace navegable en línea, sin instalación ni complementos. Evidencia: Revisión del prototipo según Cap. 25 en los navegadores de la matriz. \\
RT-25.07 Obligatorio & El acceso al prototipo estará garantizado por un mínimo de seis meses desde su entrega. & Se acredita en la instancia que fija el T-22 & Prototipo interactivo UX/UI (Cap. 25), entregado con el Informe 3 & Prototipo interactivo UX/UI (Cap. 25; entrega con el Informe 3). Falta: acceso al prototipo garantizado por un mínimo de seis meses desde su entrega. Evidencia: Revisión del prototipo según Cap. 25. \\
RT-25.08 Obligatorio & El prototipo permitirá navegación completa entre pantallas, simulación de ingreso de datos, transiciones y animaciones básicas, estados de posado y activo, simulación de carga de archivos, elementos desplegables funcionales y simulación de validaciones. & Cumple parcialmente & Prototipos navegables de los seis canales, validados con al menos cinco usuarios por perfil & Formulario T-14, paquete 1.3.4.3. Falta: simulación de ingreso de datos, transiciones y animaciones básicas, estados de posado y activo, carga de archivos, elementos desplegables y validaciones. Evidencia: Revisión del prototipo según Cap. 25. \\
RT-25.09 Obligatorio & Se entregará la arquitectura de información del prototipo: mapa del sitio completo, diagrama de navegación, taxonomía y nomenclatura, estructura de menús y jerarquía de la información. & Se acredita en la instancia que fija el T-22 & Prototipos navegables de los seis canales dentro del prototipo interactivo UX/UI (Cap. 25), entregado con el Informe 3 & Prototipo interactivo UX/UI (Cap. 25; entrega con el Informe 3); Formulario T-14, paquete 1.3.4.3. Falta: arquitectura de información del prototipo: mapa del sitio completo, diagrama de navegación, taxonomía y nomenclatura, estructura de menús y jerarquía de la información. Evidencia: Revisión del prototipo según Cap. 25. \\
RT-25.10 Deseable & Se valorará la posibilidad de exportar el prototipo a un documento interactivo para revisión sin conexión. & Se acredita en la instancia que fija el T-22 & Prototipo interactivo UX/UI (Cap. 25), entregado con el Informe 3 & Prototipo interactivo UX/UI (Cap. 25; entrega con el Informe 3); Formulario T-14, paquete 1.3.4.3. Falta: exportación del prototipo a un documento interactivo para revisión sin conexión. Evidencia: Revisión del prototipo según Cap. 25. \\
\end{tabla}
\begin{tabla}{I{0.143}I{0.243}I{0.076}I{0.236}I{0.302}}{Capítulo 26. Innovaciones.}{tab:capitulo-26-innovaciones}
  {\cab{ID requerimiento} & \cab{Descripción} & \cab{Estado} & \cab{Componente que lo satisface} & \cab{Sección de la propuesta y evidencia}}
  {\emph{Fuente:} elaboración propia a partir de las Bases técnicas transversales.}
RT-26.01 Obligatorio & Cada innovación se ubicará explícitamente en la arquitectura: qué capa la contiene, qué componentes la implementan y qué interfaces consume o expone. & Cumple & Ficha de cada innovación (Formulario T-19) con su capa, sus componentes e interfaces, dentro de la vista lógica de capas y componentes & Subdoc. 13, 13.1.3, 13.2.3, 13.3.3, 13.4.3 y 13.5.3; Subdoc. 4, 4.1.4.4 y 4.1.6.5 (capa e interfaces de los módulos de la innovación 4); Formulario T-14, paquete 1.3.1.1. Evidencia: Ficha de cada innovación (Formulario T-19). \\
RT-26.02 Obligatorio & Cada innovación identificará los paquetes de la estructura de descomposición del trabajo que la ejecutan y el mes del cronograma en que se materializa. & Cumple & Cartera de cinco innovaciones, cada una con sus paquetes de la EDT (rama 1.7) y su mes de materialización & Subdoc. 13, 13.1 a 13.5 y Tabla \refext{S13}{tab:cartera-innovaciones}; Formulario T-19; Formulario T-14, paquetes 1.7.1.1, 1.7.2.1, 1.7.3.2, 1.7.4.1 y 1.7.5.1. Evidencia: Ficha de cada innovación (Formulario T-19) \\
RT-26.03 Obligatorio & Las innovaciones de base tecnológica declararán el nivel de madurez de la tecnología con la escala utilizada y citarán las fuentes en norma APA 7.ª edición. & Cumple & Declaración del nivel TRL de cada innovación de base tecnológica, en la escala de ISO 16290, con fuentes en APA 7 en Referencias & Subdoc. 13, 13.1.2, 13.2.2, 13.3.2 y 13.5.2 (nivel TRL en la escala de ISO 16290, con fuentes en APA 7 en Referencias); Subdoc. 13, 13.4.2 (práctica consolidada, fuera de la escala TRL); Formulario T-19. Evidencia: Ficha de cada innovación (Formulario T-19). \\
RT-26.04 Obligatorio & Cada innovación declarará su riesgo de adopción, su probabilidad, su impacto, la estrategia de mitigación y el plan de contingencia si no rinde lo esperado. & Cumple & Riesgo de adopción de cada innovación con probabilidad, impacto, mitigación y plan de contingencia & Subdoc. 13, 13.1.6, 13.2.6, 13.3.6, 13.4.6 y 13.5.6; Formulario T-19; Formulario T-14, paquete 1.7.3.1. Evidencia: Ficha de cada innovación (Formulario T-19) \\
RT-26.05 Obligatorio & Cada innovación declarará su indicador de verificación con línea base, meta y momento de medición, y su impacto en inversión, costo operacional y beneficio esperado. & Cumple & Indicador de verificación de cada innovación con línea base, meta y momento de medición, e impacto en inversión, costo operacional y beneficio & Subdoc. 13, 13.1.4-13.1.5, 13.2.4-13.2.5, 13.3.4-13.3.5, 13.4.4-13.4.5 y 13.5.4-13.5.5; Formulario T-19; Formulario T-14, paquetes 1.7.1.2, 1.7.2.2, 1.7.3.3, 1.7.4.1 y 1.7.5.2. Evidencia: Ficha de cada innovación (Formulario T-19) \\
RT-26.06 Obligatorio & Las innovaciones que incorporen inteligencia artificial cumplirán íntegramente el Capítulo 18 de este documento. & Cumple & Clasificador de detenciones (única innovación con IA) con gobierno conforme al Cap. 18: modelo y ubicación, datos del cliente, validación humana, riesgos, registro, desactivación, transparencia y reentrenamiento & Subdoc. 13, 13.2.2; Subdoc. 13 Anexos, B.5 (Tabla B.4); Formulario T-19; Formulario T-14, paquetes 1.7.2.1, 1.7.2.3 y 1.7.2.4. Evidencia: Ficha de cada innovación (Formulario T-19) \\
RT-26.07 Obligatorio & Las innovaciones que modifiquen la arquitectura de seguridad requerirán su propio modelado de amenazas. & Cumple & Modelado de amenazas STRIDE propio del bloque 3 del traspaso (administración de usuarios y roles internos) & Subdoc. 13, 13.4.3; Subdoc. 13 Anexos, D.5 (Tabla D.2); Formulario T-19; Formulario T-14, paquetes 1.3.3.2 y 1.7.4.1. Evidencia: Ficha de cada innovación (Formulario T-19) \\
RT-26.08 Deseable & Se valorará que al menos una innovación sea verificable durante la marcha blanca de la Etapa 1, es decir, que su beneficio pueda medirse antes del mes 16. & Cumple & Innovaciones 1 y 3 con indicadores medidos en la marcha blanca de la Etapa 1 (meses 13 a 15) & Subdoc. 13, 13.1.5 y 13.3.5 (Tabla \refext{S13}{tab:indicadores-factores-conciliacion-estimados-incertidumbre} para la innovación 3); Formulario T-19; Formulario T-14, paquetes 1.7.1.2 y 1.7.3.3. Evidencia: Ficha de cada innovación (Formulario T-19) \\
\end{tabla}

\seccion{Requerimientos propios del caso}
La columna «Descripción» resume cada requerimiento; su texto completo, con el mismo identificador, está en el Anexo A de OVERFORE-Subdocumento2-Anexos.

\begin{tabla}{I{0.128}I{0.343}I{0.108}I{0.159}I{0.262}}{Requerimientos funcionales propios del caso.}{tab:requerimientos-funcionales-propios-caso}
  {\cab{ID requerimiento} & \cab{Descripción} & \cab{Estado} & \cab{Componente que lo satisface} & \cab{Sección de la propuesta y evidencia}}
  {\emph{Fuente:} elaboración propia a partir de las Bases técnicas del caso Minería.}
RF-TRZ-001 Propio - Funcional & El sistema debe registrar cada ciclo de carguío con el polígono de extracción vigente, el viaje… & Cumple & Módulo Operación Mina: ciclo de carguío con polígono, viaje, pala, operador, inicio y fin, recibido del conector de despacho sin pedir datos al operador & Subdoc. 4, 4.1.6.1.1; Subdoc. 5, 5.1.6; Subdoc. 5 Anexos, A.3 (entidad Ciclo de carguío); Subdoc. 2 Anexos, A.1 (RF-TRZ-001); Formulario T-14, paquetes 1.4.2.1 y 1.4.3.2. Evidencia: Prueba de aceptación del resultado esperado (Anexo A.1) en la marcha blanca de la Etapa 1 \\
RF-TRZ-002 Propio - Funcional & El sistema debe mostrar, para cualquier turno, lote de chancado o lote comercial, los viajes… & Cumple & Módulo Trazabilidad y Embarque & Subdoc. 3, 3.4.2.2; Formulario T-14, paquetes 1.4.3.2 y 1.4.3.3. Evidencia: Prueba de aceptación del resultado esperado (Anexo A.1) en la marcha blanca de la Etapa 1 \\
RF-TRZ-003 Propio - Funcional & El sistema debe permitir reconstruir, con evidencia, el origen de cualquier lote embarcado: de qué polígonos… & Cumple & Módulo Trazabilidad y Embarque (Portal para Compradores en Etapa 2) & Subdoc. 3, 3.4.2.2; Formulario T-14, paquete 1.4.3.3. Evidencia: Prueba de aceptación del resultado esperado (Anexo A.1) en la marcha blanca de la Etapa 1 \\
RF-TRZ-004 Propio - Funcional & El sistema debe permitir hacer el camino inverso: a partir de un polígono extraído… & Cumple parcialmente & Módulo Trazabilidad y Embarque (Portal para Compradores en Etapa 2) & Subdoc. 3, 3.4.1 y 3.4.2.2; Subdoc. 5, 5.1.6 y 5.2.4; Formulario T-14, paquete 1.4.3.3 (camino inverso con proporción por lote y embarque). Falta: historial de diez años consultable desde la entrega; hoy se compromete cinco años migrados y diez sólo con la operación. Evidencia: Prueba de aceptación del resultado esperado (Anexo A.1) en la marcha blanca de la Etapa 1. \\
RF-TRZ-005 Propio - Funcional & El sistema debe indicar, en cada etapa de la cadena de origen de un lote… & Cumple parcialmente & Módulo Trazabilidad y Embarque (Portal para Compradores en Etapa 2) & Subdoc. 3, 3.4.2.2; Subdoc. 5, Anexo A.3, Tabla A.5 (origen medido o estimado y versión de la regla); Subdoc. 4, 4.1.6.1.3; Formulario T-14, paquete 1.4.3.3 (marca de medido o supuesto por dato). Falta: versión de la regla aplicada en cada supuesto, mostrada en pantalla. Evidencia: Prueba de aceptación del resultado esperado (Anexo A.1) en la marcha blanca de la Etapa 1. \\
RF-TRZ-006 Propio - Funcional & El sistema debe llevar el registro de cada stock intermedio: lo que entra… & Cumple & Módulo Trazabilidad y Embarque (Portal para Compradores en Etapa 2) & Subdoc. 3, 3.4.2.2; Formulario T-14, paquete 1.4.3.3. Evidencia: Prueba de aceptación del resultado esperado (Anexo A.1) en la marcha blanca de la Etapa 1 \\
RF-TRZ-007 Propio - Funcional & El sistema debe mantener actualizado, para cada stock intermedio y para el stock de baja ley… & Cumple & Módulo Trazabilidad y Embarque (Portal para Compradores en Etapa 2) & Subdoc. 3, 3.4.2.2; Subdoc. 5, 5.1.6 y 5.4.5; Subdoc. 5, Anexo A.3, Tabla A.5 (Stock); Formulario T-14, paquete 1.4.3.3. Evidencia: Prueba de aceptación del resultado esperado (Anexo A.1) en la marcha blanca de la Etapa 1. \\
RF-TRZ-008 Propio - Funcional & El sistema debe calcular la ley del material que se retoma de un stock… & Cumple & Módulo Trazabilidad y Embarque (Portal para Compradores en Etapa 2) & Subdoc. 3, 3.4.2.2; Formulario T-14, paquete 1.4.3.3. Evidencia: Prueba de aceptación del resultado esperado (Anexo A.1) en la marcha blanca de la Etapa 1 \\
RF-TRZ-009 Propio - Funcional & El sistema debe registrar, en cada campaña de verificación de un stock, la ley… & Cumple & Módulo Trazabilidad y Embarque (Portal para Compradores en Etapa 2) & Subdoc. 3, 3.4.2.2; Formulario T-14, paquete 1.4.3.3. Evidencia: Prueba de aceptación del resultado esperado (Anexo A.1) en la marcha blanca de la Etapa 1 \\
RF-TRZ-010 Propio - Funcional & El sistema debe asignar cada viaje de camión a un polígono de origen. & Cumple & Módulo Operación Mina (Botonera de Pala, App SITCO Terreno y conector de despacho de flota) & Subdoc. 3, 3.4.2.1 y 3.3.3; Formulario T-14, paquete 1.4.3.2. Evidencia: Prueba de aceptación del resultado esperado (Anexo A.1) en la marcha blanca de la Etapa 1 \\
RF-TRZ-011 Propio - Funcional & El sistema debe registrar, en cada ciclo de carguío, dónde está la pala respecto del contacto… & Cumple & Módulo Operación Mina: el ciclo de carguío registra la posición de la pala recibida del despacho y la marca de proximidad al borde, sin acción de la persona operadora & Subdoc. 4, 4.1.6.5, Tabla \refext{S4}{tab:llamadas-sincronas-entre-componentes-umbral}; Subdoc. 5, Anexo A.3, Tabla A.4 (Ciclo de carguío: posicion\_carguio y marca\_proximidad\_borde); Subdoc. 3, 3.4.2.1; Formulario T-14, paquete 1.4.3.2. Evidencia: Prueba de aceptación del resultado esperado (Anexo A.1) en la marcha blanca de la Etapa 1 \\
RF-TRZ-012 Propio - Funcional & El sistema debe preparar y entregar a SAP S/4HANA, para su emisión electrónica, los datos… & Cumple & Módulo Trazabilidad y Embarque (Portal para Compradores en Etapa 2) & Subdoc. 3, 3.4.2.2; Formulario T-14, paquete 1.4.3.3. Evidencia: Prueba de aceptación del resultado esperado (Anexo A.1) en la marcha blanca de la Etapa 1 \\
RF-TRZ-013 Propio - Funcional & El sistema debe registrar la recepción y el acopio del concentrado en el recinto portuario… & Cumple & Módulo Trazabilidad y Embarque (Portal para Compradores en Etapa 2) & Subdoc. 3, 3.4.2.2; Formulario T-14, paquete 1.4.3.3. Evidencia: Prueba de aceptación del resultado esperado (Anexo A.1) en la marcha blanca de la Etapa 1 \\
RF-TRZ-014 Propio - Funcional & El sistema debe registrar, con marca de tiempo y responsable, cada evento del embarque… & Cumple & Módulo Trazabilidad y Embarque: hito de embarque con tipo (programación a zarpe), fecha del hecho, fecha de registro y autor, con acceso del operador portuario solo a esos eventos & Subdoc. 5, Anexo A.3, Tabla A.5 (Hito de embarque); Subdoc. 5, 5.1.6; Subdoc. 4, 4.1.7.7 y 4.1.8.3; Formulario T-14, paquete 1.4.3.3. Evidencia: Prueba de aceptación del resultado esperado (Anexo A.1) en la marcha blanca de la Etapa 1 \\
RF-TRZ-015 Propio - Funcional & El sistema debe conformar el lote comercial a partir de los lotes de concentrado… & Cumple & Módulo Trazabilidad y Embarque (Portal para Compradores en Etapa 2) & Subdoc. 3, 3.4.2.2; Formulario T-14, paquete 1.4.3.3. Evidencia: Prueba de aceptación del resultado esperado (Anexo A.1) en la marcha blanca de la Etapa 1 \\
RF-TRZ-016 Propio - Funcional & El sistema debe poner a disposición del comprador, en un portal con acceso autenticado, el certificado… & Cumple & Módulo Trazabilidad y Embarque (Portal para Compradores en Etapa 2) & Subdoc. 3, 3.4.2.2; Formulario T-14, paquete 1.4.4.5. Evidencia: Prueba de aceptación del resultado esperado (Anexo A.1) en la marcha blanca de la Etapa 2 \\
RF-TRZ-017 Propio - Funcional & El sistema debe registrar el acta de aceptación del embarque con firma electrónica avanzada y vincularla… & Cumple & Módulo Trazabilidad y Embarque, con firma electrónica avanzada gestionada por Administración y Servicios Comunes & Subdoc. 4, 4.1.6.1.2 y 4.1.7.7; Subdoc. 3, 3.4.3; Formulario T-14, paquetes 1.4.3.1 y 1.4.3.3. Evidencia: Prueba de aceptación del resultado esperado (Anexo A.1) en la marcha blanca de la Etapa 1 \\
RF-TRZ-018 Propio - Funcional & Cuando el laboratorio rechace o invalide una muestra, el sistema debe bloquear el lote… & Cumple & Módulo Trazabilidad y Embarque (Portal para Compradores en Etapa 2) & Subdoc. 3, 3.4.2.2; Formulario T-14, paquetes 1.4.2.1 y 1.4.3.3. Evidencia: Prueba de aceptación del resultado esperado (Anexo A.1) en la marcha blanca de la Etapa 1 \\
RF-CON-001 Propio - Funcional & El sistema debe registrar y conservar las tres mediciones de tonelaje tal como se generan… & Cumple & Módulo Conciliación Metalúrgica & Subdoc. 3, 3.4.2.3; Formulario T-14, paquetes 1.4.2.1 y 1.4.3.4. Evidencia: Prueba de aceptación del resultado esperado (Anexo A.1) en la marcha blanca de la Etapa 1 \\
RF-CON-002 Propio - Funcional & El sistema debe usar en la conciliación metalúrgica una única fuente de verdad del tonelaje, mostrar… & Cumple & Módulo Conciliación Metalúrgica & Subdoc. 3, 3.4.2.3; Formulario T-14, paquetes 1.2.3.2 y 1.4.3.4. Evidencia: Prueba de aceptación del resultado esperado (Anexo A.1) en la marcha blanca de la Etapa 1 \\
RF-CON-003 Propio - Funcional & El sistema debe calcular la conciliación metalúrgica del mes dentro del sistema, con acceso al dato… & Cumple & Módulo Conciliación Metalúrgica & Subdoc. 3, 3.4.2.3; Formulario T-14, paquete 1.4.3.4. Evidencia: Prueba de aceptación del resultado esperado (Anexo A.1) en la marcha blanca de la Etapa 1 \\
RF-CON-004 Propio - Funcional & El sistema debe descomponer la diferencia de conciliación por punto de la cadena -mina, stocks intermedios… & Cumple & Módulo Conciliación Metalúrgica & Subdoc. 3, 3.4.2.3; Formulario T-14, paquete 1.4.3.4. Evidencia: Prueba de aceptación del resultado esperado (Anexo A.1) en la marcha blanca de la Etapa 1 \\
RF-CON-005 Propio - Funcional & El sistema debe alertar cuando la diferencia de conciliación del período supere el umbral comprometido, indicando… & Cumple parcialmente & Módulo Conciliación Metalúrgica & Subdoc. 3, 3.4.2.3; Subdoc. 5, Anexo A.3, Tabla A.6 (diferencia y punto atribuido); Formulario T-14, paquetes 1.4.3.4 (alerta al superar el umbral) y 1.8.7.1 (umbral de no más de 3 \% desde el mes 22 y 2 \% desde el mes 28). Falta: responsable del punto de la cadena que más aporta a la diferencia. Evidencia: Prueba de aceptación del resultado esperado (Anexo A.1) en la marcha blanca de la Etapa 1. \\
RF-DAT-001 Propio - Funcional & El sistema debe cuidar la calidad de los datos según la norma ISO/IEC 25012: validarlos… & Cumple & Módulo Administración y Servicios Comunes: validación en el punto de captura y tablero de calidad por dominio y mes (completitud, exactitud y consistencia) & Subdoc. 5, 5.2.3 (Tabla \refext{S5}{tab:indicadores-calidad-datos-metas}); módulo asignado en Subdoc. 4, 4.1.3.2, Tabla \refext{S4}{tab:ubicacion-cada-capa-faena-nodo-faena}; Formulario T-14, paquetes 1.4.3.1 y 1.4.4.6. Evidencia: Prueba de aceptación del resultado esperado (Anexo A.1) en la marcha blanca de la Etapa 1 \\
RF-CON-006 Propio - Funcional & El sistema debe publicar la diferencia entre las tres mediciones de tonelaje por período, por equipo… & Cumple & Módulo Conciliación Metalúrgica & Subdoc. 3, 3.4.2.3; Formulario T-14, paquete 1.4.3.4. Evidencia: Prueba de aceptación del resultado esperado (Anexo A.1) en la marcha blanca de la Etapa 1 \\
RF-CON-007 Propio - Funcional & El sistema debe guardar cada factor de ajuste de la conciliación como registro con versiones: fórmula… & Cumple & Módulo Conciliación Metalúrgica & Subdoc. 3, 3.4.2.3; Formulario T-14, paquetes 1.2.6.1 y 1.4.3.4. Evidencia: Prueba de aceptación del resultado esperado (Anexo A.1) en la marcha blanca de la Etapa 1 \\
RF-CON-008 Propio - Funcional & El sistema debe calcular el balance metalúrgico de cada turno con los mismos registros… & Cumple & Módulo Conciliación Metalúrgica & Subdoc. 3, 3.4.2.3; Formulario T-14, paquete 1.4.3.4. Evidencia: Prueba de aceptación del resultado esperado (Anexo A.1) en la marcha blanca de la Etapa 1 \\
RF-CON-009 Propio - Funcional & El sistema debe recibir los resultados de análisis directamente del sistema de gestión de laboratorio, conectado… & Cumple & Módulo Conciliación Metalúrgica & Subdoc. 3, 3.4.2.3; Formulario T-14, paquete 1.4.2.1. Evidencia: Prueba de aceptación del resultado esperado (Anexo A.1) en la marcha blanca de la Etapa 1 \\
RF-CON-010 Propio - Funcional & El sistema debe mostrar junto a cada resultado de conciliación su memoria de cálculo: el detalle… & Cumple parcialmente & Módulo Conciliación Metalúrgica & Subdoc. 5, 5.1.6 y Anexo A.3, Tabla A.6 (paso de cálculo); Subdoc. 3, 3.4.2.3; Formulario T-14, paquete 1.4.3.4 (memoria de cálculo al reproducir un período). Falta: detalle paso a paso con datos de entrada, factores e intermedios junto a cada resultado. Evidencia: Prueba de aceptación del resultado esperado (Anexo A.1) en la marcha blanca de la Etapa 1. \\
RF-CON-011 Propio - Funcional & El sistema debe capturar la muestra en los puntos de transferencia sin agregar pasos al ciclo… & Cumple parcialmente & Módulo Conciliación Metalúrgica & Subdoc. 3, 3.2.3.2 y 3.4.5; Formulario T-14, paquete 1.4.3.4; Subdoc. 2 Anexos, SUP-12 (cada punto agrega no más de 2 segundos en el percentil 95). Falta: toneladas menos al año por cada punto de muestreo. Evidencia: Prueba de aceptación del resultado esperado (Anexo A.1) en la marcha blanca de la Etapa 1. \\
RF-CON-012 Propio - Funcional & El sistema debe publicar por polígono la diferencia entre la ley estimada del plan… & Cumple & Módulo Conciliación Metalúrgica; tablero de calidad de datos del módulo Indicadores y Centro Integrado & Subdoc. 5, 5.2.3, Tabla \refext{S5}{tab:indicadores-calidad-datos-metas} (desviación entre ley estimada y ley medida por polígono); Subdoc. 3, 3.4.2.2; Subdoc. 5, Anexo A.3, Tabla A.5 (origen\_ley medida\_laboratorio); Formulario T-14, paquetes 1.4.2.1 y 1.4.3.4. Evidencia: Prueba de aceptación del resultado esperado (Anexo A.1) en la marcha blanca de la Etapa 1 \\
RF-CON-013 Propio - Funcional & El sistema debe calcular el costo por tonelada movida y por libra de cobre fino todos… & Cumple & Módulo Indicadores y Centro Integrado & Subdoc. 3, 3.4.3; Formulario T-14, paquetes 1.4.2.3 y 1.4.4.6. Evidencia: Prueba de aceptación del resultado esperado (Anexo A.1) en la marcha blanca de la Etapa 2 \\
RF-CON-014 Propio - Funcional & El sistema debe permitir ir desde el costo por tonelada hasta cada transacción de origen… & Cumple & Módulo Indicadores y Centro Integrado & Subdoc. 3, 3.4.3; Formulario T-14, paquete 1.4.4.6. Evidencia: Prueba de aceptación del resultado esperado (Anexo A.1) en la marcha blanca de la Etapa 2 \\
RF-CON-015 Propio - Funcional & El sistema debe entregar a SAP los datos operacionales que este necesita para su registro contable… & Cumple & Módulo Conciliación Metalúrgica & Subdoc. 3, 3.4.2.3; Formulario T-14, paquete 1.4.2.1. Evidencia: Prueba de aceptación del resultado esperado (Anexo A.1) en la marcha blanca de la Etapa 1 \\
RF-TER-001 Propio - Funcional & El sistema debe permitir a la persona operadora de la pala avisar, con una sola pulsación… & Cumple & Módulo Operación Mina (Botonera de Pala, App SITCO Terreno y conector de despacho de flota) & Subdoc. 3, 3.4.2.1 y 3.3.3; Formulario T-14, paquetes 1.4.3.2 y 1.4.5.1. Evidencia: Prueba de aceptación del resultado esperado (Anexo A.1) en la marcha blanca de la Etapa 1 \\
RF-TER-002 Propio - Funcional & El sistema debe indicar en la cabina de la pala, con una luz por cada destino… & Cumple & Módulo Operación Mina (Botonera de Pala, App SITCO Terreno y conector de despacho de flota) & Subdoc. 3, 3.4.2.1 y 3.3.3; Formulario T-14, paquetes 1.4.5.1 y 1.6.2.1. Evidencia: Prueba de aceptación del resultado esperado (Anexo A.1) en la marcha blanca de la Etapa 1 \\
RF-TER-003 Propio - Funcional & El sistema debe avisar a la persona despachadora del aviso de cambio de material mediante… & Cumple & Módulo Operación Mina (Botonera de Pala, App SITCO Terreno y conector de despacho de flota) & Subdoc. 3, 3.4.2.1 y 3.3.3; Formulario T-14, paquetes 1.4.2.1 y 1.4.3.2. Evidencia: Prueba de aceptación del resultado esperado (Anexo A.1) en la marcha blanca de la Etapa 1 \\
RF-TER-004 Propio - Funcional & El sistema debe registrar el cambio de destino que la persona despachadora ejecuta en el sistema… & Cumple & Módulo Operación Mina (Botonera de Pala, App SITCO Terreno y conector de despacho de flota) & Subdoc. 3, 3.4.2.1 y 3.3.3; Formulario T-14, paquetes 1.4.2.1 y 1.4.3.2. Evidencia: Prueba de aceptación del resultado esperado (Anexo A.1) en la marcha blanca de la Etapa 1 \\
RF-TER-005 Propio - Funcional & El sistema debe conservar el aviso de cambio de material cuando el equipo opere hasta 48… & Cumple & Módulo Operación Mina (Botonera de Pala, App SITCO Terreno y conector de despacho de flota) & Subdoc. 3, 3.3.6; Subdoc. 4, 4.1.5.1 y 4.1.9.2; Formulario T-14, paquetes 1.4.2.2 y 1.4.5.1. Evidencia: Prueba de aceptación del resultado esperado (Anexo A.1) en la marcha blanca de la Etapa 1 \\
RF-TER-006 Propio - Funcional & El sistema debe capturar los datos de cada pozo perforado: de forma automática en las perforadoras… & Cumple & Módulo Operación Mina (Botonera de Pala, App SITCO Terreno y conector de despacho de flota) & Subdoc. 3, 3.4.2.1 y 3.3.3; Formulario T-14, paquetes 1.4.3.2 y 1.4.5.1. Evidencia: Prueba de aceptación del resultado esperado (Anexo A.1) en la marcha blanca de la Etapa 1 \\
RF-TER-007 Propio - Funcional & El sistema debe entregar a geología los datos de perforación en formato estructurado y exportable… & Cumple parcialmente & Módulo Operación Mina (Botonera de Pala, App SITCO Terreno y conector de despacho de flota) & Subdoc. 4, 4.1.6.1.1; Subdoc. 5, Anexo A.3, Tabla A.4 (Pozo de perforación); Formulario T-14, paquete 1.4.3.2 (registro de cada pozo). Falta: entrega a geología en CSV con pozo, posición, profundidad y parámetros, y estándar abierto cuando el fabricante lo soporte. Evidencia: Prueba de aceptación del resultado esperado (Anexo A.1) en la marcha blanca de la Etapa 1. \\
RF-TER-008 Propio - Funcional & El sistema debe registrar cada tronadura con fecha, hora, sector, pozos y desplazamiento declarado del contacto… & Cumple & Módulo Operación Mina (Botonera de Pala, App SITCO Terreno y conector de despacho de flota) & Subdoc. 4, 4.1.6.1.1; Subdoc. 5, Anexo A.3, Tabla A.4 (Tronadura y Polígono afectado por tronadura); Formulario T-14, paquete 1.4.3.2. Evidencia: Prueba de aceptación del resultado esperado (Anexo A.1) en la marcha blanca de la Etapa 1 \\
RF-PLA-001 Propio - Funcional & El sistema debe consumir automáticamente las salidas del software de planificación minera -polígonos, leyes estimadas, destinos… & Cumple & Módulo Operación Mina (Botonera de Pala, App SITCO Terreno y conector de despacho de flota) & Subdoc. 3, 3.4.2.1 y 3.3.3; Formulario T-14, paquetes 1.4.2.1 y 1.4.3.2. Evidencia: Prueba de aceptación del resultado esperado (Anexo A.1) en la marcha blanca de la Etapa 1 \\
RF-PLA-002 Propio - Funcional & El sistema debe guardar el plan minero con versiones: para cada una, su fecha y hora… & Cumple & Módulo Operación Mina (Botonera de Pala, App SITCO Terreno y conector de despacho de flota) & Subdoc. 3, 3.4.2.1 y 3.3.3; Formulario T-14, paquete 1.4.3.2. Evidencia: Prueba de aceptación del resultado esperado (Anexo A.1) en la marcha blanca de la Etapa 1 \\
RF-PLA-003 Propio - Funcional & El sistema debe calcular el cumplimiento del plan minero comparando lo ejecutado contra la versión… & Cumple & Módulo Operación Mina (Botonera de Pala, App SITCO Terreno y conector de despacho de flota) & Subdoc. 3, 3.4.2.1 y 3.3.3; Formulario T-14, paquete 1.4.3.2. Evidencia: Prueba de aceptación del resultado esperado (Anexo A.1) en la marcha blanca de la Etapa 1 \\
RF-PLA-004 Propio - Funcional & El sistema debe registrar el motivo de cada cambio de plan eligiéndolo de una lista cerrada… & Cumple parcialmente & Módulo Operación Mina (Botonera de Pala, App SITCO Terreno y conector de despacho de flota) & Subdoc. 3, 3.4.2.1; Subdoc. 5, Anexo A.3, Tabla A.4 (versión de plan con motivo\_cambio de catálogo); Formulario T-14, paquetes 1.4.3.2 (motivo de lista cerrada y notificación en 15 minutos) y 1.4.3.1 (canales de notificación del proyecto). Falta: que el plazo de 15 minutos se cuente desde la entrada en vigencia del cambio. Evidencia: Prueba de aceptación del resultado esperado (Anexo A.1) en la marcha blanca de la Etapa 1. \\
RF-PLA-005 Propio - Funcional & El sistema debe mantener un catálogo cerrado y con versiones de los tipos de detención… & Cumple & Módulo Operación Mina (Botonera de Pala, App SITCO Terreno y conector de despacho de flota) & Subdoc. 3, 3.4.2.1 y 3.3.3; Formulario T-14, paquetes 1.2.2.5 y 1.4.3.2. Evidencia: Prueba de aceptación del resultado esperado (Anexo A.1) en la marcha blanca de la Etapa 1 \\
RF-PLA-006 Propio - Funcional & El sistema debe clasificar automáticamente las detenciones de equipos y de planta usando la telemetría… & Cumple & Módulo Equipos y Mantenimiento (vista única del equipo) & Subdoc. 4, 4.1.6.2; Subdoc. 3, 3.2.5, Tabla \refext{S3}{tab:criterios-aceptacion-etapa-1} y \refext{S3}{tab:sistemas-existentes-relacion-sitco-aranda}.2.4; Subdoc. 5, Anexo A.3, Tabla A.4 (Detención, taxonomía cerrada sin texto libre); Formulario T-14, paquetes 1.4.3.2 y 1.4.4.1. Evidencia: Prueba de aceptación del resultado esperado (Anexo A.1) en la marcha blanca de la Etapa 2 \\
RF-PLA-007 Propio - Funcional & El sistema debe registrar la causa raíz de cada detención no programada de equipo… & Cumple & Módulo Equipos y Mantenimiento (vista única del equipo) & Subdoc. 4, 4.1.6.2; Subdoc. 5, 5.1.6; Subdoc. 5, Anexo A.3, Tabla A.7 (responsable de diagnóstico); Formulario T-14, paquete 1.4.4.1. Evidencia: Prueba de aceptación del resultado esperado (Anexo A.1) en la marcha blanca de la Etapa 2 \\
RF-PLA-008 Propio - Funcional & El sistema debe calcular y publicar por equipo, área y turno la disponibilidad física, la utilización… & Cumple & Módulo Operación Mina (Botonera de Pala, App SITCO Terreno y conector de despacho de flota) & Subdoc. 5, 5.2.2 y Anexo A.3, Tabla A.17 (hecho Productividad por detención); Formulario T-14, paquetes 1.4.3.6 (indicadores con definición y fórmula publicadas) y 1.4.3.2 (tiempo sin clasificar bajo 5 \%). Evidencia: Prueba de aceptación del resultado esperado (Anexo A.1) en la marcha blanca de la Etapa 1. \\
RF-PLA-009 Propio - Funcional & El sistema debe registrar las horas de detención no programada de planta leyendo del historiador… & Cumple & Módulo Operación Mina (Botonera de Pala, App SITCO Terreno y conector de despacho de flota) & Subdoc. 3, 3.4.2.1 y 3.3.3; Formulario T-14, paquetes 1.4.2.1 y 1.4.4.1. Evidencia: Prueba de aceptación del resultado esperado (Anexo A.1) en la marcha blanca de la Etapa 1 \\
RF-ACR-001 Propio - Funcional & El sistema debe mantener un expediente digital por persona con sus contratos, exámenes, cursos, protocolos… & Cumple & Módulo Acreditación y Acceso (Portal de Contratistas y Consola de Portería) & Subdoc. 3, 3.4.2.5; Formulario T-14, paquete 1.4.3.5. Evidencia: Prueba de aceptación del resultado esperado (Anexo A.1) en la marcha blanca de la Etapa 1 \\
RF-ACR-002 Propio - Funcional & El sistema debe ejecutar el flujo de acreditación de una persona contratista con etapas, responsables… & Cumple & Módulo Acreditación y Acceso (Portal de Contratistas y Consola de Portería) & Subdoc. 3, 3.4.2.5; Formulario T-14, paquete 1.4.3.5. Evidencia: Prueba de aceptación del resultado esperado (Anexo A.1) en la marcha blanca de la Etapa 1 \\
RF-ACR-003 Propio - Funcional & El sistema debe validar automáticamente la completitud y la vigencia de la documentación cargada e indicar… & Cumple & Módulo Acreditación y Acceso (Portal de Contratistas y Consola de Portería) & Subdoc. 3, 3.4.2.5; Formulario T-14, paquete 1.4.3.5. Evidencia: Prueba de aceptación del resultado esperado (Anexo A.1) en la marcha blanca de la Etapa 1 \\
RF-ACR-004 Propio - Funcional & El sistema debe permitir consultar en cualquier momento el estado de habilitación de cada persona presente… & Cumple & Módulo Acreditación y Acceso (Portal de Contratistas y Consola de Portería): consulta de habilitación, expediente con cada versión documental y tabla de lectura de personas presentes & Subdoc. 4, 4.1.6.5, Tabla \refext{S4}{tab:llamadas-sincronas-entre-componentes-umbral}; Subdoc. 5, 5.1.6 y 5.4.5; Subdoc. 3, 3.4.2.5 y 3.4.3; Formulario T-14, paquetes 1.4.3.1, 1.4.3.5 y 1.4.5.2. Evidencia: Prueba de aceptación del resultado esperado (Anexo A.1) en la marcha blanca de la Etapa 1 \\
RF-ACR-005 Propio - Funcional & El sistema debe recibir la documentación de acreditación a través de un portal autenticado para empresas… & Cumple & Módulo Acreditación y Acceso (Portal de Contratistas y Consola de Portería) & Subdoc. 3, 3.4.2.5; Formulario T-14, paquete 1.4.5.2. Evidencia: Prueba de aceptación del resultado esperado (Anexo A.1) en la marcha blanca de la Etapa 1 \\
RF-ACR-006 Propio - Funcional & El sistema debe avisar del vencimiento próximo de cada documento de habilitación con 30, 15… & Cumple & Módulo Acreditación y Acceso (Portal de Contratistas y Consola de Portería) & Subdoc. 3, 3.4.2.5; Formulario T-14, paquete 1.4.3.5. Evidencia: Prueba de aceptación del resultado esperado (Anexo A.1) en la marcha blanca de la Etapa 1 \\
RF-ACR-007 Propio - Funcional & El sistema debe calcular el vencimiento de cada protocolo de vigilancia de salud según su periodicidad… & Cumple & Módulo Acreditación y Acceso (Portal de Contratistas y Consola de Portería) & Subdoc. 3, 3.4.2.5; Formulario T-14, paquetes 1.4.3.5 y 1.4.4.2. Evidencia: Prueba de aceptación del resultado esperado (Anexo A.1) en la marcha blanca de la Etapa 1 \\
RF-ACR-008 Propio - Funcional & El sistema debe verificar si la persona que presenta su credencial en el torniquete está habilitada… & Cumple & Módulo Acreditación y Acceso (Portal de Contratistas y Consola de Portería) & Subdoc. 3, 3.4.2.5; Formulario T-14, paquete 1.4.3.5. Evidencia: Prueba de aceptación del resultado esperado (Anexo A.1) en la marcha blanca de la Etapa 1 \\
RF-ACR-009 Propio - Funcional & El sistema debe verificar la identidad de la persona que entra o sale por portería… & Cumple & Módulo Acreditación y Acceso (Portal de Contratistas y Consola de Portería) & Subdoc. 3, 3.4.2.5 y 3.4.5; Formulario T-14, paquetes 1.1.5.2 y 1.4.4.4. Evidencia: Prueba de aceptación del resultado esperado (Anexo A.1) en la marcha blanca de la Etapa 2 \\
RF-ACR-010 Propio - Funcional & El sistema debe registrar el consentimiento informado de la persona cuando el tratamiento de un dato… & Cumple & Módulo Acreditación y Acceso (Portal de Contratistas y Consola de Portería): entidad Consentimiento por categoría sensible, con versión del texto, fecha de otorgamiento y fecha de revocación & Subdoc. 5, 5.2.6 y 5.1.6; Subdoc. 5, Anexo A.3, Tabla A.8 (Consentimiento); Formulario T-14, paquetes 1.4.4.2 y 1.4.4.4. Evidencia: Prueba de aceptación del resultado esperado (Anexo A.1) en la marcha blanca de la Etapa 2 \\
RF-ACR-011 Propio - Funcional & El sistema debe identificar a la persona operadora de un equipo compartido en el relevo… & Cumple & Módulo Acreditación y Acceso (Portal de Contratistas y Consola de Portería) & Subdoc. 3, 3.4.2.5; Formulario T-14, paquete 1.4.5.1. Evidencia: Prueba de aceptación del resultado esperado (Anexo A.1) en la marcha blanca de la Etapa 1 \\
RF-ACR-012 Propio - Funcional & El sistema debe autenticar a las personas contratistas en terreno sin correo corporativo ni contraseña escrita… & Cumple & Módulo Acreditación y Acceso (Portal de Contratistas y Consola de Portería) & Subdoc. 3, 3.4.2.5; Formulario T-14, paquetes 1.3.3.3, 1.4.3.5 y 1.4.5.1. Evidencia: Prueba de aceptación del resultado esperado (Anexo A.1) en la marcha blanca de la Etapa 1 \\
RF-ACR-013 Propio - Funcional & El sistema debe registrar el vínculo entre persona, empresa contratista y faena y, al término… & Cumple & Módulo Acreditación y Acceso (Portal de Contratistas y Consola de Portería) & Subdoc. 3, 3.4.2.5; Subdoc. 5, Anexo A.3, Tabla A.12 (Asignación de persona, vigencia\_hasta); Formulario T-14, paquetes 1.4.3.5 (baja automática) y 1.4.5.2 (cada empresa contratista administra la identidad de su personal). Evidencia: Prueba de aceptación del resultado esperado (Anexo A.1) en la marcha blanca de la Etapa 1. \\
RF-ACR-014 Propio - Funcional & El sistema debe conservar en cada registro operacional y de seguridad la identidad de la persona… & Cumple & Módulo Acreditación y Acceso (Portal de Contratistas y Consola de Portería) & Subdoc. 3, 3.4.2.5; Formulario T-14, paquete 1.4.3.1. Evidencia: Prueba de aceptación del resultado esperado (Anexo A.1) en la marcha blanca de la Etapa 1 \\
RF-ACR-015 Propio - Funcional & El sistema debe permitir a las personas usuarias externas resolver por su cuenta en el portal… & Cumple & Módulo Acreditación y Acceso (Portal de Contratistas y Consola de Portería) & Subdoc. 4, 4.1.8.3, Tabla \refext{S4}{tab:mecanismo-cifrado-custodia-claves-ambito}, \refext{S4}{tab:selecciones-principales-capas-1-3}.6.3 y 4.1.7.7; Subdoc. 3, 3.4.2.5; Formulario T-14, paquete 1.4.5.2 (consultas y recuperación de acceso de contratistas y compradores) y 1.4.4.5 (portal de compradores). Evidencia: Prueba de aceptación del resultado esperado (Anexo A.1) en la marcha blanca de la Etapa 1. \\
RF-ACR-016 Propio - Funcional & El sistema debe controlar el acceso por roles, complementado con atributos -sitio, empresa contratista y área-… & Cumple parcialmente & Módulo Acreditación y Acceso (Portal de Contratistas y Consola de Portería) & Subdoc. 4, 4.1.4.7 y 4.1.8.1; OVERFORE-Subdocumento4-Anexos, Anexo E (E.1, control 5.3); Formulario T-14, paquete 1.3.3.3 (modelo de roles y control por atributos). Falta: matriz de segregación de funciones documentada y verificable como entregable. Evidencia: Prueba de aceptación del resultado esperado (Anexo A.1) en la marcha blanca de la Etapa 1. \\
RF-ACR-017 Propio - Funcional & El sistema debe verificar la identidad de cada persona que ingresa a la sala técnica… & Cumple & Control de acceso a la sala técnica (Módulo Acreditación y Acceso) & Subdoc. 3, 3.3.5; detalle en Subdoc. 2, Anexo A.1; Formulario T-14, paquete 1.4.3.5. Evidencia: Prueba de aceptación del resultado esperado (Anexo A.1) en la marcha blanca de la Etapa 1 \\
RF-ACR-018 Propio - Funcional & El sistema debe registrar cada ingreso y cada egreso de la sala técnica en una bitácora… & Cumple & Control de acceso a la sala técnica (Módulo Acreditación y Acceso) & Subdoc. 3, 3.3.5; detalle en Subdoc. 2, Anexo A.1; Formulario T-14, paquete 1.4.3.5. Evidencia: Prueba de aceptación del resultado esperado (Anexo A.1) en la marcha blanca de la Etapa 1 \\
RF-CUM-001 Propio - Funcional & El sistema debe permitir registrar incidentes, cuasi accidentes, observaciones e inspecciones de seguridad en el lugar… & Cumple & Módulo Seguridad y Salud Ocupacional (App SITCO Terreno) & Subdoc. 3, 3.4.2.6; Formulario T-14, paquetes 1.4.4.2 y 1.4.5.1. Evidencia: Prueba de aceptación del resultado esperado (Anexo A.1) en la marcha blanca de la Etapa 2 \\
RF-CUM-002 Propio - Funcional & El sistema debe registrar, junto a cada registro de seguridad, la fecha y hora del hecho… & Cumple & Módulo Seguridad y Salud Ocupacional (App SITCO Terreno): incidentes e inspecciones con fecha del hecho y fecha de registro separadas & Subdoc. 5, Anexo A.3, Tabla A.9 (Incidente Inspección Acción sobre alerta); Subdoc. 3, 3.4.2.6; Formulario T-14, paquete 1.4.4.2. Evidencia: Prueba de aceptación del resultado esperado (Anexo A.1) en la marcha blanca de la Etapa 2 \\
RF-CUM-003 Propio - Funcional & El sistema debe generar el expediente de fiscalización en formato exportable, sin armar carpetas a mano. & Cumple & Módulo Seguridad y Salud Ocupacional (App SITCO Terreno) & Subdoc. 3, 3.4.2.6; Formulario T-14, paquete 1.4.4.2. Evidencia: Prueba de aceptación del resultado esperado (Anexo A.1) en la marcha blanca de la Etapa 2 \\
RF-ADM-005 Propio - Funcional & El sistema debe registrar toda operación que cree, modifique o elimine información: quién, cuándo, desde dónde… & Cumple & Módulo Administración y Servicios Comunes & Subdoc. 3, 3.4.3; Formulario T-14, paquetes 1.4.3.1 y 1.6.1.1. Evidencia: Prueba de aceptación del resultado esperado (Anexo A.1) en la marcha blanca de la Etapa 1 \\
RF-ADM-006 Propio - Funcional & El sistema debe permitir consultar y exportar la auditoría desde la interfaz, con filtros por persona… & Cumple parcialmente & Módulo Administración y Servicios Comunes & Subdoc. 3, 3.4.3; Subdoc. 4, 4.1.8.5, Tabla \refext{S4}{tab:modulo-plataforma-sostiene-cada-funcion-durante}; Subdoc. 5, Anexo A.3, Tabla A.12 (asiento de auditoría); Formulario T-14, paquete 1.4.3.1 (consulta y exportación de la auditoría con filtros). Falta: filtros por persona, período, entidad y tipo de operación, y registro de cada exportación de información sensible en la auditoría. Evidencia: Prueba de aceptación del resultado esperado (Anexo A.1) en la marcha blanca de la Etapa 1. \\
RF-ADM-011 Propio - Funcional & El sistema debe generar documentos -guías, actas, certificados, expedientes y reportes- a partir de plantillas administrables… & Cumple & Módulo Administración y Servicios Comunes & Subdoc. 3, 3.4.3; Subdoc. 5, Anexo A.4, Tabla A.15 (Plantilla documental o de notificación) y Anexo A.3, Tabla A.12 (documento con plantilla); Formulario T-14, paquete 1.4.3.1. Evidencia: Prueba de aceptación del resultado esperado (Anexo A.1) en la marcha blanca de la Etapa 1 \\
RF-NOT-001 Propio - Funcional & El sistema debe enviar notificaciones por correo electrónico, notificación en aplicación, mensajería instantánea para personas contratistas… & Cumple & Módulo Administración y Servicios Comunes & Subdoc. 3, 3.4.3; Formulario T-14, paquetes 1.4.2.1 y 1.4.3.1. Evidencia: Prueba de aceptación del resultado esperado (Anexo A.1) en la marcha blanca de la Etapa 1 \\
RF-DAT-002 Propio - Funcional & El sistema debe eliminar o anonimizar de forma segura y verificable los datos cuyo plazo… & Cumple & Módulo Administración y Servicios Comunes: procedimiento de eliminación segura y verificable por categoría de retención & Subdoc. 5, 5.2.4, Tabla \refext{S5}{tab:retencion-categoria-entidad} y \refext{S5}{tab:sistema-maestro-operacion-prohibida-entidad}.5, Tabla \refext{S5}{tab:procedimiento-eliminacion-segura-evidencia-cada-paso}; Formulario T-14, paquetes 1.3.2.3 y 1.4.3.1. Evidencia: Prueba de aceptación del resultado esperado (Anexo A.1) en la marcha blanca de la Etapa 1 \\
RF-CUM-004 Propio - Funcional & El sistema debe soportar el programa de inspecciones planificadas de seguridad, cargado mensualmente, con su calendario… & Cumple parcialmente & Módulo Seguridad y Salud Ocupacional (App SITCO Terreno) & Subdoc. 5, Anexo A.3, Tabla A.9 (Inspección: fecha programada, responsable y resultado); Formulario T-14, paquete 1.4.4.2. Falta: calendario diario de cada inspección. Evidencia: Prueba de aceptación del resultado esperado (Anexo A.1) en la marcha blanca de la Etapa 2 \\
RF-ADM-010 Propio - Funcional & El sistema debe gestionar documentos con versiones, metadatos, control de acceso, búsqueda por contenido… & Cumple & Módulo Administración y Servicios Comunes & Subdoc. 3, 3.4.3; Formulario T-14, paquete 1.4.3.1. Evidencia: Prueba de aceptación del resultado esperado (Anexo A.1) en la marcha blanca de la Etapa 1 \\
RF-CUM-005 Propio - Funcional & El sistema debe recibir las alertas de fatiga y somnolencia desde la plataforma de las cámaras… & Cumple & Módulo Seguridad y Salud Ocupacional (App SITCO Terreno) & Subdoc. 3, 3.4.2.6; Formulario T-14, paquetes 1.4.2.3 y 1.4.4.2. Evidencia: Prueba de aceptación del resultado esperado (Anexo A.1) en la marcha blanca de la Etapa 2 \\
RF-CUM-006 Propio - Funcional & El sistema debe conservar durante 10 años el histórico completo de alertas de fatiga… & Cumple & Módulo Seguridad y Salud Ocupacional (App SITCO Terreno): entidades Alerta de fatiga y Acción sobre alerta con retención de 10 años & Subdoc. 5, 5.2.4, Tabla \refext{S5}{tab:retencion-categoria-entidad}; Subdoc. 4, 4.1.6.3; Formulario T-14, paquetes 1.3.2.3 y 1.4.4.2. Evidencia: Prueba de aceptación del resultado esperado (Anexo A.1) en la marcha blanca de la Etapa 2 \\
RF-CUM-007 Propio - Funcional & El sistema debe restringir el acceso al dato individual de las alertas de fatiga a perfiles… & Cumple & Módulo Seguridad y Salud Ocupacional (App SITCO Terreno): cifrado de campo, acceso por rol y asiento de auditoría en cada lectura del dato sensible & Subdoc. 5, 5.2.6; Subdoc. 4, 4.1.8.1 y 4.1.8.5, Tabla \refext{S4}{tab:modulo-plataforma-sostiene-cada-funcion-durante}; Subdoc. 5, Anexo A.3, Tabla A.9 (Alerta de fatiga); Formulario T-14, paquete 1.4.4.2. Evidencia: Prueba de aceptación del resultado esperado (Anexo A.1) en la marcha blanca de la Etapa 2 \\
RF-CUM-008 Propio - Funcional & El sistema debe capturar automáticamente los datos de la instrumentación de monitoreo ambiental y geotécnico… & Cumple & Módulo Cumplimiento Ambiental & Subdoc. 3, 3.4.2.7; Formulario T-14, paquete 1.4.4.3. Evidencia: Prueba de aceptación del resultado esperado (Anexo A.1) en la marcha blanca de la Etapa 2 \\
RF-CUM-009 Propio - Funcional & El sistema debe recibir por integración los resultados de los laboratorios externos de calidad del aire… & Cumple & Módulo Cumplimiento Ambiental: informes de laboratorios externos por la puerta de exposición, con identidad externa y credencial de cliente & Subdoc. 4, 4.1.4.5, 4.1.6.3 y 4.1.8.3, Tabla \refext{S4}{tab:modulo-plataforma-sostiene-cada-funcion-durante}; Subdoc. 3, 3.4.2.7; Formulario T-14, paquete 1.4.4.3. Evidencia: Prueba de aceptación del resultado esperado (Anexo A.1) en la marcha blanca de la Etapa 2 \\
RF-CUM-010 Propio - Funcional & El sistema debe registrar el consumo de agua por punto de medición y permitir declarar… & Cumple & Módulo Cumplimiento Ambiental & Subdoc. 3, 3.4.2.7; Formulario T-14, paquete 1.4.4.3. Evidencia: Prueba de aceptación del resultado esperado (Anexo A.1) en la marcha blanca de la Etapa 2 \\
RF-CUM-011 Propio - Funcional & El sistema debe registrar el combustible efectivamente consumido por equipo y por área, obtenido… & Cumple & Módulo Cumplimiento Ambiental & Subdoc. 3, 3.4.2.7; Formulario T-14, paquete 1.4.4.3. Evidencia: Prueba de aceptación del resultado esperado (Anexo A.1) en la marcha blanca de la Etapa 2 \\
RF-CUM-012 Propio - Funcional & El sistema debe generar los reportes a la autoridad ambiental desde los datos registrados, con trazabilidad… & Cumple & Módulo Cumplimiento Ambiental & Subdoc. 3, 3.4.2.7; Formulario T-14, paquete 1.4.4.3. Evidencia: Prueba de aceptación del resultado esperado (Anexo A.1) en la marcha blanca de la Etapa 2 \\
RF-CUM-013 Propio - Funcional & El sistema debe registrar los compromisos de la Resolución de Calificación Ambiental con su frecuencia… & Cumple & Módulo Cumplimiento Ambiental & Subdoc. 3, 3.4.2.7; Formulario T-14, paquete 1.4.4.3. Evidencia: Prueba de aceptación del resultado esperado (Anexo A.1) en la marcha blanca de la Etapa 2 \\
RF-CUM-014 Propio - Funcional & El sistema debe construir los indicadores de empleo local y de compras locales del convenio… & Cumple & Módulo Cumplimiento Ambiental & Subdoc. 3, 3.4.2.7; Formulario T-14, paquete 1.4.4.6. Evidencia: Prueba de aceptación del resultado esperado (Anexo A.1) en la marcha blanca de la Etapa 2 \\
RF-CUM-015 Propio - Funcional & El sistema debe calcular cada mes la huella de carbono por tonelada de cobre fino… & Cumple & Módulo Cumplimiento Ambiental & Subdoc. 3, 3.4.2.7; Formulario T-14, paquetes 1.4.4.3 y 1.8.7.3. Evidencia: Prueba de aceptación del resultado esperado (Anexo A.1) en la marcha blanca de la Etapa 2 \\
RF-CUM-016 Propio - Funcional & El sistema debe guardar con versiones los factores de emisión y la evidencia de cada dato… & Cumple & Módulo Cumplimiento Ambiental: huella con datos de actividad y factores versionados & Subdoc. 5, 5.1.6; Subdoc. 5, Anexo A.3, Tabla A.10 (Huella: dato\_actividad e id\_version\_factor); Subdoc. 3, 3.4.2.7; Formulario T-14, paquete 1.4.4.3. Evidencia: Prueba de aceptación del resultado esperado (Anexo A.1) en la marcha blanca de la Etapa 2 \\
RF-CUM-017 Propio - Funcional & El sistema debe generar el expediente de verificación por tercero independiente con los datos de actividad… & Cumple & Módulo Cumplimiento Ambiental & Subdoc. 3, 3.4.2.7; Formulario T-14, paquetes 1.4.4.3 y 1.8.7.3. Evidencia: Prueba de aceptación del resultado esperado (Anexo A.1) en la marcha blanca de la Etapa 2 \\
RF-EQU-001 Propio - Funcional & El sistema debe mostrar en una sola vista, para cada equipo, su estado actual, su historial… & Cumple & Módulo Equipos y Mantenimiento (vista única del equipo) & Subdoc. 3, 3.4.2.4; Formulario T-14, paquete 1.4.4.1. Evidencia: Prueba de aceptación del resultado esperado (Anexo A.1) en la marcha blanca de la Etapa 2 \\
RF-EQU-002 Propio - Funcional & El sistema debe integrar la telemetría de los portales de los tres fabricantes de equipos… & Cumple & Módulo Equipos y Mantenimiento (vista única del equipo) & Subdoc. 3, 3.4.2.4; Formulario T-14, paquetes 1.2.4.1 y 1.4.2.3. Evidencia: Prueba de aceptación del resultado esperado (Anexo A.1) en la marcha blanca de la Etapa 2 \\
RF-EQU-003 Propio - Funcional & El sistema debe permitir cambiar la fuente de telemetría de un fabricante sin afectar la vista… & Cumple & Módulo Equipos y Mantenimiento (vista única del equipo) & Subdoc. 3, 3.4.2.4; Formulario T-14, paquetes 1.2.4.2, 1.3.5.2 y 1.4.4.1. Evidencia: Prueba de aceptación del resultado esperado (Anexo A.1) en la marcha blanca de la Etapa 2 \\
RF-EQU-004 Propio - Funcional & El sistema debe exponer las pautas de mantenimiento y las órdenes de trabajo pendientes… & Cumple & Módulo Equipos y Mantenimiento (vista única del equipo) & Subdoc. 3, 3.4.2.4; Formulario T-14, paquetes 1.4.2.3 y 1.4.4.1. Evidencia: Prueba de aceptación del resultado esperado (Anexo A.1) en la marcha blanca de la Etapa 2 \\
RF-EQU-005 Propio - Funcional & El sistema debe permitir levantar un aviso de falla desde la tablet, en no más… & Cumple & Módulo Equipos y Mantenimiento (vista única del equipo) & Subdoc. 3, 3.4.2.4; Formulario T-14, paquetes 1.4.4.1 y 1.4.5.1. Evidencia: Prueba de aceptación del resultado esperado (Anexo A.1) en la marcha blanca de la Etapa 2 \\
RF-EQU-006 Propio - Funcional & El sistema debe conservar el aviso de falla levantado en zona sin cobertura, con integridad… & Cumple & Módulo Equipos y Mantenimiento (vista única del equipo) & Subdoc. 3, 3.4.2.4; Formulario T-14, paquetes 1.4.2.2 y 1.4.5.1. Evidencia: Prueba de aceptación del resultado esperado (Anexo A.1) en la marcha blanca de la Etapa 2 \\
RF-CNX-001 Propio - Funcional & El sistema debe sincronizar automáticamente, sin intervención manual, los registros hechos en faena durante un corte… & Cumple & Nodo de Faena, Integrador de Sistemas y dispositivos de terreno & Subdoc. 3, 3.3.6 y 3.4.3; Formulario T-14, paquetes 1.4.2.2 y 1.4.5.1. Evidencia: Prueba de aceptación del resultado esperado (Anexo A.1) en la marcha blanca de la Etapa 1 \\
RF-CNX-002 Propio - Funcional & Sin enlace, el sistema debe validar la credencial y las autorizaciones de una persona… & Cumple & Nodo de Faena, Integrador de Sistemas y dispositivos de terreno & Subdoc. 3, 3.3.6 y 3.4.3; Formulario T-14, paquete 1.4.2.2. Evidencia: Prueba de aceptación del resultado esperado (Anexo A.1) en la marcha blanca de la Etapa 1 \\
RF-CNX-003 Propio - Funcional & El dispositivo de terreno debe operar sin enlace desde el origen: guardar cada registro cifrado… & Cumple & Nodo de Faena, Integrador de Sistemas y dispositivos de terreno & Subdoc. 3, 3.3.6 y 3.4.3; Formulario T-14, paquete 1.4.5.1. Evidencia: Prueba de aceptación del resultado esperado (Anexo A.1) en la marcha blanca de la Etapa 1 \\
RF-CNX-004 Propio - Funcional & El sistema debe administrar de forma remota y centralizada los dispositivos de terreno y los gabinetes… & Cumple & Nodo de Faena, Integrador de Sistemas y dispositivos de terreno & Subdoc. 3, 3.3.6 y 3.4.3; Formulario T-14, paquetes 1.4.2.2 y 1.4.5.1. Evidencia: Prueba de aceptación del resultado esperado (Anexo A.1) en la marcha blanca de la Etapa 1 \\
RF-ADM-001 Propio - Funcional & El sistema debe mantener accesibles, en un repositorio de consulta de solo lectura, los datos históricos… & Cumple & Módulo Administración y Servicios Comunes: repositorio de consulta histórica de solo lectura por módulo & Subdoc. 5, 5.3.5; Subdoc. 4, 4.1.4.6; Formulario T-14, paquetes 1.4.6.5 (repositorio de consulta de solo lectura, meses 10 a 20) y 1.3.2.3 (plazos de retención por categoría y método de eliminación). Evidencia: Prueba de aceptación del resultado esperado (Anexo A.1) en la marcha blanca de la Etapa 1 \\
RF-ADM-002 Propio - Funcional & El sistema debe permitir al Cliente, sin intervención del Adjudicatario, administrar personas usuarias, roles, permisos, unidades… & Cumple & Módulo Administración y Servicios Comunes & Subdoc. 3, 3.4.3; Formulario T-14, paquete 1.4.3.1. Evidencia: Prueba de aceptación del resultado esperado (Anexo A.1) en la marcha blanca de la Etapa 1 \\
RF-ADM-003 Propio - Funcional & El sistema debe permitir configurar desde la interfaz de administración las reglas de negocio parametrizables -umbrales… & Cumple & Módulo Administración y Servicios Comunes & Subdoc. 3, 3.4.3; Formulario T-14, paquete 1.4.3.1. Evidencia: Prueba de aceptación del resultado esperado (Anexo A.1) en la marcha blanca de la Etapa 1 \\
RF-ADM-004 Propio - Funcional & El sistema debe exigir la aprobación de una segunda persona para todo cambio de parámetro… & Cumple & Módulo Administración y Servicios Comunes: versión de parámetro con segundo perfil aprobador y justificación & Subdoc. 4, 4.1.8.5, Tabla \refext{S4}{tab:modulo-plataforma-sostiene-cada-funcion-durante}; Subdoc. 4, Anexos, E.1 control 5.3; Subdoc. 5, Anexo A.3, Tabla A.12 (Parámetro: justificacion e id\_aprobador\_segundo\_perfil); Formulario T-14, paquete 1.4.3.1. Evidencia: Prueba de aceptación del resultado esperado (Anexo A.1) en la marcha blanca de la Etapa 1 \\
RF-ADM-007 Propio - Funcional & El sistema debe registrar el ciclo de vida completo de cada cuenta de usuario de modo… & Cumple & Módulo Administración y Servicios Comunes: auditoría inalterable con sello de integridad y estados de cuenta activa, bloqueada y baja & Subdoc. 4, 4.1.8.5, Tabla \refext{S4}{tab:volumetria-referencia-capacidad-diseno}; Subdoc. 5, 5.2.4, Tabla \refext{S5}{tab:retencion-categoria-entidad}; Subdoc. 5, Anexo A.3, Tabla A.12 (Persona, Asiento de auditoría); Subdoc. 4, Anexos, E.4 control 8.2; Formulario T-14, paquetes 1.3.3.3 y 1.4.3.1. Evidencia: Prueba de aceptación del resultado esperado (Anexo A.1) en la marcha blanca de la Etapa 1 \\
RF-ADM-008 Propio - Funcional & El sistema debe soportar flujos de trabajo con estados, transiciones, responsables, plazos, escalamiento automático por vencimiento… & Cumple & Módulo Administración y Servicios Comunes & Subdoc. 3, 3.4.3; Formulario T-14, paquete 1.4.3.1. Evidencia: Prueba de aceptación del resultado esperado (Anexo A.1) en la marcha blanca de la Etapa 1 \\
RF-DAT-003 Propio - Funcional & El sistema debe permitir la carga y descarga masiva de información en formatos abiertos, con validación… & Cumple parcialmente & Módulo Administración y Servicios Comunes & Subdoc. 5, 5.2.2, Tabla \refext{S5}{tab:vias-explotacion-informacion}; Subdoc. 5, Anexo A.4, Tabla A.15 (Trabajo de importación o exportación); Formulario T-14, paquete 1.4.3.1. Falta: formato abierto de los archivos de carga masiva. Evidencia: Prueba de aceptación del resultado esperado (Anexo A.1) en la marcha blanca de la Etapa 1 \\
RF-ADM-009 Propio - Funcional & El sistema debe mostrar a cada persona usuaria, en una bandeja de tareas unificada, todas… & Cumple & Módulo Administración y Servicios Comunes & Subdoc. 3, 3.4.3; Subdoc. 5, Anexo A.3, Tabla A.12 (Caso de flujo, Tarea y Transición); Formulario T-14, paquete 1.4.3.1. Evidencia: Prueba de aceptación del resultado esperado (Anexo A.1) en la marcha blanca de la Etapa 1 \\
RF-ADM-012 Propio - Funcional & El sistema debe ofrecer una búsqueda global por el contenido de los registros, que encuentre resultados… & Cumple parcialmente & Módulo Administración y Servicios Comunes & Subdoc. 4, 4.1.1.2 (Tabla \refext{S4}{tab:tecnologias-capa-6-alternativa-evaluada}, OpenSearch); Subdoc. 5, 5.2.1 (Tabla \refext{S5}{tab:naturaleza-dato-motor-razon-eleccion}); Formulario T-14, paquete 1.4.3.1. Falta: filtro por categorías en la búsqueda global. Evidencia: Prueba de aceptación del resultado esperado (Anexo A.1) en la marcha blanca de la Etapa 1 \\
RF-ADM-013 Propio - Funcional & El sistema debe entregar listados ordenables, filtrables y paginados, exportables en formatos abiertos con el filtro… & Cumple parcialmente & Módulo Administración y Servicios Comunes & Subdoc. 5, Anexo A.4, Tabla A.15 (Trabajo de importación o exportación); Subdoc. 5, 5.2.2, Tabla \refext{S5}{tab:vias-explotacion-informacion} y \refext{S5}{tab:posicion-cap-dominio}.5; Formulario T-14, paquete 1.4.3.1. Falta: listados ordenables, filtrables y paginados; umbral de 50 MB y aviso al terminar la exportación. Evidencia: Prueba de aceptación del resultado esperado (Anexo A.1) en la marcha blanca de la Etapa 1 \\
RF-NOT-002 Propio - Funcional & El sistema debe permitir a cada persona elegir el canal y la frecuencia con que recibe… & Cumple & Módulo Administración y Servicios Comunes: preferencias de canal y frecuencia con baja de informativas y avisos obligatorios no cancelables & Subdoc. 5, Anexo A.4, Tabla A.15 (Preferencia e informe programado); Subdoc. 3, 3.4.3; Formulario T-14, paquete 1.4.3.1. Evidencia: Prueba de aceptación del resultado esperado (Anexo A.1) en la marcha blanca de la Etapa 1 \\
RF-NOT-003 Propio - Funcional & El sistema debe permitir al Cliente administrar las plantillas de notificación, con variables de la transacción… & Cumple & Módulo Administración y Servicios Comunes: plantillas de notificación versionadas, con variables, modificables sin despliegue de código & Subdoc. 5, Anexo A.4, Tabla A.15 (Plantilla documental o de notificación); Subdoc. 5, Anexo A.3, Tabla A.12; Formulario T-14, paquete 1.4.3.1. Evidencia: Prueba de aceptación del resultado esperado (Anexo A.1) en la marcha blanca de la Etapa 1 \\
RF-ANA-001 Propio - Funcional & El sistema debe ofrecer tableros operacionales y de gestión con los indicadores del Caso, filtrables… & Cumple & Módulo Indicadores y Centro Integrado & Subdoc. 3, 3.4.3; Formulario T-14, paquetes 1.4.3.6 y 1.4.4.6. Evidencia: Prueba de aceptación del resultado esperado (Anexo A.1) en la marcha blanca de la Etapa 1 \\
RF-ANA-002 Propio - Funcional & El sistema debe permitir al Cliente construir sus propios informes sin intervención del Adjudicatario, mediante… & Cumple & Módulo Indicadores y Centro Integrado & Subdoc. 3, 3.4.3; Formulario T-14, paquete 1.4.4.6. Evidencia: Prueba de aceptación del resultado esperado (Anexo A.1) en la marcha blanca de la Etapa 2 \\
RF-ANA-003 Propio - Funcional & El sistema debe permitir exportar todo informe en formatos abiertos y programar su envío automático… & Cumple & Módulo Indicadores y Centro Integrado: informes programados con exportación en formatos abiertos & Subdoc. 5, 5.2.2, Tabla \refext{S5}{tab:vias-explotacion-informacion}; Subdoc. 5, Anexo A.4, Tabla A.15 (Preferencia e informe programado); Formulario T-14, paquete 1.4.4.6. Evidencia: Prueba de aceptación del resultado esperado (Anexo A.1) en la marcha blanca de la Etapa 2 \\
RF-TER-009 Propio - Funcional & El sistema debe calcular, a partir de los parámetros de perforación de cada pozo, un índice… & Cumple & Módulo Operación Mina (Botonera de Pala, App SITCO Terreno y conector de despacho de flota) & Subdoc. 4, 4.1.6.1.1; Subdoc. 5, Anexo A.3, Tabla A.4 (Pozo: indice\_dureza, id\_version\_regla\_dureza y parametros\_perforacion); Formulario T-14, paquete 1.4.3.2. Evidencia: Prueba de aceptación del resultado esperado (Anexo A.1) en la marcha blanca de la Etapa 1 \\
RF-EQU-007 Propio - Funcional & El sistema debe registrar el tiempo transcurrido entre la detención de un equipo de mina… & Cumple parcialmente & Módulo Equipos y Mantenimiento (vista única del equipo) & Subdoc. 4, 4.1.6.2; Subdoc. 5, Anexo A.3, Tabla A.7 (Detención de planta, fecha\_diagnostico); Formulario T-14, paquete 1.4.4.1. Falta: tiempo medio entre falla y diagnóstico frente a la meta de 1 hora del Cliente. Evidencia: Prueba de aceptación del resultado esperado (Anexo A.1) en la marcha blanca de la Etapa 2 \\
RF-PLA-010 Propio - Funcional & El sistema debe medir la pérdida efectiva de tiempo en el cambio de turno… & Cumple & Módulo Operación Mina (Botonera de Pala, App SITCO Terreno y conector de despacho de flota) & Subdoc. 3, 3.4.2.1 y 3.3.3; Formulario T-14, paquetes 1.4.3.2 y 1.4.3.6. Evidencia: Prueba de aceptación del resultado esperado (Anexo A.1) en la marcha blanca de la Etapa 1 \\
RF-CON-016 Propio - Funcional & El sistema debe registrar el tiempo transcurrido entre la toma de cada muestra y la disponibilidad… & Cumple & Módulo Conciliación Metalúrgica & Subdoc. 5, Anexo A.3, Tabla A.6 (Muestra: hora\_toma; Análisis: hora\_publicacion); Subdoc. 4, 4.1.6.1.3; Formulario T-14, paquete 1.4.3.4. Evidencia: Prueba de aceptación del resultado esperado (Anexo A.1) en la marcha blanca de la Etapa 1 \\
RF-CON-017 Propio - Funcional & El sistema debe enviar al sistema de gestión de laboratorio la identificación de cada muestra tomada… & Cumple & Módulo Conciliación Metalúrgica & Subdoc. 3, 3.3.2.2; Formulario T-14, paquete 1.4.2.1. Evidencia: Prueba de aceptación del resultado esperado (Anexo A.1) en la marcha blanca de la Etapa 1 \\
RF-CNX-005 Propio - Funcional & El sistema debe registrar cada corte del enlace con Antofagasta, con su inicio, su término… & Cumple & Nodo de Faena (registro de cada corte) e indicadores de consola operativa para el área de TI del cliente & Subdoc. 4, 4.1.3.2 y 4.1.4.8; Subdoc. 3, 3.3.6; Formulario T-14, paquete 1.4.2.2. Evidencia: Prueba de aceptación del resultado esperado (Anexo A.1) en la marcha blanca de la Etapa 1 \\
RF-TRZ-019 Propio - Funcional & El sistema debe registrar la liquidación que informa el comprador de cada embarque con sus componentes… & Cumple & Módulo Trazabilidad y Embarque (Portal para Compradores en Etapa 2) & Subdoc. 4, 4.1.6.1.2; Subdoc. 5, 5.1.6, Figura \refext{S5}{fig:trazabilidad-embarque-parte-2b-embarque-hitos}; Subdoc. 5, Anexo A.3, Liquidación y Componente de liquidación; Subdoc. 3, 3.3.2.2; Formulario T-14, paquete 1.4.3.3. Evidencia: Prueba de aceptación del resultado esperado (Anexo A.1) en la marcha blanca de la Etapa 1 \\
RF-TRZ-020 Propio - Funcional & El sistema debe calcular, para cada lote comercial en conformación, el contenido de arsénico… & Cumple & Módulo Trazabilidad y Embarque (Portal para Compradores en Etapa 2) & Subdoc. 3, 3.4.2.2; Formulario T-14, paquete 1.7.1.1. Evidencia: Prueba de aceptación del resultado esperado (Anexo A.1) en la marcha blanca de la Etapa 1 \\
RF-EQU-008 Propio - Funcional & El sistema debe generar alertas de mantenimiento basadas en reglas definidas por el propio equipo… & Cumple & Módulo Equipos y Mantenimiento (alertas por reglas propias del equipo de mantenimiento) & Subdoc. 3, 3.4.2.4; Formulario T-14, paquete 1.4.4.1. Evidencia: Prueba de aceptación del resultado esperado (Anexo A.1) en la marcha blanca de la Etapa 2 \\
\end{tabla}
\begin{tabla}{I{0.128}I{0.304}I{0.108}I{0.147}I{0.314}}{Requerimientos no funcionales propios del caso.}{tab:requerimientos-no-funcionales-propios-caso}
  {\cab{ID requerimiento} & \cab{Descripción} & \cab{Estado} & \cab{Componente que lo satisface} & \cab{Sección de la propuesta y evidencia}}
  {\emph{Fuente:} elaboración propia a partir de las Bases técnicas del caso Minería.}
RNF-CAP-001 Propio - No funcional & El sistema debe poder crecer hasta 3 veces la actividad actual de movimiento de material… & Cumple & Arquitectura escalable de la solución (factor de diseño de 3 veces, crecimiento a 600 TPS sin cambiar la topología) & Subdoc. 4, 4.2.4.1 y 4.2.4.2, Tabla \refext{S4}{tab:volumetria-referencia-capacidad-diseno}; Subdoc. 4, 4.2.4.6; OVERFORE-Subdocumento4-Anexos, Anexo D (D.1 y D.3.1); Formulario T-14, paquete 1.3.1.2. Evidencia: Prueba de carga a 1,5 veces el peak declarado (momento de mayor actividad, RT-09.06) y demostración en Preproducción del escalamiento hasta 3 veces la volumetría actual sin cambio de arquitectura. Prueba de carga con el perfil de concurrencia derivado del numeral 14.1 y declarado conforme al numeral 14.2. \\
RNF-CAP-002 Propio - No funcional & El sistema debe absorber hasta 3 veces el volumen actual de muestras y resultados de laboratorio… & Cumple & Arquitectura escalable de la solución (factor de diseño de 3 veces sobre los hechos de negocio, incluidas las muestras de laboratorio) & Subdoc. 4, 4.2.4.1 y Tabla \refext{S4}{tab:volumetria-referencia-capacidad-diseno}; Subdoc. 4, 4.2.4.6; Subdoc. 4, Anexo D.3.1 (195 muestras por día) y D.1 fila 1 (capacidad de diseño de 3 veces); Formulario T-14, paquete 1.3.1.2. Evidencia: Prueba de carga a 1,5 veces el peak declarado (RT-09.06) y demostración en Preproducción del escalamiento hasta 3 veces la volumetría actual sin cambio de arquitectura. Prueba de carga sobre la integración con el sistema de gestión de laboratorio en el peak de recepción de resultados. \\
RNF-CAP-003 Propio - No funcional & El sistema debe absorber hasta 3 veces el volumen actual de control de acceso, es decir… & Cumple parcialmente & Arquitectura escalable de la solución (factor de diseño de 3 veces; 2.900 pasos de portería por día como base) & Subdoc. 4, 4.2.4.1, 4.2.4.2 y Tabla \refext{S4}{tab:volumetria-referencia-capacidad-diseno}; Formulario T-14, paquete 1.3.1.2. Falta: el Anexo D de OVERFORE-Subdocumento4-Anexos dimensiona 8.895 personas y 106 contrapartes, y no 25.500 personas y 270 empresas. Evidencia: Prueba de carga a 1,5 veces el peak declarado (RT-09.06) y demostración en Preproducción del escalamiento hasta 3 veces la volumetría actual sin cambio de arquitectura. Prueba de carga en el peak de cambio de turno y simulación del escenario de detención mayor de planta. \\
RNF-CAP-004 Propio - No funcional & El sistema debe poder crecer hasta 3 veces la telemetría y las señales de proceso actuales… & Cumple parcialmente & Arquitectura escalable de la solución (series de tiempo con capacidad de diseño de 3.269 lecturas por segundo, 3 veces las 1.090 actuales) & Subdoc. 4, 4.2.4.1 y 4.2.4.4; OVERFORE-Subdocumento4-Anexos, Anexo D (D.3.1; Tabla D.3, fila C1); Subdoc. 4, 4.1.7.2, Tabla \refext{S4}{tab:campos-obligatorios-evento-dominio} (vista única en 5 s); Formulario T-14, paquetes 1.3.1.2 y 1.5.3.1. Falta: demostración de que la vista única de equipo se mantiene en 5 segundos o menos con 3 veces el volumen. Evidencia: Prueba de carga a 1,5 veces el peak declarado (RT-09.06) y demostración en Preproducción del escalamiento hasta 3 veces la volumetría actual sin cambio de arquitectura. Prueba de carga con el volumen proyectado y medición de la latencia de la vista de equipo. \\
RNF-CAP-005 Propio - No funcional & El sistema debe absorber hasta 3 veces el volumen actual de registros de seguridad, es decir… & Cumple & Arquitectura escalable de la solución (factor de diseño de 3 veces sobre los hechos de seguridad: incidentes, inspecciones y alertas de fatiga) & Subdoc. 4, 4.2.4.1 y Tabla \refext{S4}{tab:volumetria-referencia-capacidad-diseno}; Subdoc. 4, 4.2.4.4; Subdoc. 4, Anexo D.3.1 (32 registros, 21 inspecciones y 13 alertas por día); Formulario T-14, paquete 1.3.1.2. Evidencia: Prueba de carga a 1,5 veces el peak declarado (RT-09.06) y demostración en Preproducción del escalamiento hasta 3 veces la volumetría actual sin cambio de arquitectura. Prueba de carga con el volumen proyectado y verificación del tiempo de generación del expediente de fiscalización. \\
RNF-CAP-006 Propio - No funcional & El sistema debe almacenar localmente, en faena y en los dispositivos de terreno, hasta 150.000 eventos… & Cumple & Nodo de Faena (sala técnica, Etapa 1) & Subdoc. 4, 4.1.4.5 y 4.2.4.5 (cola persistente de 150.000 eventos en el Nodo de Faena; cada gabinete conserva su cola local durante un corte); Formulario T-14, paquete 1.4.2.2. Evidencia: Corte controlado de 48 horas con medición del volumen acumulado por nodo y del margen restante. \\
RNF-CAP-007 Propio - No funcional & El sistema debe sostener 40 transacciones por segundo en régimen normal y 200 en el momento… & Cumple & Dimensionamiento transaccional: 40 TPS en régimen, 200 TPS de punta declarada, prueba a 300 TPS y crecimiento a 600 TPS & Subdoc. 4, 4.2.4 y 4.2.4.2, Tabla \refext{S4}{tab:volumetria-referencia-capacidad-diseno}; Subdoc. 5, 5.4.1; Subdoc. 9, 9.2.4.5; Formulario T-14, paquete 1.5.3.1 (300 TPS, 1,5 veces el pico de 200 TPS, durante 60 minutos). Evidencia: Prueba de carga con el perfil derivado del numeral 14.1 y declarado conforme al numeral 14.2 \\
RNF-CAP-008 Propio - No funcional & El sistema debe escalar horizontal y automáticamente sus capas de aplicación e integración, sin intervención manual… & Cumple & Clúster AKS de servicios con 4 a 9 nodos y réplicas por servicio, con escalado anticipado antes de cada cambio de turno & Subdoc. 4, 4.2.2.2; Subdoc. 4, 4.2.4.3 y 4.2.4.6; OVERFORE-Subdocumento4-Anexos, Anexo D (D.3.3, Tabla D.6); Formulario T-14, paquetes 1.3.1.2 y 1.5.3.1. Evidencia: Prueba de carga creciente en Preproducción hasta gatillar el escalamiento, verificando las transacciones durante el evento. \\
RNF-CUM-001 Propio - No funcional & El tratamiento de los datos personales sensibles debe cumplir la Ley N° 21.719 y el artículo… & Cumple & Módulos Seguridad y Salud Ocupacional y Acreditación y Acceso (catálogo de tratamientos con base autorizada por el cliente, cifrado de campo y eliminación al vencer la retención) & Subdoc. 5, 5.2.6; Subdoc. 5, 5.2.4 y 5.2.5; Subdoc. 5, Anexo A.3 (id\_version\_base\_tratamiento por registro); Subdoc. 4, 4.1.8.4, Tabla \refext{S4}{tab:amenazas-mayor-consecuencia-componentes-expuestos-control}; Formulario T-14, paquetes 1.3.2.3 y 1.3.6.2. Evidencia: Revisión con el área legal del Cliente, sobre cada categoría de dato, de que el 100 \% de los registros tenga su base legal de tratamiento documentada y trazable, más prueba de eliminación o anonimización verificable sobre una muestra de registros con plazo de retención vencido. \\
RNF-CUM-002 Propio - No funcional & El sistema debe firmar con firma electrónica avanzada los registros de seguridad exigidos por la autoridad… & Cumple & Firma electrónica avanzada (Ley N° 19.799) del módulo Administración y Servicios Comunes, con prestador acreditado & Subdoc. 3, 3.4.2.6, 3.4.2.7 y 3.4.3; Subdoc. 4, 4.1.7.1, Tabla \refext{S4}{tab:llamadas-sincronas-entre-componentes-umbral}; Subdoc. 5, Anexo A.3, Documento firmado (serie del certificado, sello de tiempo y evidencia de firma); Formulario T-14, paquete 1.4.3.1. Evidencia: Verificación de la firma avanzada y de la validez del certificado al emitirse, y de la firma y del sello de tiempo sobre una muestra de documentos de cada tipo, incluidos documentos cuyo certificado de firma ya esté vencido (real o de prueba con vencimiento configurado), dentro del plazo de retención. \\
RNF-DES-001 Propio - No funcional & El sistema debe registrar un movimiento de material desde el equipo en terreno, incluido el aviso… & Cumple & Botonera de Pala y Nodo de Faena & Subdoc. 3, 3.2.4.3 y 3.3.3; Formulario T-14, paquetes 1.4.2.2 y 1.6.3.1. Evidencia: Medición instrumentada sobre al menos 1.000 registros reales en operación normal, con el dispositivo definitivo de terreno. La medición incluye el aviso de cambio de material y se realiza durante la marcha blanca (período de prueba en paralelo con el método actual), en turno de día y de noche. \\
RNF-DES-002 Propio - No funcional & El sistema debe completar el registro en la báscula de camiones en un máximo de 3… & Cumple & Registro de báscula sin hardware adicional (lectura de la báscula y QR en tableta) & Subdoc. 3, 3.2.4.3 y 3.3.3; Formulario T-14, paquete 1.4.3.3. Evidencia: Medición sobre 200 pesajes consecutivos en operación normal. \\
RNF-DES-003 Propio - No funcional & El sistema debe verificar la habilitación de una persona en portería en no más de 1,5… & Cumple & Gabinete de borde de portería con copia local de personas habilitadas & Subdoc. 3, 3.2.4.3 y 3.3.5; Formulario T-14, paquete 1.4.2.2. Evidencia: Medición sobre 500 ingresos en el peak de cambio de turno, con y sin enlace disponible. \\
RNF-DES-004 Propio - No funcional & El sistema debe completar, desde la lectura de la credencial de proximidad en cabina… & Cumple & Botonera de Pala (toque de credencial al inicio del turno) con copia local de identidades en el Nodo de Faena & Subdoc. 4, 4.1.5.1, Tabla \refext{S4}{tab:interfaces-sincronas-eventos-publicados-datos-consumidos}; Subdoc. 4, 4.1.7.2, Tabla \refext{S4}{tab:campos-obligatorios-evento-dominio}; Subdoc. 4, 4.1.8.3; Subdoc. 3, 3.3.3 y 3.4.1; Formulario T-14, paquete 1.4.5.1. Evidencia: Prueba en un cambio de turno de noche en invierno, con guantes de invierno puestos y sin retirarlos, cronometrando el tiempo entre la lectura de la credencial y la apertura de sesión y verificando que no supere 2 segundos en el 95 \% de los intentos. \\
RNF-DES-005 Propio - No funcional & El sistema debe completar la reconstrucción del origen de un lote ya embarcado en menos… & Cumple & Módulo Trazabilidad y Embarque & Subdoc. 3, 3.2.5 y 3.4.2.2; Formulario T-14, paquete 1.4.3.3. Evidencia: Ejercicio cronometrado sobre tres lotes ya embarcados durante la marcha blanca, midiendo la reconstrucción en menos de 2 horas, y revisión de que cada etapa muestra su evidencia o su supuesto declarado y se distingue una de otro, con verificación por el área comercial. \\
RNF-DES-006 Propio - No funcional & El sistema debe cerrar la conciliación metalúrgica del mes en no más de 3 días hábiles… & Cumple & Módulo Conciliación Metalúrgica & Subdoc. 3, 3.2.5 y 3.4.2.3; Formulario T-14, paquete 1.4.3.4. Evidencia: Medición durante seis cierres consecutivos, contra la línea base de 18 días hábiles. \\
RNF-DES-007 Propio - No funcional & El sistema debe dejar disponible, todos los días, el costo por tonelada movida y por libra… & Cumple & Módulo Indicadores y Centro Integrado (costo por tonelada, Etapa 2) & Subdoc. 3, 3.2.2.3 y 3.4.3; Formulario T-14, paquete 1.4.4.6. Evidencia: Medición durante 30 días consecutivos, contra la situación actual en que el costo no existe con esa frecuencia y llega 18 días hábiles después del cierre, verificando que cada día esté disponible dentro de las 4 horas desde la entrega de SAP. \\
RNF-DES-008 Propio - No funcional & El sistema debe dejar disponible el cumplimiento del plan minero, comparado contra la versión vigente… & Cumple & Módulo Operación Mina & Subdoc. 3, 3.2.5 y 3.4.1; Formulario T-14, paquete 1.4.3.2. Evidencia: Medición durante 30 turnos consecutivos, contra la línea base de 10 a 15 días de desfase. \\
RNF-DES-009 Propio - No funcional & El sistema debe clasificar automáticamente las detenciones a medida que ocurren y dejar el resto disponible… & Cumple & Módulo Operación Mina y Equipos y Mantenimiento & Subdoc. 3, 3.2.5 y 3.4.2.1; Formulario T-14, paquetes 1.4.3.2 y 1.4.4.1. Evidencia: Medición del retardo entre fin de detención y disponibilidad de la clasificación, y del porcentaje mensual sin clasificar, contra la línea base de 23 \%. Prueba previa a producción (SUP-05): al reclasificar tres meses de historia, un residuo automático superior al 5 \% indica que las reglas de clasificación están mal calibradas y deben ajustarse antes de la puesta en producción. \\
RNF-DES-010 Propio - No funcional & El sistema debe mostrar los indicadores con un retraso máximo, respecto del hecho registrado, de 15… & Cumple & Modelo analítico de Indicadores y Centro Integrado (ClickHouse alimentado por eventos, separado del registro) y caché de consultas de portales & Subdoc. 4, 4.1.2.3 (ADR-08) y 4.1.6.1.1; Subdoc. 5, 5.2.1 y 5.2.2; Subdoc. 5, 5.4.4; Subdoc. 5, Anexo A.3 (latencia\_comprometida); Formulario T-14, paquetes 1.4.3.6, 1.4.4.1 y 1.4.4.6. Evidencia: Comparación entre marca de tiempo del evento y su disponibilidad en la capa analítica, sobre una muestra de 30 días. \\
RNF-DES-011 Propio - No funcional & El sistema debe cargar la página de la vista única de equipo en no más… & Cumple & Vista materializada de Equipos y Mantenimiento servida por el Portal Web SITCO, sin consulta en línea a las fuentes & Subdoc. 4, 4.1.7.2, Tabla \refext{S4}{tab:campos-obligatorios-evento-dominio}; Subdoc. 4, 4.1.2.4 y 4.1.6.2; Subdoc. 3, 3.4.2.4; Subdoc. 9, 9.2.4.5; Formulario T-14, paquete 1.4.4.1. Evidencia: Medición sobre 200 consultas en operación normal y 50 con una integración de telemetría caída. \\
RNF-DES-012 Propio - No funcional & El sistema debe completar en un máximo de 4 horas la sincronización posterior a 48 horas… & Cumple & Nodo de Faena (sala técnica, Etapa 1) & Subdoc. 3, 3.3.6 y 3.2.5; Formulario T-14, paquetes 1.4.2.2 y 1.5.5.2. Evidencia: Corte controlado de 48 horas con cambio de turno dentro de la ventana, midiendo el tiempo hasta la consistencia completa. \\
RNF-DES-013 Propio - No funcional & El sistema debe dejar disponible el cálculo mensual de la huella de carbono por tonelada… & Cumple & Módulo Cumplimiento Ambiental (huella de carbono mensual) & Subdoc. 3, 3.2.5 y 3.4.2.7; Formulario T-14, paquetes 1.4.4.3 y 1.8.7.3. Evidencia: Medición durante tres cierres consecutivos, contra la línea base de un cálculo anual de dos meses de duración. \\
RNF-DES-014 Propio - No funcional & El sistema debe sostener la diferencia de la conciliación metalúrgica mensual en no más de 3… & Cumple & Módulo Conciliación Metalúrgica & Subdoc. 3, 3.2.5; Formulario T-14, paquetes 1.4.3.4 y 1.8.7.1. Evidencia: Medición durante doce cierres consecutivos, contra la línea base de 3,1 \% a 6,8 \% declarada en el numeral 7.1. \\
RNF-DES-015 Propio - No funcional & El flujo de acreditación soportado por el sistema debe permitir habilitar a una persona contratista nueva… & Cumple & Módulo Acreditación y Acceso y Portal de Contratistas & Subdoc. 3, 3.2.5 y 3.4.2.5; Formulario T-14, paquete 1.4.3.5. Evidencia: Medición sobre 50 acreditaciones consecutivas de empresas contratistas de distinto tamaño, contra la línea base de 9 días. \\
RNF-DES-016 Propio - No funcional & El sistema debe registrar el aviso de cambio de material sin detener la pala: 0 segundos… & Cumple & Botonera de Pala (reemplaza la llamada por radio) & Subdoc. 3, 3.3.3; Formulario T-14, paquetes 1.4.5.1 y 1.6.3.5. Evidencia: Medición del tiempo de ciclo de carguío (el tiempo de una palada) con y sin el aviso habilitado, durante la marcha blanca, en turno de día y de noche, sobre al menos 500 ciclos. \\
RNF-DES-017 Propio - No funcional & El sistema debe mantener los tiempos de registro en terreno, de báscula y de consulta… & Cumple & Separación entre el registro transaccional (PostgreSQL) y el modelo analítico (ClickHouse), con prueba de carga concurrente & Subdoc. 5, 5.2.2 y 5.4.6; Subdoc. 4, 4.1.2.3, ADR-08; Subdoc. 9, 9.2.4.5; Formulario T-14, paquetes 1.4.3.6 y 1.5.3.1. Evidencia: Prueba de carga concurrente de operación y analítica en Preproducción. \\
RNF-DES-018 Propio - No funcional & El sistema debe enviar las notificaciones en segundo plano, con reintento si falla, sin duplicados… & Cumple & Mensajería de notificaciones del módulo Administración y Servicios Comunes (envío asíncrono con reintento y registro de entrega) & Subdoc. 4, 4.1.7.1, Tabla \refext{S4}{tab:servicios-alcanzables-desde-fuera-red-compania}; Subdoc. 4, 4.1.6.4, Tabla \refext{S4}{tab:llamadas-sincronas-entre-componentes-umbral}; Subdoc. 5, Anexo A.4, Notificación y entrega (estado e historial\_entregas con intento, entrega, apertura o error); Formulario T-14, paquete 1.4.3.1. Evidencia: Prueba de envío con caída simulada del proveedor de un canal. \\
RNF-DIS-001 Propio - No funcional & El sistema debe operar y registrar en faena durante un corte completo del enlace con Antofagasta… & Cumple & Nodo de Faena (sala técnica, Etapa 1) & Subdoc. 3, 3.3.6; Formulario T-14, paquetes 1.4.2.2 y 1.5.5.2. Evidencia: Corte controlado de 48 horas ejecutado en faena, verificando que ninguna de esas funciones se degrada ni se interrumpe y que ningún registro se pierde ni se rechaza. \\
RNF-DIS-002 Propio - No funcional & El sistema debe sostener la disponibilidad que corresponde según la clasificación de servicios del Artículo 78.2… & Cumple & Servicios críticos y de gestión clasificados por severidad & Subdoc. 3, 3.2.2.6; Formulario T-14, paquetes 1.6.3.1 y 1.8.2.1. Evidencia: Medición mensual con el registro de indisponibilidad del sistema, excluyendo las ventanas de intervención aprobadas. \\
RNF-DIS-003 Propio - No funcional & El sistema no debe perder ningún registro generado en faena durante un corte de enlace: 0… & Cumple & Dispositivos de terreno y Nodo de Faena & Subdoc. 3, 3.3.6; Formulario T-14, paquetes 1.4.2.2 y 1.5.5.2. Evidencia: Conciliación entre los registros generados localmente y los registros presentes tras la sincronización, en tres cortes controlados. \\
RNF-DIS-004 Propio - No funcional & El Adjudicatario debe intervenir el sistema solo dentro de la ventana operacional protegida de la faena… & Cumple & Ventana operacional protegida & Subdoc. 3, 3.2.3.3; Subdoc. 4, 4.1.3.4; Formulario T-14, paquetes 1.1.5.1, 1.6.4.1 y 1.6.4.2. Evidencia: Revisión del registro de intervenciones del periodo contra el calendario operacional de la faena, verificando la ventana y el aviso previo de 10 días hábiles. \\
RNF-DIS-005 Propio - No funcional & El sistema debe mantener operativa la vista única de equipo cuando una o más integraciones… & Cumple & Vista única del equipo (Equipos y Mantenimiento) que indica la fuente de telemetría ausente y sigue disponible & Subdoc. 4, 4.1.7.1, Tabla \refext{S4}{tab:conectores-integrador-sistemas-etapa-lugar-ejecucion}; Subdoc. 4, 4.1.7.2, Tabla \refext{S4}{tab:campos-obligatorios-evento-dominio}; Subdoc. 3, 3.4.2.4; Formulario T-14, paquete 1.4.4.1. Evidencia: Prueba de caída controlada de cada una de las tres integraciones de fabricante, por separado y en conjunto, midiendo el tiempo de presentación de la vista completa contra el umbral en modo degradado de RNF-DES-011 y verificando que la vista indica la fuente que falta y que las demás funciones siguen accesibles. \\
RNF-DIS-006 Propio - No funcional & El sistema debe cumplir para los componentes de terreno un RTO no superior a 1 hora… & Cumple & Nodo de Faena (tres servidores en la sala técnica, bandeja de salida y registro firmado en el dispositivo), ante la falla de un nodo o gabinete. La pérdida total de la sala técnica se trata en RNF-OPE-005. & Subdoc. 5, 5.2.1, Tabla \refext{S5}{tab:objetivos-recuperacion-mecanismo-datos-sostiene}; Subdoc. 4, 4.2.3.4; Subdoc. 9, 9.2.4.5; Formulario T-14, paquetes 1.4.2.2 y 1.5.5.2. Evidencia: Prueba de recuperación de un nodo de borde con corte súbito, cronometrando el tiempo hasta la recuperación del servicio (máximo 1 hora) y verificando que ningún evento del registro operacional se pierde. \\
RNF-DIS-007 Propio - No funcional & El sistema debe degradar de forma controlada ante la indisponibilidad de un componente no crítico… & Cumple parcialmente & Conectores con cortacircuito, cola y límite de tasa; modo reducido del Nodo de Faena informado al usuario & Subdoc. 4, 4.1.2.4; Subdoc. 4, 4.1.7.1; Subdoc. 4, 4.1.9.1, Tabla \refext{S4}{tab:reglas-resolucion-conflictos-sincronizacion}; Subdoc. 3, 3.3.6, Tabla \refext{S3}{tab:funciones-sitco-aranda-durante-corte-enlace}; Subdoc. 9, 9.1.2, Tabla \refext{S9}{tab:caracteristicas-iso-iec-25010-2023-umbral} y \refext{S9}{tab:umbrales-bloqueantes-calidad-codigo}.4.5; Formulario T-14, paquetes 1.4.2.1 y 1.6.1.1. Evidencia: Inyección controlada de fallas y prueba de estrés hasta el punto de quiebre en Preproducción, verificando cada transacción enviada contra las registradas. \\
RNF-DIS-008 Propio - No funcional & El sistema debe permitir desplegar cambios sin interrupción del servicio: 0 minutos de indisponibilidad… & Cumple & Estrategia de despliegue sin interrupción: azul-verde en la nube, extracción en el Nodo de Faena y anillos en terreno & Subdoc. 4, 4.1.3.4 y 4.1.2.2; Formulario T-14, paquete 1.4.1.2. Evidencia: Despliegue demostrado en Preproducción antes de cada paso a producción, con medición de disponibilidad durante el despliegue. \\
RNF-DIS-009 Propio - No funcional & El sistema debe mantener los tiempos de registro en terreno, báscula y portería cuando los portales… & Cumple & Capa de borde y exposición: los canales de terreno no pasan por el borde público, por lo que sus tiempos no dependen de la carga de los portales & Subdoc. 4, 4.1.4.2; Subdoc. 9, 9.2.4.5; Formulario T-14, paquete 1.5.3.1. Evidencia: Prueba de carga sobre los portales en Preproducción con medición simultánea de los tiempos de terreno. \\
RNF-INF-001 Propio - No funcional & El Proponente debe especificar, y el sistema debe operar sobre, 1 sala técnica principal en faena… & Cumple & Sala técnica nueva, Nodo de Faena y cuatro gabinetes de borde & Subdoc. 3, 3.3.5 y 3.2.2.2; Formulario T-14, paquete 1.6.1.2. Evidencia: Revisión de la arquitectura física contra RT-06.01 y verificación en terreno de cada emplazamiento. \\
RNF-INT-001 Propio - No funcional & El sistema debe leer el historiador de proceso de la planta mediante OPC UA. & Cumple & Integrador de Sistemas (conector del historiador, OPC UA por diodo de datos) & Subdoc. 3, 3.3.4 y 3.3.2.2; Formulario T-14, paquete 1.4.2.1. Evidencia: Revisión de la integración y verificación del protocolo efectivamente empleado. \\
RNF-INT-002 Propio - No funcional & La integración con el sistema de gestión de laboratorio debe usar una interfaz documentada del propio… & Cumple & Integrador de Sistemas (conector de laboratorio) & Subdoc. 3, 3.3.2.2; Formulario T-14, paquete 1.4.2.1. Evidencia: Prueba de ida y vuelta sobre una muestra de 200 muestras: se envía la identificación de cada una al sistema de gestión de laboratorio y se recibe su resultado, contando en cada sentido los registros perdidos y los duplicados (se exige 0 y 0). \\
RNF-INT-003 Propio - No funcional & El sistema debe integrarse en ambos sentidos con el sistema de despacho de flota, cuyo contrato… & Cumple & Integrador de Sistemas (conector de despacho de flota) & Subdoc. 3, 3.3.2.2; Formulario T-14, paquete 1.4.2.1. Evidencia: Revisión funcional conjunta con el proveedor del sistema de despacho. \\
RNF-INT-004 Propio - No funcional & El sistema debe entregar datos a SAP sin reemplazarlo ni modificarlo, con 0 emisiones contables fuera… & Cumple & Integrador de Sistemas (conector de SAP) & Subdoc. 3, 3.3.2.2 y 3.4.3; Formulario T-14, paquete 1.4.2.1. Evidencia: Revisión de las interfaces con SAP y conciliación entre el dato entregado y el dato contabilizado. \\
RNF-INT-005 Propio - No funcional & El sistema debe integrar operación y gestión con un modelo de referencia estándar de la industria… & Cumple & Integrador de Sistemas (ISA-95) y conector de telemetría (ISO 15143-3) & Subdoc. 3, 3.4.3 y 3.4.2.4; Formulario T-14, paquetes 1.2.4.1 y 1.3.1.3. Evidencia: Revisión del modelo de integración y del mapeo de señales por fabricante. \\
RNF-INT-006 Propio - No funcional & El sistema debe documentar el 100 \% de sus interfaces, con servicios síncronos en OpenAPI 3.1… & Cumple & Integrador de Sistemas (un conector por sistema) y catálogo de interfaces generado desde el código & Subdoc. 4, 4.1.7.4, 4.1.7.5 y 4.1.7.6; Formulario T-14, paquetes 1.3.5.2 y 1.4.2.1. Evidencia: Revisión del catálogo de interfaces y de los contratos contra el código desplegado en cada entrega: servicios en OpenAPI 3.1, flujos por eventos en AsyncAPI 2.6 o superior, versionado semántico, pruebas de compatibilidad hacia atrás y registro de avisos de obsolescencia con 6 meses de preaviso. \\
RNF-POR-001 Propio - No funcional & El sistema debe poder replicarse a una nueva unidad operacional mediante parametrización, con catálogos, unidades organizacionales… & Cumple & Módulo Administración y Servicios Comunes (parametrización) & Subdoc. 3, 3.4.3; Formulario T-14, paquete 1.4.3.1. Evidencia: Ejercicio cronometrado de configuración de una unidad operacional adicional en ambiente de prueba, verificando que se completa en el plazo declarado y sin modificar código fuente. \\
RNF-POR-002 Propio - No funcional & El sistema debe permitir exportar el 100 \% de la información del Cliente, incluidos los registros… & Cumple & Exportación total en formatos abiertos por el cliente, sin intervención del adjudicatario (Administración y Servicios Comunes) & Subdoc. 5, 5.2.2; Subdoc. 5, 5.2.5; Subdoc. 9, 9.2.4.5; Formulario T-14, paquetes 1.4.3.1 y 1.4.6.5. Evidencia: Prueba de exportación y lectura por un tercero sin acceso a la solución. \\
RNF-POR-003 Propio - No funcional & El sistema debe distribuir su aplicación móvil por las tiendas oficiales o por gestión de flota… & Cumple parcialmente & App SITCO Terreno distribuida por la gestión de dispositivos (Microsoft Intune) & Subdoc. 4, 4.1.1.3; Subdoc. 4, Anexos, T11-25 (Intune, distribución firmada); Formulario T-14, paquetes 1.4.1.3 y 1.4.5.1. Falta: rechazo automático de las versiones no firmadas. Evidencia: Revisión del proceso de distribución e intento de instalar una versión no firmada. \\
RNF-POR-004 Propio - No funcional & El sistema debe soportar, en su aplicación móvil, la versión vigente del sistema operativo móvil… & Cumple & App SITCO Terreno sobre Windows 11 IoT Enterprise LTSC con actualizaciones mensuales & Subdoc. 4, 4.1.2.3, ADR-06; Subdoc. 4, 4.1.1.4; Formulario T-14, paquete 1.4.5.1. Evidencia: Prueba de la aplicación en cada versión soportada. \\
RNF-SEG-001 Propio - No funcional & El sistema debe cifrar campo por campo los datos de salud ocupacional, exámenes preocupacionales, resultados… & Cumple & Módulo Seguridad y Salud Ocupacional (datos con cifrado de campo y clave separada por categoría) & Subdoc. 4, 4.1.8.4, Tabla \refext{S4}{tab:amenazas-mayor-consecuencia-componentes-expuestos-control}; Subdoc. 5, 5.2.6; Subdoc. 3, 3.4.2.6; Formulario T-14, paquetes 1.3.3.1 y 1.4.4.2. Evidencia: Revisión de arquitectura y prueba de lectura directa de la base de datos, verificando que los campos de las 4 categorías sensibles no revelan su contenido. \\
RNF-SEG-002 Propio - No funcional & El sistema debe registrar el 100 \% de las consultas a información sensible, además de toda… & Cumple & Módulo Seguridad y Salud Ocupacional (registro de cada consulta) & Subdoc. 3, 3.4.2.6; Formulario T-14, paquetes 1.3.2.3 y 1.4.4.2. Evidencia: Auditoría sobre una muestra de accesos, contrastando el registro con los accesos efectivamente realizados. \\
RNF-SEG-003 Propio - No funcional & El sistema debe usar las cuentas del directorio corporativo del grupo para el personal propio… & Cumple & Identidad federada con directorio corporativo y doble factor; Módulo Administración y Servicios Comunes & Subdoc. 3, 3.2.3.3 y 3.4.3; Formulario T-14, paquetes 1.3.3.3 y 1.4.5.2. Evidencia: Revisión del registro de accesos y de la configuración de autenticación; prueba de intento de uso de credencial compartida; prueba de acceso administrativo, privilegiado y externo sin segundo factor; y prueba de alta y baja de una persona contratista desde el portal de empresas. \\
RNF-SEG-004 Propio - No funcional & El sistema debe conectarse al historiador de proceso de la planta como cliente OPC UA… & Cumple & Diodo de datos y zona desmilitarizada industrial (IEC 62443) & Subdoc. 3, 3.3.4 y 3.3.2.2; Formulario T-14, paquetes 1.4.2.1 y 1.6.1.3. Evidencia: Revisión de arquitectura de red y prueba de intento de escritura sobre el control de proceso. \\
RNF-SEG-005 Propio - No funcional & El sistema debe conservar los registros, con integridad verificable y disponibilidad para fiscalización, por los siguientes… & Cumple & Retención por categoría de registro con repositorio de consulta histórica, inmutabilidad y auditoría & Subdoc. 5, 5.2.4, Tabla \refext{S5}{tab:retencion-categoria-entidad}; Subdoc. 5, 5.3.5; Formulario T-14, paquetes 1.3.2.3 y 1.4.6.5. Evidencia: Prueba de recuperación de registros de cada categoría en el extremo de su plazo de retención. \\
RNF-SEG-006 Propio - No funcional & El sistema debe revocar el acceso de una persona en un máximo de 24 horas… & Cumple parcialmente & Módulo Acreditación y Acceso con copia local de identidades en el Nodo de Faena & Subdoc. 4, 4.1.8.1; OVERFORE-Subdocumento4-Anexos, Anexo E (6.5); Formulario T-14, paquetes 1.4.2.2 (copia local de identidades con vigencia de 72 horas) y 1.3.3.3 (baja en 24 horas o menos). Falta: durante un corte la baja se registra en la consola local del Nodo y el Nodo rechaza la credencial al vencer su copia de 72 horas; el plazo de 24 horas no se garantiza en faena. Evidencia: Prueba de revocación cronometrada sobre 20 casos, incluida la propagación a los nodos de borde \\
RNF-SEG-007 Propio - No funcional & El sistema debe autenticar cada integración entre sistemas con credenciales seguras y no con claves fijas… & Cumple & Puerta de exposición y conectores del Integrador de Sistemas (OAuth 2.1 con credenciales de cliente o TLS mutuo) & Subdoc. 4, 4.1.8.2, Tabla \refext{S4}{tab:mecanismo-identidad-cada-poblacion-usuaria}; Subdoc. 4, 4.1.7.4; Formulario T-14, paquete 1.4.2.1. Evidencia: Revisión de la configuración de cada integración (OAuth 2.1 con credenciales de cliente o autenticación mutua TLS, según RT-05.18) y prueba de acceso sin credencial válida. \\
RNF-SEG-008 Propio - No funcional & El sistema debe cifrar el 100 \% de las comunicaciones en tránsito con la versión vigente… & Cumple & Capa de borde (TLS 1.3) y malla de servicios con TLS mutuo & Subdoc. 4, 4.1.8.4, Tabla \refext{S4}{tab:amenazas-mayor-consecuencia-componentes-expuestos-control}; Subdoc. 4, 4.1.4.2; Subdoc. 4, 4.2.2.2; Formulario T-14, paquete 1.6.1.1. Evidencia: Escaneo de configuración de todos los servicios expuestos e internos: TLS 1.3, ningún servicio con TLS 1.0 o 1.1, HSTS con precarga y rotación de certificados (RT-11.08). \\
RNF-SEG-009 Propio - No funcional & El sistema debe cifrar el 100 \% de los datos en reposo, en nube, en faena… & Cumple & Cifrado del 100 \% de los almacenes con claves en HSM (bóveda de claves por región) & Subdoc. 4, 4.1.8.4, Tabla \refext{S4}{tab:amenazas-mayor-consecuencia-componentes-expuestos-control}; Subdoc. 4, 4.1.4.7; Subdoc. 4, Anexos, T11 bóveda de claves; Formulario T-14, paquete 1.3.3.1. Evidencia: Revisión de configuración de almacenamiento y de la política de claves. \\
RNF-SEG-010 Propio - No funcional & El sistema debe proteger el 100 \% de los puntos de entrada públicos de los portales… & Cumple & Entrada pública de la capa de borde (Front Door con WAF, DDoS y protección contra bots) para los portales de contratistas y de compradores & Subdoc. 4, 4.1.4.2; Subdoc. 4, 4.2.2.2; Subdoc. 4, 4.1.8.2, Tabla \refext{S4}{tab:mecanismo-identidad-cada-poblacion-usuaria}; Formulario T-14, paquete 1.6.1.1. Evidencia: Prueba de abuso automatizado y revisión de accesibilidad del reto. \\
RNF-SEG-011 Propio - No funcional & El sistema debe registrar los eventos de seguridad, incluidos los accesos a la sala técnica… & Cumple & Registro y detección de eventos de seguridad (capa 7 con SIEM) con archivo; bitácora inalterable de la sala técnica (RF-ACR-018) & Subdoc. 4, 4.1.8.5; Subdoc. 4, 4.1.4.7 y 4.1.4.8; Subdoc. 3, 3.4.2.5; Subdoc. 2, Anexo A.1 (RF-ACR-018); Formulario T-14, paquetes 1.6.1.1 y 1.8.2.2. Evidencia: Revisión de la retención configurada y recuperación de un evento archivado. \\
RNF-SEG-012 Propio - No funcional & El sistema debe ofrecer inicio de sesión único para todos sus módulos, con 1 autenticación… & Cumple & Módulo Administración y Servicios Comunes (identidad compartida: Entra ID federado y External ID) & Subdoc. 3, 3.4.3; Subdoc. 4, 4.1.4.7 y 4.1.8.3; Formulario T-14, paquete 1.3.3.3. Evidencia: Prueba de navegación entre módulos y de cierre de sesión. \\
RNF-SEG-013 Propio - No funcional & El sistema debe dar permisos de administrador solo por tiempo limitado y con aprobación previa, grabar… & Cumple & Elevación temporal de privilegios con aprobación previa y sesión grabada; consola fuera de banda con sesión grabada & Subdoc. 4, 4.1.8.3; Subdoc. 4, 4.1.8.5, Tabla \refext{S4}{tab:modulo-plataforma-sostiene-cada-funcion-durante}; Subdoc. 4, Anexos, Anexo E 8.2 y T11-36; Formulario T-14, paquete 1.3.3.3. Evidencia: Revisión de cuentas privilegiadas y de grabaciones sobre una muestra de sesiones; revisión de la declaración del procedimiento de acceso de emergencia y verificación del control y la auditoría de esa cuenta. \\
RNF-SEG-014 Propio - No funcional & El sistema debe cerrar la sesión de una persona por duración máxima o por inactividad… & Cumple parcialmente & Identidad y sesiones de los portales y canales (Entra ID, External ID, credencial con PIN en terreno) & Subdoc. 4, 4.1.8.3; Formulario T-14, paquete 1.3.3.3. Falta: duración máxima e inactividad de sesión modificables por el Cliente. Evidencia: Pruebas de caducidad, revocación y manipulación de credenciales de sesión (credenciales firmadas, de vida breve y con refresco rotatorio, según RT-12.07 y RT-12.08). \\
RNF-SEG-015 Propio - No funcional & El sistema debe excluir de sus registros técnicos los datos personales sensibles y las credenciales… & Cumple parcialmente & Capa 8 de observabilidad (OpenTelemetry y Azure Monitor) con trazas técnicas de 90 días sin datos personales sensibles & Subdoc. 4, 4.1.4.8; Subdoc. 5, 5.2.4, Tabla \refext{S5}{tab:retencion-categoria-entidad}; Formulario T-14, paquete 1.6.1.1. Evidencia: Revisión automatizada de muestras de registros técnicos. \\
RNF-SEG-016 Propio - No funcional & El sistema debe cifrar el 100 \% de la información almacenada en los dispositivos de terreno… & Cumple & Dispositivos de terreno con registro local cifrado (SQLCipher y TPM) y gestión remota por Intune & Subdoc. 4, 4.1.5.1; Subdoc. 4, 4.1.8.4, Tabla \refext{S4}{tab:amenazas-mayor-consecuencia-componentes-expuestos-control}; Subdoc. 4, 4.1.8.1; Subdoc. 4, Anexos, Anexo E 7.9; Formulario T-14, paquetes 1.4.5.1 y 1.6.2.1. Evidencia: Prueba de extracción del almacenamiento de un dispositivo y de borrado remoto. \\
RNF-USA-001 Propio - No funcional & El sistema debe permitir operar el 100 \% de las funciones de registro en terreno… & Cumple & Botonera de Pala y tabletas rugerizadas & Subdoc. 3, 3.3.3; Formulario T-14, paquetes 1.4.5.1 y 1.5.2.1. Evidencia: Prueba con personal de operación en turno de noche de invierno, con los guantes efectivamente usados en faena. \\
RNF-USA-002 Propio - No funcional & El sistema debe mantener legible la interfaz de terreno bajo radiación solar directa, de modo… & Cumple & Botonera de Pala y tabletas legibles al sol & Subdoc. 3, 3.3.3; Formulario T-14, paquetes 1.4.5.1 y 1.5.2.1. Evidencia: Prueba en cabina al mediodía en verano, con 10 personas operadoras, midiendo la proporción de registros completados sin sombrear la pantalla con la mano o el cuerpo. \\
RNF-USA-003 Propio - No funcional & El Proponente debe especificar equipamiento de terreno que funcione entre -6 °C y 34 °C… & Cumple & Tableta rugerizada (T11-67) y Botonera de Pala (T11-68) de terreno & Subdoc. 4, 4.1.5.1 (Tabla \refext{S4}{tab:conectores-integrador-sistemas-etapa-lugar-ejecucion}); OVERFORE-Subdocumento4-Anexos, Anexo A (T11-67 y T11-68); Formulario T-11, T11-67 y T11-68; Formulario T-14, paquetes 1.5.2.1 y 1.6.2.1. Evidencia: Prueba en cabina de pala y de camión, con el equipamiento especificado, durante un turno completo de invierno y uno de verano. \\
RNF-USA-004 Propio - No funcional & El sistema debe permitir al supervisor de turno clasificar el residuo de detenciones no resueltas automáticamente… & Cumple & Módulo Operación Mina (taxonomía única; el supervisor clasifica solo el residuo no resuelto por reglas) & Subdoc. 3, Tablas \refext{S3}{tab:criterios-aceptacion-etapa-1} y \refext{S3}{tab:criterios-aceptacion-etapa-2}; Subdoc. 3, 3.4.2.1; Subdoc. 13, 13.2; Formulario T-14, paquetes 1.2.2.5 y 1.4.3.2. Evidencia: Medición del tiempo efectivo de clasificación durante 30 turnos, con el volumen real de detenciones. \\
RNF-USA-005 Propio - No funcional & El sistema debe lograr que la persona operadora de pala registre por su cuenta… & Cumple & Botonera de Pala & Subdoc. 3, 3.2.5; Formulario T-14, paquetes 1.6.3.1 y 1.6.6.4. Evidencia: Medición en las cuatro últimas semanas de la marcha blanca de la Etapa 1, contra la línea base de dos iniciativas previas abandonadas. \\
RNF-USA-006 Propio - No funcional & El sistema debe presentar en español el 100 \% de sus interfaces, mensajes, alertas y documentación… & Cumple & Interfaces de SITCO, documentación y mesa de ayuda en español & Subdoc. 4, 4.1.4.1; Subdoc. 3, 3.2.2.6; Subdoc. 9, Anexos, puerta de calidad de documentación; Formulario T-14, paquete 1.3.4.2. Evidencia: Revisión de la totalidad de las pantallas y mensajes en cada entrega. \\
RNF-USA-007 Propio - No funcional & El sistema debe permitir completar un registro de seguridad en terreno, con firma, sin conexión… & Cumple & App SITCO Terreno (módulo Seguridad y Salud Ocupacional) & Subdoc. 3, 3.4.2.6; Formulario T-14, paquetes 1.4.4.2 y 1.5.2.2. Evidencia: Medición con inspectores de seguridad en terreno, en zona sin cobertura, sobre 50 registros. \\
RNF-USA-008 Propio - No funcional & El sistema no debe requerir de la persona operadora de pala ninguna acción de registro distinta… & Cumple & Botonera de Pala (una sola acción de registro) & Subdoc. 3, 3.3.3 y 3.4.2.1; Formulario T-14, paquete 1.4.5.1. Evidencia: Revisión funcional de la interfaz de cabina y medición del tiempo de ciclo durante la marcha blanca. \\
RNF-USA-009 Propio - No funcional & El sistema no debe requerir de la persona operadora de equipo ninguna acción de clasificación… & Cumple & Módulo Operación Mina (clasifica el supervisor de turno) & Subdoc. 3, Tablas \refext{S3}{tab:criterios-aceptacion-etapa-1} y \refext{S3}{tab:criterios-aceptacion-etapa-2}; Subdoc. 3, 3.4.2.1; Formulario T-14, paquetes 1.4.3.2 y 1.7.2.1. Evidencia: Revisión funcional de las interfaces de cabina y de terreno durante la marcha blanca. \\
RNF-USA-010 Propio - No funcional & El sistema debe cumplir WCAG 2.2 nivel AA en todas sus interfaces destinadas a personas usuarias… & Cumple & Capa de presentación de los seis canales con puerta de calidad de accesibilidad WCAG 2.2 AA & Subdoc. 4, 4.1.4.1; Subdoc. 9, 9.1.2, Tabla \refext{S9}{tab:caracteristicas-iso-iec-25010-2023-umbral} y \refext{S9}{tab:umbrales-bloqueantes-calidad-codigo}.4.3; Formulario T-14, paquetes 1.3.4.2, 1.4.5.2 y 1.5.1.3. Evidencia: Verificación con herramientas automatizadas y pruebas manuales, con informe de conformidad. \\
RNF-USA-011 Propio - No funcional & El sistema debe funcionar con diseño responsivo en escritorio, tableta y teléfono, adaptando el contenido… & Cumple & Portales web y consolas del sistema de diseño de SITCO & Prototipo interactivo UX/UI (Cap. 25); Formulario T-14, paquete 1.4.5.2. Evidencia: Revisión de cada pantalla en las cuatro resoluciones. \\
RNF-USA-012 Propio - No funcional & El sistema debe permitir alcanzar cada funcionalidad principal desde la pantalla de inicio del perfil… & Cumple & App SITCO Terreno, Consola de Portería y portales con máximo de tres interacciones por función & Subdoc. 4, 4.1.5.1; Subdoc. 9, 9.1.2, Tabla \refext{S9}{tab:caracteristicas-iso-iec-25010-2023-umbral}; Formulario T-14, paquetes 1.3.4.2 y 1.4.5.1. Evidencia: Recorrido de cada funcionalidad principal por perfil, contando interacciones. \\
RNF-USA-013 Propio - No funcional & El sistema debe entregar retroalimentación visual ante cada acción y presentar los errores indicando qué ocurrió… & Cumple & App SITCO Terreno y Consola de Portería con confirmación visual de cada acción & Subdoc. 4, 4.1.5.1; Subdoc. 4, 4.1.4.3; Formulario T-14, paquete 1.3.4.2. Evidencia: Revisión del catálogo de mensajes y prueba de usabilidad con personas usuarias del Cliente. \\
RNF-USA-014 Propio - No funcional & El sistema debe soportar a personas usuarias con baja alfabetización digital mediante alto contraste, objetivos táctiles… & Cumple & Botonera de Pala (botones grandes, sin pantalla) y tabletas con guantes & Subdoc. 3, 3.3.3; Formulario T-14, paquete 1.3.4.2. Evidencia: Revisión de interfaces de terreno y prueba con personas operadoras. \\
RNF-USA-015 Propio - No funcional & El sistema debe ser navegable íntegramente por teclado, con el 100 \% de las funciones… & Cumple & Navegación completa por teclado en las pantallas de escritorio, con orden de foco lógico y atajos para las operaciones frecuentes & Subdoc. 9, 9.2.4.3; Formulario T-14, paquetes 1.3.4.2 y 1.5.1.3. Evidencia: Prueba manual de navegación por teclado en cada pantalla de escritorio. \\
RNF-OPE-001 Propio - No funcional & El sistema debe mostrar en una sola consola el estado de sus componentes, de sus integraciones… & Cumple & Centro de operaciones 24x7 & Subdoc. 3, 3.2.2.5; Formulario T-14, paquetes 1.6.1.1 y 1.8.2.1. Evidencia: Prueba de diagnóstico de una falla de integración inducida, cronometrada, ejecutada por personal de TI del Cliente desde la consola única, con un tiempo de diagnóstico no superior a 30 minutos. \\
RNF-OPE-002 Propio - No funcional & El sistema debe emitir alertas anticipadas de degradación al equipo de operación: colas de sincronización cerca… & Cumple parcialmente & Centro de operaciones 24x7 (alertas antes de que un recurso llegue a su límite) & Subdoc. 3, 3.2.2.5; Formulario T-14, paquetes 1.4.2.2, 1.4.3.6 y 1.4.4.6. Falta: emisión de la alerta en no más de 5 minutos desde que la cola cruza 120.000 eventos; alerta de integraciones de fabricante caídas; alerta de indicadores atrasados sobre sus umbrales (15 minutos de turno y 4 horas de gestión). Evidencia: Prueba de inyección controlada de degradación (crecimiento de cola, caída de integración) verificando el tiempo de emisión de la alerta y el margen respecto del umbral comprometido. Se verifica que la alerta se emite al cruzar 120.000 eventos, en no más de 5 minutos, y que el margen restante equivale a al menos 4 horas de acumulación al ritmo medio. \\
RNF-OPE-003 Propio - No funcional & El sistema debe contar con una política de respaldo de 3 copias en 2 medios distintos… & Cumple parcialmente & Respaldo 3-2-1-1-0 con copia inmutable en Brazil South y restauración mensual verificada & Subdoc. 4, 4.2.3.5; OVERFORE-Subdocumento4-Anexos, Anexo E (E.4, control 8.13); OVERFORE-Subdocumento4-Anexos, Anexo A (T11-20, T11-46); Formulario T-14, paquetes 1.3.1.2 y 1.8.8.5. Falta: declaración por dominio de datos para laboratorio y para mantenimiento, con su tiempo de restauración propio; 0 errores en las verificaciones de restauración. Evidencia: Revisión de la política y de la configuración de respaldo contra el esquema 3-2-1-1-0 y de los resultados de verificación de restauración; revisión del cifrado y de la gestión de claves; intento de borrado o modificación de una copia inmutable con credenciales administrativas; revisión de la declaración por dominio de datos. \\
RNF-OPE-004 Propio - No funcional & El sistema debe permitir probar la restauración de los respaldos al menos una vez al mes… & Cumple & Pruebas mensuales de restauración de respaldos en Preproducción, con el tiempo efectivo medido en un registro mensual & Subdoc. 4, 4.2.3.5 (verificación mensual de la restauración en Preproducción); Formulario T-14, paquete 1.8.8.5. Evidencia: Ejecución y acta de la prueba mensual de restauración, con el tiempo medido contra el tiempo de restauración declarado en RNF-OPE-003. \\
RNF-OPE-005 Propio - No funcional & El sistema debe recuperarse ante una falla mayor del sitio central con un tiempo máximo… & Cumple parcialmente & Región secundaria São Paulo con réplica continua & Subdoc. 3, 3.2.2.5 y 3.3.5; Subdoc. 4, 4.3.2.2 y 4.2.3.4; Formulario T-14, paquetes 1.5.5.1 (conmutación real a la región secundaria en M15 y M21, antes de cada paso), 1.3.1.4 (RTO de 1 hora ante la falla de un servidor o de un gabinete) y 1.8.8.1 (dos conmutaciones reales al año). Se distinguen tres escenarios: sitio central, RTO máximo de 4 horas y RPO máximo de 15 minutos; falla de un nodo o gabinete, RTO máximo de 1 hora y cero eventos perdidos; pérdida total de sala, RTO máximo de 4 horas y cero eventos confirmados perdidos, conforme a S8 G.4. Falta: completar el inventario del kit, la custodia exterior de eventos y su trazabilidad a T-13/T-17. Evidencia: Prueba de conmutación real con informe del RTO y RPO efectivamente alcanzados y plan de corrección de brechas \\
RNF-OPE-006 Propio - No funcional & El sistema debe permitir seguir cualquier operación de punta a punta por todos los sistemas… & Cumple & Integrador de Sistemas (un conector por sistema) (identificador común de mensaje) & Subdoc. 3, 3.4.3; Subdoc. 4, 4.1.4.8 (Figura \refext{S4}{fig:flujo-telemetria-capa-observabilidad-transversal}); Subdoc. 4, 4.1.7.1; Formulario T-14, paquetes 1.6.1.1 y 1.8.7.1. Evidencia: Seguimiento de una operación de negocio de punta a punta a través de todos los sistemas involucrados, en la nube y en faena, y verificación, sobre una muestra de transacciones e integraciones, de que todas portan el identificador único y registran su entrada y su salida, y de que el monitoreo del componente on-premise aparece en la misma plataforma de observabilidad que el de la nube. Revisión de la definición y del origen de los datos de cada indicador de nivel de servicio. \\
RNF-TRV-02.06 Propio - No funcional & El sistema debe garantizar que el 100 \% de las operaciones de escritura expuestas a reintentos… & Cumple parcialmente & Clave de idempotencia exigida en toda escritura reintentable y consumidores idempotentes & Subdoc. 4, 4.1.7.2 (Tabla \refext{S4}{tab:campos-obligatorios-evento-dominio}); Subdoc. 4, 4.1.2.4; Formulario T-14, paquete 1.4.2.1. Falta: prueba de reenvío de la misma solicitud a las 72 horas, dentro de la ventana de deduplicación. Evidencia: Prueba: reenviar cada operación de escritura expuesta a reintentos con la misma clave y verificar 0 efectos duplicados; inspección de la documentación de la ventana de deduplicación, incluido un reenvío a las 72 horas. \\
RNF-TRV-02.07 Propio - No funcional & El sistema debe entregar cada evento de negocio al menos una vez, con 0 duplicados… & Cumple & Bus de eventos (Azure Service Bus Premium) con entrega al menos una vez, orden por agregado y deduplicación en el consumidor & Subdoc. 4, 4.1.4.5 (Figura \refext{S4}{fig:componentes-capa-integracion-eventos-nube-faena}); Subdoc. 4, 4.1.7.3; Subdoc. 4, 4.1.9.2 (Tabla \refext{S4}{tab:reglas-resolucion-conflictos-sincronizacion}); Formulario T-14, paquete 1.4.2.1. Evidencia: Diseño del bus de eventos y de sus consumidores descrito en S4 (entrega al menos una vez, orden por agregado, cola de fallidos y deduplicación en el consumidor), que se verifica con la prueba: inyectar fallas y reenvíos en el bus de eventos y verificar 0 eventos perdidos, 0 efectos duplicados en el consumidor y orden por partición/agregado. \\
RNF-TRV-02.08 Propio - No funcional & El sistema debe aplicar en el 100 \% de las llamadas remotas un tiempo de espera… & Cumple parcialmente & Conectores del Integrador de Sistemas con tiempo de espera, reintento, cortacircuito y mamparo & Subdoc. 4, 4.1.7.1; Subdoc. 4, 4.1.2.4; OVERFORE-Subdocumento4-Anexos, Anexo A (T11-07); Formulario T-14, paquete 1.6.1.1. Falta: ráfaga de 100 llamadas, y las cifras de reintento (base de 200 milisegundos, variación de ±20 \%) y de límite (50 llamadas por segundo) para las llamadas servicio a servicio. Evidencia: Inspección de configuración (1 s en consultas, 2 s en escrituras, 3 reintentos, 50 llamadas por segundo) de todas las llamadas remotas y prueba de falla inyectada de cada dependencia verificando el disparo del cortacircuitos, el aislamiento y el retroceso. \\
RNF-TRV-04.11 Propio - No funcional & El Adjudicatario debe mantener una cobertura de pruebas automatizadas de al menos 70 \%… & Cumple & Pruebas unitarias y de contrato automatizadas que se ejecutan en cada solicitud de integración, con cobertura de líneas bloqueante en 70 \% & Subdoc. 9, 9.1.4 (Tabla \refext{S9}{tab:declaracion-uso-ia-seccion}, cobertura de líneas bloqueante en 70 \%); Formulario T-14, paquetes 1.1.4.2 y 1.5.1.1. Evidencia: Reporte de cobertura emitido por la herramienta de pruebas sobre el código de lógica de negocio en cada corrida de integración, con resultado igual o superior a 70 \%, y prueba de que un cambio bajo el umbral es rechazado. \\
RNF-TRV-07.03 Propio - No funcional & El sistema debe replicar de forma continua los datos del sitio principal al sitio secundario, medir… & Cumple & Réplica continua a Brazil South con monitor de retraso y alerta a los 5 minutos & Subdoc. 4, 4.3.2.1 (Figura \refext{S4}{fig:recuperacion-brazil-south}); OVERFORE-Subdocumento4-Anexos, Anexo A (T11-12); Formulario T-14, paquete 1.3.1.4. Evidencia: Prueba en Preproducción y en el sitio secundario: inducir retraso de replicación por sobre el umbral de 5 minutos y verificar que la medición se registra y la alerta llega al equipo de operación; revisión de la métrica de retraso en la consola operativa durante la marcha blanca. \\
RNF-TRV-11.04 Propio - No funcional & El Adjudicatario debe mantener un programa de gestión de vulnerabilidades con escaneo continuo del sistema… & Cumple & Informe mensual del escáner de vulnerabilidades & Subdoc. 3, 3.2.2.5; Subdoc. 4, 4.1.7.6; Subdoc. 4, 4.1.1.5 (Tabla \refext{S4}{tab:version-fin-soporte-plan-actualizacion-componentes}); OVERFORE-Subdocumento4-Anexos, Anexo E (E.4, control 8.8); Formulario T-14, paquetes 1.5.1.3, 1.5.4.3 y 1.8.2.2. Evidencia: Informe mensual del escáner con fecha de detección y de cierre por vulnerabilidad, contrastado por muestreo con la evidencia de remediación; auditoría del Cliente sobre el registro de vulnerabilidades. \\
RNF-TRV-11.07 Propio - No funcional & El sistema debe publicar el 100 \% de sus servicios alcanzables desde fuera de la red… & Cumple parcialmente & Entrada pública con CDN, WAF y DDoS en capas 3, 4 y 7; acceso de confianza cero para el personal propio & Subdoc. 4, 4.1.4.2 (Figura \refext{S4}{fig:entradas-capa-borde-exposicion}); Subdoc. 4, 4.1.8.2 (Tabla \refext{S4}{tab:mecanismo-identidad-cada-poblacion-usuaria}); Subdoc. 4, 4.2.2.2 (Figura \refext{S4}{fig:actores-caminos-acceso-sobre-topologia-fisica}); Formulario T-14, paquete 1.6.1.1. Falta: que el acceso de confianza cero del Portal Web SITCO (Microsoft Entra Private Access) pase también por CDN, WAF y DDoS en capas 3, 4 y 7. Evidencia: Escaneo externo de dominios, puertos y servicios contrastado con la superficie de exposición declarada (RT-11.13); pruebas de ataque controlado en capas 3, 4 y 7 y de evasión de la capa de borde, con 0 accesos directos al origen. \\
RNF-TRV-11.11 Propio - No funcional & El sistema debe aplicar, en la puerta de enlace de servicios y sobre el 100 \%… & Cumple parcialmente & Puertas de enlace de exposición y operacional con autenticación, autorización, cuotas, límites de tasa, validación de esquema y carga útil sobre el 100 \% de las llamadas & Subdoc. 4, 4.1.4.3 (Figura \refext{S4}{fig:consumidores-controles-puerta-exposicion-puerta-operacional}; Tabla \refext{S4}{tab:comparacion-aplicacion-web-progresiva-nativa-hibrida}); OVERFORE-Subdocumento4-Anexos, Anexo A (T11-07); Formulario T-14, paquete 1.6.1.1. Falta: cuota de 200.000 llamadas diarias y límite de 50 llamadas por segundo, con ráfaga de 100, por cliente autenticado. Evidencia: Pruebas de casos negativos por servicio (sin credencial, sin permiso, sobre la cuota, sobre el límite de 50 llamadas por segundo, sobre la cuota diaria de 200.000, esquema inválido, carga útil maliciosa) y revisión de la configuración de la puerta de enlace, con 0 llamadas no controladas. \\
RNF-TRV-21.11 Propio - No funcional & El sistema debe proveer capacitación en línea en modalidad de autoformación para las personas usuarias… & Cumple & Autoformación en línea por perfil, con registro de avance por persona consultable por el Cliente y disponibilidad mensual de 99,5 \% & Formulario T-14, paquete 1.6.5.2 (autoformación en línea por perfil, avance por persona consultable por el Cliente y disponibilidad mensual de 99,5 \%). Evidencia: Inspección de demostración: completar un curso de autoformación con un usuario de prueba y verificar que el avance queda registrado y es consultable por el Cliente, y medición de la disponibilidad mensual contra 99,5 \% y de la cobertura de cursos por perfil. \\
RNF-TRV-21.13 Propio - No funcional & El sistema de mesa de ayuda del Adjudicatario debe exponer al Cliente los indicadores clave… & Cumple & Informe mensual de nivel de servicio & Subdoc. 3, 3.2.2.5; Formulario T-14, paquetes 1.8.1.1 y 1.8.7.1. Evidencia: Inspección de demostración: verificar visibilidad de indicadores clave con perfil Cliente y restricción de los específicos por perfil; verificar la emisión del reporte mensual por severidad con análisis de incumplimientos. \\
RNF-TRV-21.15 Propio - No funcional & El sistema de mesa de ayuda del Adjudicatario debe proveer un canal único de registro… & Cumple & Mesa de ayuda de nivel 1 & Subdoc. 3, 3.2.2.5; Formulario T-14, paquetes 1.8.1.1 y 1.8.3.1. Evidencia: Prueba de aceptación: crear casos de cada severidad por el canal y verificar número único, clasificación y estados del ciclo hasta el cierre; verificar que un ticket no pasa a cerrado sin confirmación o sin vencer el plazo de 7 días corridos declarado; consultar el registro histórico y el análisis de tendencia mensual. \\
\end{tabla}
\begin{tabla}{I{0.134}I{0.300}I{0.109}I{0.150}I{0.308}}{Requerimientos de transición y preparación.}{tab:requerimientos-transicion-preparacion}
  {\cab{ID requerimiento} & \cab{Descripción} & \cab{Estado} & \cab{Componente que lo satisface} & \cab{Sección de la propuesta y evidencia}}
  {\emph{Fuente:} elaboración propia a partir de las Bases técnicas del caso Minería.}
TRA-001 Propio - Transición y preparación & El sistema debe registrar el tiempo que cada punto de muestreo obligatorio agrega al ciclo productivo… & Cumple & Módulo Conciliación Metalúrgica (muestreo en la planta) & Subdoc. 3, 3.4.5; Formulario T-14, paquete 1.6.3.1. Evidencia: Medición cronometrada en la marcha blanca del tiempo agregado por cada punto de muestreo, con y sin el paso, informada por punto de muestreo y por turno. Se verifica que el tiempo agregado por cada punto no supera 2 segundos en el percentil 95. \\
TRA-002 Propio - Transición y preparación & El sistema debe contener documentados, antes del cierre de la marcha blanca de la Etapa 1… & Cumple & Módulo Conciliación Metalúrgica (factores documentados y versionados) & Subdoc. 3, 3.4.2.3; Formulario T-14, paquetes 1.2.6.1 y 1.4.3.4. Evidencia: Recálculo de doce cierres mensuales históricos con los factores reconstruidos, contrastando el resultado contra el cierre original de cada mes. \\
TRA-003 Propio - Transición y preparación & El Adjudicatario debe capacitar a la totalidad del personal de cada área, programando la capacitación… & Cumple & Capacitación por ola & Subdoc. 3, 3.1.3; Formulario T-14, paquete 1.6.5.1. Evidencia: Revisión del plan de capacitación (programación en al menos 2 ciclos por área y capacitación del turno de noche en su turno) y del registro de cobertura efectiva por área y turno. \\
TRA-004 Propio - Transición y preparación & El sistema debe registrar el tiempo que el aviso de cambio de material agrega al ciclo… & Cumple parcialmente & Botonera de Pala & Subdoc. 3, 3.3.3; Formulario T-14, paquetes 1.6.3.1 y 1.6.3.5. Falta: desglose por equipo y por turno del tiempo del aviso en el ciclo de carguío. Evidencia: Registro del tiempo del aviso en el ciclo de carguío por equipo y por turno durante la marcha blanca, contrastado con cronometraje en terreno de una muestra de ciclos con y sin aviso. \\
TRA-005 Propio - Transición y preparación & El sistema debe medir y publicar la adopción de cada función de registro en terreno… & Cumple & Indicadores de adopción por ola & Subdoc. 3, 3.2.5 y 3.4.5; Formulario T-14, paquete 1.6.6.4. Evidencia: Contraste mensual entre eventos registrados y eventos esperados por función, área y turno, con muestra auditada, contra la meta de RNF-USA-005. \\
TRA-006 Propio - Transición y preparación & El sistema debe permitir revertir la puesta en producción de una función de registro en terreno… & Cumple & Implantación en olas reversibles & Subdoc. 3, 3.2.2.4; Formulario T-14, paquetes 1.6.4.1 y 1.6.4.2. Evidencia: Ensayo de reversión de una función en un área durante la marcha blanca, verificando que las demás áreas siguen operando y midiendo el tiempo de reversión contra el máximo de 30 minutos. \\
TRA-007 Propio - Transición y preparación & El sistema debe registrar la capacitación efectivamente recibida por cada persona usuaria, con área, turno… & Cumple & Registro de capacitación por persona con área, turno y fecha, y cobertura por área y turno visible para el Cliente & Formulario T-14, paquete 1.6.5.1 (registro por persona con área, turno y fecha, y cobertura por área y turno visible para el Cliente). Evidencia: Contraste del registro de capacitación contra la nómina por área y turno, verificando la cobertura calculada. \\
TRA-008 Propio - Transición y preparación & El sistema debe permitir operar en paralelo con el método actual durante la marcha blanca, registrar… & Cumple & Marcha blanca en paralelo con planillas y radio & Subdoc. 3, 3.1.3 y 3.2.2.4; Formulario T-14, paquete 1.6.3.1. Evidencia: Informe de conciliación diaria de la marcha blanca entre la solución y el registro vigente, con las diferencias explicadas, y acta de apagado del método anterior. \\
TRA-MIG-001 Propio - Transición y preparación & El Proponente debe presentar en su propuesta un plan de migración de datos históricos con alcance… & Cumple & Migración por etapa & Subdoc. 3, 3.2.2.2 y 3.2.2.3; Subdoc. 5, 5.3 (Tabla \refext{S5}{tab:volumen-migracion-origen} y \refext{S5}{tab:naturaleza-dato-motor-razon-eleccion}.4; cálculo en la Tabla C.1 del Anexo C de OVERFORE-Subdocumento5-Anexos); Formulario T-14, paquetes 1.4.6.1 y 1.4.6.2. Evidencia: Revisión documental del plan de migración contra los siete elementos exigidos (alcance, origen, volumen, transformación, calidad, ejecución, reversión). \\
TRA-MIG-002 Propio - Transición y preparación & El Adjudicatario debe incluir en la migración una etapa de perfilado y saneamiento previo… & Cumple & Perfilado y saneamiento previo de los datos de origen, con informe de defectos y decisión registrada sobre cada uno & Subdoc. 5, 5.3.3 (perfilado y saneamiento, con la decisión sobre cada defecto); Formulario T-14, paquete 1.4.6.1. Evidencia: Revisión del informe de perfilado y saneamiento y verificación de que cada defecto detectado tiene una decisión registrada. \\
TRA-MIG-003 Propio - Transición y preparación & El Adjudicatario debe ejecutar al menos 2 ensayos completos de migración sobre el ambiente de Preproducción… & Cumple & Dos ensayos de migración en Preproducción, con tiempo total y resultado de conciliación medidos, antes de la carga definitiva & Subdoc. 5, 5.3.2 y 5.3.3; Formulario T-14, paquetes 1.4.6.2 y 1.4.6.4. Evidencia: Revisión de las bitácoras de los dos ensayos y contraste del resultado de conciliación de cada uno. \\
TRA-MIG-004 Propio - Transición y preparación & El Adjudicatario debe conciliar la migración definitiva con recuentos y sumas de control sobre el 100… & Cumple & Conciliación de la migración definitiva con completitud del 100 \% de los registros, sumas de control y muestra dirigida, por etapa & Subdoc. 5, 5.3.4 (conciliación y aceptación); Formulario T-14, paquetes 1.4.6.2 y 1.4.6.4. Evidencia: Recuento y suma de control sobre la totalidad de los registros migrados, más muestreo dirigido sobre los casos de mayor riesgo, contrastando origen contra destino. \\
TRA-MIG-005 Propio - Transición y preparación & El Adjudicatario debe migrar los datos históricos de cada dominio, por los años que fija… & Cumple & Migración de cinco años de producción, trazabilidad y laboratorio, y de diez años de acreditación (Etapa 1); tres años de mantenimiento y diez de seguridad, salud y ambiente (Etapa 2) & Subdoc. 3, 3.2.2.2 y 3.2.2.3; Subdoc. 5, 5.3.1, 5.3.3 y 5.3.5; Formulario T-14, paquetes 1.4.6.1, 1.4.6.3 y 1.4.6.5. Evidencia: Por cada dominio, recuento y suma de control origen contra destino sobre el 100 \% de los registros de la ventana, más muestreo dirigido (conciliación según TRA-MIG-004); verificación de que los registros fuera de la ventana son consultables en el repositorio durante el período de retención. Para los registros de seguridad digitalizados, verificación de legibilidad e indexación de las imágenes sobre una muestra. \\
TRA-009 Propio - Transición y preparación & El Adjudicatario debe presentar un plan de capacitación estructurado por perfil; entregar todo el material… & Cumple & Plan de capacitación por los cinco perfiles, material en español editable y de propiedad del Cliente, y certificación de administradores y equipo técnico & Formulario T-14, paquetes 1.6.5.1 y 1.6.5.2. Evidencia: Revisión documental del plan de capacitación contra los perfiles exigidos; revisión del material entregado en cuanto a idioma, formato editable y cesión de propiedad; registro de asistencia y evaluación por persona y acta de certificación. Cálculo del porcentaje certificado por perfil y por turno contra el mínimo de 95 \%. \\
TRA-010 Propio - Transición y preparación & El sistema debe permitir declarar la meta de adopción por función y por área y alertar… & Cumple & Metas de adopción por ola & Subdoc. 3, 3.4.5; Formulario T-14, paquete 1.6.6.4. Evidencia: Prueba de alerta con datos sintéticos de adopción bajo la meta durante el número de turnos parametrizado, y revisión de que el plan de acción de adopción se activa. \\
\end{tabla}
\begin{tabla}{I{0.134}I{0.257}I{0.110}I{0.141}I{0.358}}{Requerimientos del proyecto.}{tab:requerimientos-proyecto}
  {\cab{ID requerimiento} & \cab{Descripción} & \cab{Estado} & \cab{Componente que lo satisface} & \cab{Sección de la propuesta y evidencia}}
  {\emph{Fuente:} elaboración propia a partir de las Bases técnicas del caso Minería.}
RP-001 Propio - Del proyecto & El Adjudicatario debe mantener vigentes las certificaciones ISO 45001 e ISO 14001 y acreditar conocimiento… & Cumple parcialmente & Certificaciones ISO 45001 e ISO 14001; habilitación previa del personal que ingresa a faena & Subdoc. 1, sección 1.4 y Anexo D del Subdocumento 1; Formulario T-14, paquetes 1.1.7.1 y 1.8.7.2. Falta: experiencia comprobable en faena minera en operación. Evidencia: Revisión de los certificados ISO 45001 e ISO 14001 y de su fecha de vigencia, de la acreditación de conocimiento del Reglamento de Seguridad Minera y de la evidencia de experiencia en faena minera en operación (contratos o certificados de clientes). \\
RP-002 Propio - Del proyecto & Si la propuesta incorpora analítica predictiva, el Adjudicatario debe declarar, en la propuesta, con qué datos… & Cumple & Clasificador asistido de detenciones (innovación 2 del Subdocumento 13), uno de los dos componentes de inteligencia artificial, junto con la verificación facial de portería & Subdoc. 13, 13.2.2 (datos de entrenamiento y reentrenamiento trimestral) y 13.2.5 (precisión mínima); OVERFORE-Subdocumento13-Anexos, Anexo B, B.5 y B.6; Formulario T-14, paquetes 1.7.2.1, 1.7.2.2 y 1.7.2.3. Evidencia: Revisión documental, con el área de mantenimiento del Cliente, de los datos de entrenamiento, de la tasa de falsos positivos declarada y del mecanismo de mejora. \\
RP-003 Propio - Del proyecto & El Adjudicatario debe declarar, componente por componente, qué se ejecuta en la faena y qué… & Cumple parcialmente & Emplazamiento faena y nube: cinco módulos en el Nodo de Faena y el resto en la nube & Subdoc. 3, 3.3.2.1 y 3.3.5; Subdoc. 4, 4.2; OVERFORE-Subdocumento4-Anexos, Anexo B (Tabla B.1); Formulario T-14, paquete 1.3.1.2. Falta: clasificar en la Tabla B.1 el diodo de datos y los demás componentes físicos y de red de Subdoc. 3, 3.3.3 a 3.3.5 que no aparecen allí. Evidencia: Revisión de la documentación de arquitectura en cada hito de entrega. \\
RP-004 Propio - Del proyecto & El Adjudicatario debe entregar el estudio de cobertura del rajo, identificando las zonas de sombra… & Cumple & Estudio de cobertura LTE del rajo & Subdoc. 3, 3.3.4 y 3.2.2.2; Formulario T-14, paquetes 1.3.1.5 y 1.8.2.3. Evidencia: Revisión del estudio y contraste con la medición de cobertura en terreno, y verificación de una actualización cada seis meses. \\
RP-005 Propio - Del proyecto & El Adjudicatario debe levantar y mantener un inventario técnico versionado de los componentes de la solución… & Cumple & Inventario técnico de los sistemas existentes & Subdoc. 3, 3.2.2.2; Subdoc. 4, 4.1.1.5 (Tabla \refext{S4}{tab:version-fin-soporte-plan-actualizacion-componentes}); OVERFORE-Subdocumento4-Anexos, Anexo A (Tabla A.1, responsable por T11); Formulario T-14, paquetes 1.2.1.1 y 1.4.1.4. Evidencia: Revisión documental en cada hito de entrega. \\
RP-006 Propio - Del proyecto & El Adjudicatario debe ofrecer como servicio costeado toda función de operación de la solución que requiera… & Cumple & Operación como servicio (sala técnica y Nodo de Faena operados por Overfore) & Subdoc. 3, 3.2.3.3, 3.2.2.5 y 3.1.3; Formulario T-14, paquetes 1.6.5.3, 1.8.1.1 y 1.8.2.1. Evidencia: Revisión del modelo de operación propuesto contra la dotación declarada del Cliente. \\
RP-007 Propio - Del proyecto & El Adjudicatario debe operar un centro de atención en español, 24×7×365, con capacidad de respuesta presencial… & Cumple & Mesa de ayuda 24x7 en español con respuesta presencial en faena & Subdoc. 3, 3.2.2.5; Formulario T-14, paquetes 1.8.1.1 y 1.8.1.2. Evidencia: Revisión del horario y el idioma de atención; prueba de escalamiento de un incidente crítico provocado, con respuesta presencial en faena; revisión de la memoria de cálculo del dimensionamiento (método, volumen de contactos estimado y supuestos) contra los 1.250 contactos mensuales declarados \\
RP-008 Propio - Del proyecto & El personal del Adjudicatario que asista en faena debe estar acreditado como contratista y contar… & Cumple parcialmente & Personal de soporte acreditado y con examen de altura vigente & Subdoc. 3, 3.2.2.5 y 3.2.3.3; Formulario T-14, paquetes 1.1.7.1 y 1.8.1.2. Falta: tiempo de traslado desde Antofagasta (145 km) del personal que entra de turno; su costo se declara en la Oferta Económica. Evidencia: Revisión de las acreditaciones y de los exámenes de altura vigentes del personal en faena; revisión del plan de dotación, de los tiempos de traslado declarados y del costo de traslado en la oferta. \\
RP-009 Propio - Del proyecto & El Adjudicatario debe documentar los procedimientos de operación críticos de la solución y lograr… & Cumple parcialmente & Transferencia de conocimiento al equipo de TI del Cliente & Subdoc. 3, 3.1.3; Formulario T-14, paquetes 1.1.10.2 (documentación y ejecución sin asistencia de arranque, monitoreo, respaldo, restauración, escalamiento y atención de fallas de integración) y 1.3.1.4 (procedimientos de conmutación y retorno, ejecutables por el personal del Cliente tras la transferencia de conocimiento). Falta: que el personal del Cliente ejecute sin asistencia el procedimiento de conmutación antes del cierre de la marcha blanca (RP-009). Evidencia: Prueba de ejecución de cada procedimiento crítico por personal del Cliente sin asistencia, observada y registrada \\
RP-010 Propio - Del proyecto & El Adjudicatario debe proveer acompañamiento y mentoría al equipo técnico del Cliente durante al menos 6… & Cumple & Transferencia de conocimiento durante la operación & Subdoc. 3, 3.1.3; Formulario T-14, paquetes 1.6.5.3 (mentoría de seis meses a las once personas del área de TI del Cliente, meses 22 a 27), 1.8.9.1 (seis jornadas, dos por año), 1.8.9.2 (inducción del personal nuevo sin cargo adicional) y 1.8.1.1 (base de conocimiento mantenida durante el contrato, con métricas de uso y de utilidad). Evidencia: Registro de las sesiones de acompañamiento y mentoría (fecha, participantes y temas); actas de las jornadas anuales de actualización; registro de capacitación del personal nuevo y verificación de que no se cobró cargo adicional; revisión de la base de conocimiento y de sus métricas \\
RP-011 Propio - Del proyecto & El Adjudicatario debe planificar las validaciones que requieran al personal clave de operaciones fuera… & Cumple & Validaciones fuera de las detenciones de marzo y septiembre & Subdoc. 3, 3.2.2.5; Formulario T-14, paquete 1.2.2.5. Evidencia: Revisión del cronograma de validaciones contra las fechas de las detenciones programadas de marzo y septiembre. \\
\end{tabla}
\end{horizontal}
\end{document}
